%% Volume II, Chapter 14: Series Solutions
\chapter{Series Solutions}
\label{ch:series-solutions}

\abstract{The equations of the previous chapters had constant coefficients, and their solutions were built from exponentials, sines and cosines. When the coefficients vary, as they do for the vibrations of a circular drum, the potential around a sphere or the light near a caustic, the solutions are almost never elementary functions: the equation itself defines new functions. Power series are the first and most direct way to obtain them. Substituting a series turns the differential equation into a recurrence for the coefficients, which produces the solution term by term and tells how far it converges. The same recurrence is also the first algorithm for computing the new functions, and the chapter ends by asking when it works.}

\Cref{ch:linear-higher-order} solved linear equations with constant coefficients completely, and \cref{ch:linear-ode-systems} did the same for systems. For variable coefficients both chapters had to stop short. \Cref{sec:variable-coefficients} described the solutions of $y''+q(t)y=0$ without formulas, through Sturm's theorems on their zeros, and \cref{prop:commuting-coefficients} showed by an example that the solutions of a system $\mathbf y'=A(t)\mathbf y$ are in general not given by any formula in the integrals of the coefficients. \Cref{ch:first-order-equations} met the same obstacle from the other side: the Riccati equation $y'=t^2+y^2$ cannot be solved by quadratures, and the substitution $y=-w'/w$ turns it into the linear equation $w''+t^2w=0$, whose solutions are no more elementary. Yet the equations with variable coefficients that occur most often, those of Airy, Legendre and Bessel among them, have coefficients that are polynomials, and for such equations there is a method that succeeds wherever the leading coefficient does not vanish, and with a modification near the mildest of its zeros: one looks for the solution as a power series, as \vcref{1}{ch:series} promised at its end.

The method rests on the results of \vcref{1}{sec:power-series}: inside its interval of convergence a power series can be differentiated term by term (\vcref{1}{thm:termwise-diff}), two power series can be multiplied by the Cauchy product (\vcref{1}{prop:formal-analytic}), and a function has at most one power series expansion about a given point (\vcref{1}{cor:power-series-unique}). Substituting an unknown series into the equation and comparing coefficients turns the differential equation into a recurrence for the coefficients. The remaining question is whether the series produced by the recurrence converges; the main theorems of the chapter say that it does, as far as the coefficients of the equation allow.

\Cref{sec:series-idea} explains the method on equations whose solutions we already know and defines ordinary and singular points. \Cref{sec:ordinary-points} proves that near an ordinary point every solution is the sum of a power series whose radius of convergence is at least that of the coefficients. \Cref{sec:airy,sec:legendre} apply the theorem to the equations of Airy and Legendre; the first gives the functions that \cref{ch:linear-higher-order} drew and the solution of the Riccati equation of \cref{ch:first-order-equations}, the second the Legendre polynomials of \cref{ch:inner-products}, now as solutions of a differential equation. \Cref{sec:frobenius} treats regular singular points, at which the coefficients become infinite in the mildest possible way, by the method of Frobenius, whose model is the Euler--Cauchy equation of \cref{ch:linear-higher-order}. \Cref{sec:bessel} applies it to Bessel's equation. \Cref{sec:classical-equations} looks briefly at the other classical equations, and \cref{sec:series-computation} asks how well the series serve for computing the functions they define.

The independent variable of this chapter is $x$, not the variable $t$ of \cref{ch:first-order-equations,ch:linear-higher-order,ch:linear-ode-systems}, and primes denote derivatives with respect to $x$: $y'=dy/dx$. The equations of this chapter usually describe a shape in space rather than a motion in time, and the literature on the functions they define, the tables, the handbooks and the software, writes them in $x$, as \vcref{1}{sec:power-series} wrote its power series and \cref{ch:inner-products} its Legendre polynomials. Only in \cref{sec:airy}, where we return to the Riccati equation of \cref{ch:first-order-equations}, do we keep that chapter's variable $t$. Everything else follows the conventions of Part~II. The coefficients of power series may be real or complex; all estimates use the modulus and hold in both cases.

\section{Power Series Solutions: The Idea}
\label{sec:series-idea}

\subsection{Two Familiar Equations}

The equation $y'=y$ has the solutions $ce^x$, by \vcref{1}{thm:exp-ode}. Suppose we did not know this and looked for a solution of the form
\[
y(x)=\sum_{n=0}^\infty a_nx^n
\]
with a positive radius of convergence. By \vcref{1}{thm:termwise-diff}, $y'(x)=\sum_{n\ge1}na_nx^{n-1}=\sum_{n\ge0}(n+1)a_{n+1}x^n$ inside the interval of convergence, so the equation $y'-y=0$ says that the power series $\sum_n\bigl((n+1)a_{n+1}-a_n\bigr)x^n$ has the sum $0$. By the uniqueness of coefficients (\vcref{1}{cor:power-series-unique}), every coefficient vanishes:
\[
(n+1)\,a_{n+1}=a_n\qquad(n\in\mathbb N).
\]
This \term{recurrence} determines all coefficients from $a_0$: $a_1=a_0$, $a_2=a_1/2=a_0/2$, $a_3=a_2/3=a_0/3!$, and by induction $a_n=a_0/n!$. The series $a_0\sum x^n/n!$ has the radius of convergence $\infty$, and its sum is $a_0e^x$ (\vcref{1}{thm:exp-series}). Conversely, the computation read backwards shows that this series is a solution. The power series method has found the exponential function, together with the fact that every solution that is the sum of a power series is a multiple of it.

For the harmonic oscillator $y''+y=0$, the same substitution gives $y''=\sum_n(n+2)(n+1)a_{n+2}x^n$, and the equation becomes
\[
(n+2)(n+1)\,a_{n+2}=-a_n\qquad(n\in\mathbb N).
\]
Now each coefficient is determined by the one two places before it, and the coefficients with even and with odd indices do not interact. The values $a_0=y(0)$ and $a_1=y'(0)$ are free, and
\[
a_{2k}=\frac{(-1)^ka_0}{(2k)!},\qquad a_{2k+1}=\frac{(-1)^ka_1}{(2k+1)!},
\]
so that $y=a_0\cos x+a_1\sin x$ by \vcref{1}{thm:sin-cos-series}. Again the method produces a two-dimensional family of solutions, parametrised by the initial values, in accordance with \cref{thm:solution-space}.

\subsection{The Method}

The two examples follow a pattern that works for every linear equation whose coefficients are power series. To solve $y''+p(x)y'+q(x)y=0$ near a point $x_0$:
\begin{enumerate}
\item write the unknown solution as $y=\sum_na_n(x-x_0)^n$ with undetermined coefficients, and the coefficients $p$ and $q$ as power series about $x_0$;
\item differentiate term by term, multiply the series by the Cauchy product, and collect the terms with the same power of $x-x_0$;
\item set every coefficient of the result equal to zero, which by the uniqueness of coefficients is necessary and sufficient for a convergent series to be a solution; this gives a recurrence that determines $a_2,a_3,\ldots$ from $a_0$ and $a_1$;
\item prove that the series so obtained converges, and determine its radius of convergence.
\end{enumerate}
The first three steps are formal manipulations, of the kind that \vcref{1}{sec:gf} carried out with generating functions, and \vcref{1}{prop:formal-analytic} guarantees that they are valid wherever the series converge. The fourth step is the real work. Without it the method can fail: \cref{sec:frobenius} meets an equation for which the recurrence produces a series that converges nowhere except at its centre. The theorem of \cref{sec:ordinary-points} settles the fourth step once and for all under a simple condition on the coefficients.

\subsection{Analytic Functions, Ordinary and Singular Points}

The condition concerns the coefficients $p$ and $q$. They must be given by power series, and the right notion is that of an analytic function, which \vcref{1}{def:taylor-series} introduced through the Taylor series.

\begin{definition}[Analytic function]
\label{def:analytic}
A real- or complex-valued function $f$ defined on an open interval containing $x_0$ is \term{analytic at $x_0$} if there is a power series $\sum_na_n(x-x_0)^n$ with a positive radius of convergence whose sum equals $f(x)$ for all $x$ in some open interval around $x_0$. It is analytic on an open interval $I$ if it is analytic at every point of $I$.
\end{definition}

By \vcref{1}{thm:termwise-diff}, the series is then the Taylor series of $f$ at $x_0$, so the definition agrees with that of \vcref{1}{def:taylor-series}. Polynomials are analytic everywhere, and so are $e^x$, $\sin x$ and $\cos x$, whose series converge on all of $\mathbb R$. The function $1/(1-x)$ is analytic at $0$, with the geometric series of radius $1$; it is in fact analytic at every point $x_0\ne1$, since $1/(1-x)=1/((1-x_0)-(x-x_0))$ is a geometric series in $(x-x_0)/(1-x_0)$. Sums and products of functions analytic at $x_0$ are analytic at $x_0$, and so is $1/f$ if $f$ is analytic at $x_0$ with $f(x_0)\ne0$, by \vcref{1}{prop:formal-analytic}. A function analytic at $x_0$ is infinitely differentiable near $x_0$, and the smooth function $e^{-1/x^2}$ of \vcref{1}{prop:flat-function}, whose Taylor series at $0$ is zero, shows that the converse fails.

\begin{definition}[Ordinary and singular points]
\label{def:ordinary-point}
Consider the equation
\begin{equation}
y''+p(x)\,y'+q(x)\,y=0 ,
\label{eq:second-order-normal}
\end{equation}
with coefficients $p$ and $q$ defined on an open interval around $x_0$, except possibly at $x_0$. The point $x_0$ is an \term{ordinary point} of the equation if $p$ and $q$, suitably defined at $x_0$, are analytic at $x_0$, and a \term{singular point} otherwise.
\end{definition}

Most equations of interest come in the form $P(x)y''+Q(x)y'+R(x)y=0$ with polynomial coefficients, and their normal form \eqref{eq:second-order-normal} has the rational coefficients $p=Q/P$ and $q=R/P$. For them the distinction is easy to make.

\begin{proposition}[Polynomial coefficients]
\label{prop:polynomial-ordinary}
Let $P$, $Q$ and $R$ be polynomials without a common zero, $P\ne0$. A point $x_0$ is an ordinary point of $y''+(Q/P)y'+(R/P)y=0$ if and only if $P(x_0)\ne0$.
\end{proposition}

\begin{proof}
A polynomial is a power series about every point, with finitely many nonzero coefficients. If $P(x_0)\ne0$, then $1/P$ is analytic at $x_0$ by \vcref{1}{prop:formal-analytic}\,(iii), and so are $Q/P$ and $R/P$, by part (i) of the same proposition. If $P(x_0)=0$, then $Q(x_0)\ne0$ or $R(x_0)\ne0$; say $R(x_0)\ne0$. Near $x_0$, $|R(x)|\ge|R(x_0)|/2$ by continuity, while $|P(x)|\le K|x-x_0|$ for some constant $K$, since $P(x)=(x-x_0)\tilde P(x)$ with a polynomial $\tilde P$ (\vcref{1}{cor:factor}). So $|R(x)/P(x)|\to\infty$ as $x\to x_0$, and no definition of $R/P$ at $x_0$ makes it continuous there; but the sum of a power series is continuous (\vcref{1}{thm:termwise-diff}). The case $Q(x_0)\ne0$ is the same.
\end{proof}

Airy's equation $y''-xy=0$ has no singular point. Legendre's equation $(1-x^2)y''-2xy'+\nu(\nu+1)y=0$ has the singular points $\pm1$, Bessel's equation $x^2y''+xy'+(x^2-\nu^2)y=0$ the singular point $0$, and the Euler--Cauchy equation $x^2y''+axy'+by=0$ of \cref{prop:euler-cauchy} also the singular point $0$, unless $a=b=0$. At a singular point the solutions may misbehave: the Euler--Cauchy solutions $x^r$ are not analytic at $0$ unless $r\in\mathbb N$, and $x^r\ln x$ never is. Away from the singular points, the next section shows, nothing of the kind can happen.

\section{Solutions Near an Ordinary Point}
\label{sec:ordinary-points}

\subsection{The Recurrence}

We first carry out the formal steps once and for all. By a translation we may take $x_0=0$; the general case follows by writing everything in powers of $x-x_0$.

\begin{lemma}[Substituting a power series]
\label{lem:series-substitution}
Let $p(x)=\sum_kp_kx^k$, $q(x)=\sum_kq_kx^k$ and $y(x)=\sum_na_nx^n$ be power series with radii of convergence at least $R>0$. Then for $|x|<R$,
\begin{gather*}
y''+p\,y'+q\,y=\sum_{n=0}^\infty c_nx^n,\\
c_n\coloneqq(n+2)(n+1)\,a_{n+2}+\sum_{k=0}^n\bigl((k+1)\,a_{k+1}\,p_{n-k}+a_k\,q_{n-k}\bigr),
\end{gather*}
and the series on the right has a radius of convergence at least $R$.
\end{lemma}

\begin{proof}
By \vcref{1}{thm:termwise-diff}, $y'=\sum_n(n+1)a_{n+1}x^n$ and $y''=\sum_n(n+2)(n+1)a_{n+2}x^n$ for $|x|<R$, with radii at least $R$. By \vcref{1}{prop:formal-analytic}\,(i), the product $p\,y'$ is the sum of the Cauchy product of the two series, whose coefficient of $x^n$ is $\sum_{k=0}^n(k+1)a_{k+1}p_{n-k}$, with radius at least $R$, and similarly $q\,y=\sum_n\bigl(\sum_{k=0}^na_kq_{n-k}\bigr)x^n$. Adding the three series (\vcref{1}{thm:series-rules}) gives the claim.
\end{proof}

So a power series $y$ with radius at least $R$ solves \eqref{eq:second-order-normal} on $(-R,R)$ if and only if all $c_n$ vanish, by \vcref{1}{cor:power-series-unique}, that is, if and only if
\begin{equation}
(n+2)(n+1)\,a_{n+2}=-\sum_{k=0}^n\bigl((k+1)\,p_{n-k}\,a_{k+1}+q_{n-k}\,a_k\bigr)\qquad(n\in\mathbb N).
\label{eq:ordinary-recurrence}
\end{equation}
The right side involves only $a_0,\ldots,a_{n+1}$, so the recurrence determines $a_2,a_3,\ldots$ uniquely from $a_0$ and $a_1$, which are free; they are $y(0)$ and $y'(0)$. What remains is to show that the series converges.

\subsection{Bounding the Coefficients}

The recurrence expresses each coefficient through all the earlier ones, with a factor of order $1/n^2$ in front. The following lemma shows that such a recurrence cannot make the coefficients grow faster than any geometric sequence $\lambda^n$ with $\lambda>1$. The idea of bounding the solution of a recurrence by the solution of a simpler one goes back to Cauchy's method of majorants~\cite[Chap.~2]{hille1976}; the form given here works with real sequences only.

\begin{lemma}[Growth of solutions of recurrences]
\label{lem:coefficient-growth}
Let $(b_n)$ be a sequence of nonnegative numbers, let $\alpha,\beta,D\ge0$ and $\lambda>1$, and let $(\gamma_n)$ be positive numbers with $(\alpha+\beta n)/\gamma_n\to0$ as $n\to\infty$. Suppose that for some $n_0$,
\begin{equation}
\gamma_n\,b_n\le\sum_{k=0}^{n-1}(\alpha+\beta k)\,b_k+(\alpha+\beta n)\,D\lambda^n\qquad(n\ge n_0).
\label{eq:majorant-inequality}
\end{equation}
Then there is a constant $C$ with $b_n\le C\lambda^n$ for all $n\in\mathbb N$.
\end{lemma}

\begin{proof}
Choose $N\ge n_0$ such that $2(\alpha+\beta n)\le(\lambda-1)\gamma_n$ for all $n\ge N$; this is possible because $(\alpha+\beta n)/\gamma_n\to0$. Let $C\ge D(\lambda-1)$ be so large that $b_n\le C\lambda^n$ for the finitely many $n<N$. We show by induction that $b_n\le C\lambda^n$ for $n\ge N$ as well. If it holds for all indices below $n$, then, since $\alpha+\beta k\le\alpha+\beta n$ for $k<n$,
\[
\sum_{k=0}^{n-1}(\alpha+\beta k)\,b_k\le C(\alpha+\beta n)\sum_{k=0}^{n-1}\lambda^k\le\frac{C(\alpha+\beta n)\,\lambda^n}{\lambda-1},
\]
by the formula for the geometric sum. With $D\le C/(\lambda-1)$, \eqref{eq:majorant-inequality} gives
\[
\gamma_nb_n\le(\alpha+\beta n)\lambda^n\Bigl(\frac C{\lambda-1}+D\Bigr)\le\frac{2(\alpha+\beta n)}{\lambda-1}\,C\lambda^n\le\gamma_nC\lambda^n,
\]
and dividing by $\gamma_n>0$ gives $b_n\le C\lambda^n$.
\end{proof}

The lemma is used as follows. If the coefficients of a power series satisfy $|a_n|\rho^n\le C\lambda^n$, then $|a_n|r^n\le C$ for $r=\rho/\lambda$, and the radius of convergence is at least $\rho/\lambda$ (\vcref{1}{def:radius}). If this holds for every $\rho<R$ and every $\lambda>1$, the radius is at least $R$.

\subsection{The Theorem}

\begin{theorem}[Solutions near an ordinary point]
\label{thm:ordinary-point}
Let $p(x)=\sum_kp_k(x-x_0)^k$ and $q(x)=\sum_kq_k(x-x_0)^k$ be power series with radii of convergence at least $R$, where $0<R\le\infty$, and let $I\coloneqq(x_0-R,x_0+R)$.
\begin{enumerate}
\item For all numbers $a_0$ and $a_1$ there is exactly one power series $y=\sum_na_n(x-x_0)^n$ with these first two coefficients whose sum solves $y''+p(x)y'+q(x)y=0$ near $x_0$. Its coefficients are given by the recurrence \eqref{eq:ordinary-recurrence}, its radius of convergence is at least $R$, and it solves the equation on all of $I$.
\item Every solution of the equation on $I$ is the sum of such a power series, the one with $a_0=y(x_0)$ and $a_1=y'(x_0)$.
\item The solutions $y_1$ and $y_2$ with $(y_1(x_0),y_1'(x_0))=(1,0)$ and $(y_2(x_0),y_2'(x_0))=(0,1)$ form a fundamental set on $I$.
\end{enumerate}
If $p$, $q$, $a_0$ and $a_1$ are real, all coefficients $a_n$ are real.
\end{theorem}

\begin{proof}
We take $x_0=0$. (i) By \cref{lem:series-substitution} and \vcref{1}{cor:power-series-unique}, a power series with a positive radius solves the equation near $0$ exactly when its coefficients satisfy \eqref{eq:ordinary-recurrence}, and the recurrence determines them uniquely from $a_0$ and $a_1$. It remains to bound the coefficients. Fix $\rho$ with $0<\rho<R$. Since $\rho<R$, the sequences $(p_k\rho^k)$ and $(q_k\rho^k)$ are bounded (\vcref{1}{def:radius}); let $M$ be a common bound, so that $|p_j|,|q_j|\le M\rho^{-j}$. Let $b_n\coloneqq|a_n|\rho^n$. Taking moduli in \eqref{eq:ordinary-recurrence} and multiplying by $\rho^{n+2}$,
\begin{align*}
(n+2)(n+1)\,b_{n+2}&\le\rho^{n+2}\sum_{k=0}^nM\rho^{\,k-n}\bigl((k+1)|a_{k+1}|+|a_k|\bigr)\\
&=M\sum_{k=0}^n\bigl((k+1)\rho\,b_{k+1}+\rho^2\,b_k\bigr).
\end{align*}
With $m\coloneqq n+2$, every index on the right is less than $m$, and each $b_j$ appears with a factor at most $M(\rho^2+\rho j)$. Hence
\[
m(m-1)\,b_m\le\sum_{j=0}^{m-1}M(\rho^2+\rho j)\,b_j\qquad(m\ge2),
\]
which is \eqref{eq:majorant-inequality} with $\gamma_m=m(m-1)$, $\alpha=M\rho^2$, $\beta=M\rho$ and $D=0$. Since $(\alpha+\beta m)/(m(m-1))\to0$, \cref{lem:coefficient-growth} gives $b_m\le C\lambda^m$ for every $\lambda>1$, and by the remark after the lemma the radius of convergence of $\sum a_nx^n$ is at least $R$. By \cref{lem:series-substitution}, the sum then solves the equation on all of $I$.

(ii) The functions $p$ and $q$ are continuous on $I$ (\vcref{1}{thm:termwise-diff}), so by \cref{cor:higher-existence} the equation has exactly one solution on $I$ with given values of $y(0)$ and $y'(0)$. The power series of (i) with $a_0=y(0)$ and $a_1=y'(0)$ is such a solution, since $a_0$ and $a_1$ are its value and derivative at $0$ (\vcref{1}{thm:termwise-diff}); so it is the solution.

(iii) The Wronskian of $y_1$ and $y_2$ at $x_0$ is $1\cdot1-0\cdot0=1\ne0$, so they form a fundamental set by \cref{thm:wronskian-solutions}. If $p$, $q$, $a_0$ and $a_1$ are real, the recurrence involves only real numbers.
\end{proof}

The theorem answers the question of the fourth step: near an ordinary point the method of \cref{sec:series-idea} always succeeds, and the series converge at least as far as the series of the coefficients do. Part (ii) says more: the theorem does not merely produce some solutions, it shows that every solution on $I$ is the sum of its Taylor series there, so that all solutions are analytic wherever the coefficients are. For the examples of \cref{sec:series-idea}, $p$ and $q$ are constants, $R=\infty$, and the series of $e^x$, $\cos x$ and $\sin x$ converge everywhere, as they must.

The equation $y''+xy'+y=0$ is a first example with a variable coefficient. Here $p(x)=x$ and $q(x)=1$, and the recurrence \eqref{eq:ordinary-recurrence} reads $(n+2)(n+1)a_{n+2}=-na_n-a_n$, that is,
\[
a_{n+2}=-\frac{a_n}{n+2}\qquad(n\in\mathbb N).
\]
The even coefficients give $y_1=1-x^2/2+x^4/8-\cdots=\sum_k(-1)^kx^{2k}/(2^kk!)$, which is $e^{-x^2/2}$, and the odd ones give
\[
y_2(x)=x-\frac{x^3}3+\frac{x^5}{3\cdot5}-\cdots=\sum_{k=0}^\infty\frac{(-1)^k\,x^{2k+1}}{1\cdot3\cdots(2k+1)} .
\]
Both converge on all of $\mathbb R$, as the theorem predicts. The second solution is not an elementary function, but it has a closed form of another kind: since the equation is $(y'+xy)'=0$, its solutions satisfy $y'+xy=c$, and \cref{thm:linear-first-order} gives $y_2(x)=e^{-x^2/2}\int_0^xe^{s^2/2}\,ds$. Expanding this product of two series reproduces the coefficients above. In most equations no such integral is available, and the series is the solution.

\subsection{How Far the Series Converge}

The theorem gives a lower bound for the radius of convergence, and the bound may be pessimistic. For the equation $(1-x^2)y''-2xy'+2y=0$ of \cref{sec:wronskian}, which is Legendre's equation with $\nu=1$, the coefficients $p(x)=-2x/(1-x^2)=-2\sum_kx^{2k+1}$ and $q(x)=2/(1-x^2)$ have radius $1$, so all solutions are power series about $0$ with radius at least $1$; one of them is the polynomial $x$, with radius $\infty$. Whether the bound is attained can only be decided by looking at the solutions.

What decides the radius of the coefficients themselves is less obvious than it seems. The equation
\[
(1+x^2)\,y''+2x\,y'=0,\qquad\text{that is,}\qquad\bigl((1+x^2)\,y'\bigr)'=0 ,
\]
has the ordinary point $x_0$ for every real $x_0$, by \cref{prop:polynomial-ordinary}: its coefficients are analytic on the whole real line. Its solutions satisfy $(1+x^2)y'=c$, so they are the functions $c_1+c\arctan x$, and the power series of $\arctan x$ about $0$ has the radius $1$ (\vcref{1}{thm:arctan-series}). The coefficient $p(x)=2x/(1+x^2)$ has a power series about $0$ of radius $1$ as well, and the theorem gives exactly the right answer, but the reason lies outside the real line, as the note on complex power series in \vcref{1}{sec:power-series} explained: $1+x^2$ vanishes at the complex points $\pm i$, at distance $1$ from $0$. The general statement belongs to the theory of functions of a complex variable.

\begin{theorem}[Radius and complex singular points]
\label{thm:rational-coefficients}
Let $P$, $Q$ and $R$ be polynomials with real or complex coefficients, let $x_0$ be real with $P(x_0)\ne0$, and let $\rho$ be the distance from $x_0$ to the nearest complex zero of $P$ ($\rho=\infty$ if $P$ is constant). Then every solution of $P(x)y''+Q(x)y'+R(x)y=0$ on an interval around $x_0$ is the sum of a power series about $x_0$ whose radius of convergence is at least $\rho$.
\end{theorem}

\begin{proof}[sketch]
By \cref{thm:ordinary-point} it suffices to show that $Q/P$ and $R/P$ are sums of power series about $x_0$ with radius at least $\rho$. Over the complex numbers, $P$ splits into linear factors $x-z_j$ with $|z_j-x_0|\ge\rho$, and a partial fraction decomposition with complex coefficients writes $Q/P$ and $R/P$ as a polynomial plus a combination of terms $1/(x-z_j)^m$. Each such term is $(-1)^m(z_j-x_0)^{-m}(1-w)^{-m}$ with $w=(x-x_0)/(z_j-x_0)$, and the binomial series of $(1-w)^{-m}$ converges for $|w|<1$, that is, for $|x-x_0|<|z_j-x_0|$. Carrying this out requires power series with complex variable, which the note in \vcref{1}{sec:power-series} describes and Volume~III develops; there it is also shown that the same bound holds for equations whose coefficients are analytic functions of a complex variable in a disc~\cite{hille1976}.
\end{proof}

For Legendre's equation, $P(x)=1-x^2$ has the zeros $\pm1$, and the bound is $1$ about $0$, and $1-|x_0|$ about a point $x_0\in(-1,1)$. For $(1+x^2)y''+2xy'=0$ it is $\sqrt{1+x_0^2}$ about $x_0$, the distance to $\pm i$; for $x_0=0$ it is attained by $\arctan x$. For Airy's equation it is $\infty$.

\subsection{Equations of Higher Order and Systems}

Nothing in the argument depends on the order two. For a linear system $\mathbf y'=A(x)\mathbf y$ whose coefficient matrix is given by power series, the proof is even simpler, and equations of order $k$ are included through their companion systems (\cref{prop:companion-system}). \Cref{prop:commuting-coefficients} showed that such systems have no general formula in terms of integrals of the coefficients; when the coefficients are analytic, they have power series solutions instead.

\begin{proposition}[Systems with analytic coefficients]
\label{prop:analytic-systems}
Let $A(x)=\sum_kA_k(x-x_0)^k$, where the $A_k$ are $n\times n$ matrices and each entry of $A$ is a power series with radius of convergence at least $R>0$. Then every solution of $\mathbf y'=A(x)\mathbf y$ on $I=(x_0-R,x_0+R)$ is $\mathbf y(x)=\sum_m\mathbf a_m(x-x_0)^m$, where $\mathbf a_0=\mathbf y(x_0)$, the coefficient vectors satisfy
\begin{equation}
(m+1)\,\mathbf a_{m+1}=\sum_{k=0}^mA_{m-k}\,\mathbf a_k\qquad(m\in\mathbb N),
\label{eq:system-recurrence}
\end{equation}
and each component series has radius of convergence at least $R$.
\end{proposition}

\begin{proof}
Take $x_0=0$, let $\mathbf a_0$ be given and define $\mathbf a_1,\mathbf a_2,\ldots$ by \eqref{eq:system-recurrence}. Fix $0<\rho<R$. Each entry of $A_k\rho^k$ is bounded in $k$ (\vcref{1}{def:radius}), so there is $M$ such that every entry of $A_k$ has modulus at most $M\rho^{-k}/n$ for all $k$. Each component of $A_k\mathbf v$ is a sum of $n$ products of an entry of $A_k$ with a component of $\mathbf v$, so $\|A_k\mathbf v\|_\infty\le M\rho^{-k}\|\mathbf v\|_\infty$ for every real or complex vector $\mathbf v$. With $b_m\coloneqq\|\mathbf a_m\|_\infty\rho^m$ the recurrence gives $(m+1)b_{m+1}\le M\rho\sum_{k=0}^mb_k$, that is, $mb_m\le\sum_{k<m}M\rho\,b_k$ for $m\ge1$. This is \eqref{eq:majorant-inequality} with $\gamma_m=m$, $\alpha=M\rho$ and $\beta=D=0$, and $\alpha/m\to0$; so $b_m\le C\lambda^m$ for every $\lambda>1$ by \cref{lem:coefficient-growth}, and each component of $\sum\mathbf a_mx^m$ has radius at least $R$. Differentiating term by term and multiplying the series entry by entry (\vcref{1}{thm:termwise-diff,prop:formal-analytic}), the coefficient of $x^m$ in $\mathbf y'-A\mathbf y$ is $(m+1)\mathbf a_{m+1}-\sum_{k\le m}A_{m-k}\mathbf a_k=\mathbf 0$, so the series is a solution on $I$. By \cref{thm:linear-existence}, it is the only solution with the value $\mathbf a_0$ at $0$.
\end{proof}

For the system of \cref{prop:commuting-coefficients}, written in the variable $x$, $A(x)=A_0+A_1x$ with
\[
A_0=\begin{bmatrix}1&0\\0&0\end{bmatrix},\qquad A_1=\begin{bmatrix}0&1\\0&0\end{bmatrix},
\]
and the solution with $\mathbf y(0)=\mathbf e_2$. The recurrence \eqref{eq:system-recurrence} gives $(m+1)a_{m+1,2}=0$ for the second components, so $y_2\equiv1$, and $(m+1)a_{m+1,1}=a_{m,1}+a_{m-1,2}$ for the first, with $a_{0,1}=0$. Hence $a_{1,1}=0$, $a_{2,1}=1/2$ and $a_{m+1,1}=a_{m,1}/(m+1)$ for $m\ge2$, so that $a_{m,1}=1/m!$ for $m\ge2$ and $y_1=\sum_{m\ge2}x^m/m!=e^x-1-x$, the solution found there by other means. The power series method treats every system with analytic coefficients in this way, whether or not the values of $A$ commute.

\section{Airy's Equation}
\label{sec:airy}

\subsection{The Series Solutions}

\term{Airy's equation}
\begin{equation}
y''=xy
\label{eq:airy}
\end{equation}
is the simplest equation of the second order with a variable coefficient, and \cref{sec:variable-coefficients} described its solutions qualitatively: they oscillate for $x<0$, ever faster as $x$ decreases, and have at most one zero for $x>0$. Here $p=0$ and $q(x)=-x$, so every point is ordinary, and \cref{thm:ordinary-point} applies with $R=\infty$. With $y=\sum a_nx^n$, the product $xy$ is $\sum_{n\ge1}a_{n-1}x^n$, and comparing the coefficients of $x^n$ on both sides gives $2a_2=0$ for $n=0$ and
\[
(n+2)(n+1)\,a_{n+2}=a_{n-1}\qquad(n\ge1).
\]
Each coefficient is determined by the one three places before it. Since $a_2=0$, the coefficients $a_5,a_8,\ldots$ all vanish, and the solution splits into two series, one with the powers $x^{3k}$, started by $a_0$, and one with the powers $x^{3k+1}$, started by $a_1$.

\begin{proposition}[Series solutions of Airy's equation]
\label{prop:airy-series}
The solutions of $y''=xy$ with $(y(0),y'(0))=(1,0)$ and $(0,1)$ are
\begin{gather*}
f(x)=1+\sum_{k=1}^\infty\frac{x^{3k}}{(2\cdot3)(5\cdot6)\cdots((3k-1)\cdot3k)},\\
g(x)=x+\sum_{k=1}^\infty\frac{x^{3k+1}}{(3\cdot4)(6\cdot7)\cdots(3k\cdot(3k+1))} .
\end{gather*}
Both series converge on all of $\mathbb R$, $f$ and $g$ form a fundamental set with $W(f,g)\equiv1$, and on $(0,\infty)$ both functions are positive and increasing, with $f(x)\ge1+x^3/6$ and $g(x)\ge x+x^4/12$.
\end{proposition}

\begin{proof}
The coefficients follow from the recurrence: $a_{3k}=a_{3k-3}/((3k-1)\,3k)$ with $a_0=1$, and $a_{3k+1}=a_{3k-2}/(3k\,(3k+1))$ with $a_1=1$. By \cref{thm:ordinary-point}, with $R=\infty$, the series converge everywhere and form a fundamental set. The equation has no first-derivative term, so by Abel's formula (\cref{thm:abel}) the Wronskian is constant, equal to its value $f(0)g'(0)-f'(0)g(0)=1$. For $x>0$ all terms of both series are positive, which gives the lower bounds, and so are all terms of the differentiated series, so $f'>0$ and $g'>0$ there.
\end{proof}

The coefficients decrease so fast that the series are practical for moderate $|x|$. For instance $f(1)=1+1/6+1/180+1/12960+\cdots=1.1722999\ldots$; the fourth term is already below $10^{-4}$. \Cref{fig:airy} of \cref{ch:linear-higher-order} was computed from these series. Summing them with the ball arithmetic of python-flint (version 0.9.0), which carries rigorous error bounds, and locating the sign changes by bisection gives the zeros of $f$ nearest to $0$ as $-1.9864$, $-3.8253$ and $-5.2956$, and those of $g$, besides $0$, as $-2.6664$, $-4.3425$ and $-5.7410$, in agreement with the values quoted in \cref{sec:variable-coefficients}.

\subsection{The Airy Functions}

Every solution of Airy's equation is a combination $c_1f+c_2g$, and two combinations are distinguished by their behaviour for large positive $x$. Since $f$ and $g$ grow without bound there, the combinations that do so as slowly as possible, or decay, have to be found by balancing the two.

\begin{definition}[Airy functions]
\label{def:airy-functions}
The \term{Airy functions} are the solutions
\[
\operatorname{Ai}(x)\coloneqq c_1f(x)-c_2g(x),\qquad\operatorname{Bi}(x)\coloneqq\sqrt3\,\bigl(c_1f(x)+c_2g(x)\bigr)
\]
of Airy's equation, with the constants
\begin{gather*}
c_1\coloneqq\operatorname{Ai}(0)=\frac1{3^{2/3}\,\Gamma(2/3)}=0.355028\ldots,\\
c_2\coloneqq-\operatorname{Ai}'(0)=\frac1{3^{1/3}\,\Gamma(1/3)}=0.258819\ldots .
\end{gather*}
\end{definition}

The values of the gamma function (\vcref{1}{def:gamma}) are $\Gamma(1/3)=2.678939\ldots$ and $\Gamma(2/3)=1.354118\ldots$. The Wronskian of the two functions is
\begin{align*}
W(\operatorname{Ai},\operatorname{Bi})&=\sqrt3\bigl(c_1c_2W(f,g)-c_2c_1W(g,f)\bigr)\\
&=2\sqrt3\,c_1c_2=\frac{2}{\sqrt3\,\Gamma(1/3)\,\Gamma(2/3)}\ne0 ,
\end{align*}
so they form a fundamental set. With the reflection formula $\Gamma(s)\Gamma(1-s)=\pi/\sin(\pi s)$ of the note in \vcref{1}{sec:general-powers}, which gives $\Gamma(1/3)\Gamma(2/3)=2\pi/\sqrt3$, the Wronskian is $1/\pi$.

The constants are not arbitrary. Airy met the function $\operatorname{Ai}$ in 1838, as an integral that a change of variable turns into
\[
\operatorname{Ai}(x)=\frac1\pi\int_0^\infty\cos\Bigl(\frac{s^3}3+xs\Bigr)\,ds ,
\]
when he computed the intensity of light near a caustic, the bright curve in which light rays reflected or refracted by a curved surface accumulate, as in a rainbow~\cite{airy1838}. One can show that this integral converges, solves \eqref{eq:airy}, and has the value and derivative at $0$ given above, and that its solution is distinguished by its decay: as $x\to\infty$,
\[
\operatorname{Ai}(x)\sim\frac{e^{-\zeta}}{2\sqrt\pi\,x^{1/4}},\quad\operatorname{Bi}(x)\sim\frac{e^{\zeta}}{\sqrt\pi\,x^{1/4}},\quad\zeta\coloneqq\frac23x^{3/2},
\]
so that $\operatorname{Ai}$ is, up to a constant factor, the only solution that tends to zero, and $\operatorname{Bi}$ is a solution that grows at the same rate as all the others, normalised so that it has the same amplitude as $\operatorname{Ai}$ for $x\to-\infty$~\cite[Chap.~9]{olver2010nist}. These facts need the methods of asymptotic analysis, which Volume~XI develops. \Cref{fig:airy-functions}A shows the two functions. Their first zeros are $-2.33811$, $-4.08795$ and $-5.52056$ for $\operatorname{Ai}$ and $-1.17371$, $-3.27109$ and $-4.83074$ for $\operatorname{Bi}$, and they interlace, as Sturm's separation theorem (\cref{thm:sturm-separation}) requires.

The integral explains why Airy's function appears wherever waves turn back. Near a point where a quantity that is oscillatory on one side becomes exponential on the other, as for a light ray at a caustic or a quantum particle at the edge of a classically forbidden region, the governing equation looks locally like $y''=xy$ after a change of scale, and the transition is described by $\operatorname{Ai}$~\cite{vallee2004}.

\begin{figure}
\sdpanel{A}{\begin{tikzpicture}
\begin{axis}[sdplot, width=0.56\textwidth, height=0.42\textwidth, xmin=-10, xmax=2.5, ymin=-0.6, ymax=1.6,
  xtick={-10,-8,-6,-4,-2,0,2}, ytick={-0.5,0,0.5,1,1.5}, xlabel={$x$}, legend pos=north west]
\addplot[black, line width=0.3pt, domain=-10:2.5, samples=2, forget plot] {0};
\addplot[black, line width=0.8pt, smooth] coordinates {(-10,0.0402) (-9.95,0.0893) (-9.9,0.1362) (-9.85,0.1798) (-9.8,0.2189) (-9.75,0.2526) (-9.7,0.2802) (-9.65,0.3011) (-9.6,0.3147) (-9.55,0.3207) (-9.5,0.3191) (-9.45,0.3099) (-9.4,0.2935) (-9.35,0.2701) (-9.3,0.2405) (-9.25,0.2052) (-9.2,0.1653) (-9.15,0.1215) (-9.1,0.0750) (-9.05,0.0267) (-9,-0.0221) (-8.95,-0.0705) (-8.9,-0.1173) (-8.85,-0.1614) (-8.8,-0.2021) (-8.75,-0.2382) (-8.7,-0.2692) (-8.65,-0.2943) (-8.6,-0.3131) (-8.55,-0.3252) (-8.5,-0.3303) (-8.45,-0.3284) (-8.4,-0.3196) (-8.35,-0.3041) (-8.3,-0.2822) (-8.25,-0.2545) (-8.2,-0.2216) (-8.15,-0.1841) (-8.1,-0.1429) (-8.05,-0.0988) (-8,-0.0527) (-7.95,-0.0056) (-7.9,0.0417) (-7.85,0.0881) (-7.8,0.1329) (-7.75,0.1750) (-7.7,0.2137) (-7.65,0.2484) (-7.6,0.2783) (-7.55,0.3029) (-7.5,0.3218) (-7.45,0.3347) (-7.4,0.3413) (-7.35,0.3417) (-7.3,0.3358) (-7.25,0.3237) (-7.2,0.3059) (-7.15,0.2825) (-7.1,0.2540) (-7.05,0.2211) (-7,0.1843) (-6.95,0.1442) (-6.9,0.1017) (-6.85,0.0574) (-6.8,0.0121) (-6.75,-0.0334) (-6.7,-0.0783) (-6.65,-0.1219) (-6.6,-0.1635) (-6.55,-0.2024) (-6.5,-0.2380) (-6.45,-0.2698) (-6.4,-0.2971) (-6.35,-0.3198) (-6.3,-0.3373) (-6.25,-0.3496) (-6.2,-0.3564) (-6.15,-0.3577) (-6.1,-0.3535) (-6.05,-0.3439) (-6,-0.3291) (-5.95,-0.3094) (-5.9,-0.2851) (-5.85,-0.2566) (-5.8,-0.2244) (-5.75,-0.1888) (-5.7,-0.1506) (-5.65,-0.1103) (-5.6,-0.0683) (-5.55,-0.0255) (-5.5,0.0178) (-5.45,0.0608) (-5.4,0.1029) (-5.35,0.1437) (-5.3,0.1826) (-5.25,0.2190) (-5.2,0.2526) (-5.15,0.2829) (-5.1,0.3095) (-5.05,0.3322) (-5,0.3508) (-4.95,0.3649) (-4.9,0.3745) (-4.85,0.3796) (-4.8,0.3800) (-4.75,0.3759) (-4.7,0.3674) (-4.65,0.3545) (-4.6,0.3375) (-4.55,0.3166) (-4.5,0.2922) (-4.45,0.2644) (-4.4,0.2337) (-4.35,0.2004) (-4.3,0.1650) (-4.25,0.1278) (-4.2,0.0892) (-4.15,0.0497) (-4.1,0.0097) (-4.05,-0.0304) (-4,-0.0703) (-3.95,-0.1094) (-3.9,-0.1474) (-3.85,-0.1840) (-3.8,-0.2189) (-3.75,-0.2516) (-3.7,-0.2820) (-3.65,-0.3098) (-3.6,-0.3348) (-3.55,-0.3567) (-3.5,-0.3755) (-3.45,-0.3910) (-3.4,-0.4032) (-3.35,-0.4119) (-3.3,-0.4172) (-3.25,-0.4190) (-3.2,-0.4174) (-3.15,-0.4125) (-3.1,-0.4044) (-3.05,-0.3931) (-3,-0.3788) (-2.95,-0.3617) (-2.9,-0.3419) (-2.85,-0.3196) (-2.8,-0.2951) (-2.75,-0.2685) (-2.7,-0.2400) (-2.65,-0.2100) (-2.6,-0.1785) (-2.55,-0.1459) (-2.5,-0.1123) (-2.45,-0.0781) (-2.4,-0.0433) (-2.35,-0.0083) (-2.3,0.0267) (-2.25,0.0616) (-2.2,0.0961) (-2.15,0.1302) (-2.1,0.1635) (-2.05,0.1959) (-2,0.2274) (-1.95,0.2577) (-1.9,0.2868) (-1.85,0.3145) (-1.8,0.3408) (-1.75,0.3655) (-1.7,0.3886) (-1.65,0.4101) (-1.6,0.4299) (-1.55,0.4479) (-1.5,0.4643) (-1.45,0.4788) (-1.4,0.4917) (-1.35,0.5028) (-1.3,0.5123) (-1.25,0.5200) (-1.2,0.5262) (-1.15,0.5308) (-1.1,0.5338) (-1.05,0.5354) (-1,0.5356) (-0.95,0.5344) (-0.9,0.5320) (-0.85,0.5283) (-0.8,0.5236) (-0.75,0.5178) (-0.7,0.5110) (-0.65,0.5033) (-0.6,0.4948) (-0.55,0.4856) (-0.5,0.4757) (-0.45,0.4652) (-0.4,0.4542) (-0.35,0.4428) (-0.3,0.4309) (-0.25,0.4187) (-0.2,0.4063) (-0.15,0.3936) (-0.1,0.3808) (-0.05,0.3680) (0,0.3550) (0.05,0.3421) (0.1,0.3292) (0.15,0.3164) (0.2,0.3037) (0.25,0.2912) (0.3,0.2788) (0.35,0.2667) (0.4,0.2547) (0.45,0.2431) (0.5,0.2317) (0.55,0.2206) (0.6,0.2098) (0.65,0.1993) (0.7,0.1892) (0.75,0.1793) (0.8,0.1698) (0.85,0.1607) (0.9,0.1519) (0.95,0.1434) (1,0.1353) (1.05,0.1275) (1.1,0.1200) (1.15,0.1129) (1.2,0.1061) (1.25,0.0996) (1.3,0.0935) (1.35,0.0876) (1.4,0.0820) (1.45,0.0768) (1.5,0.0717) (1.55,0.0670) (1.6,0.0625) (1.65,0.0583) (1.7,0.0543) (1.75,0.0506) (1.8,0.0470) (1.85,0.0437) (1.9,0.0406) (1.95,0.0377) (2,0.0349) (2.05,0.0324) (2.1,0.0300) (2.15,0.0277) (2.2,0.0256) (2.25,0.0237) (2.3,0.0218) (2.35,0.0201) (2.4,0.0186) (2.45,0.0171) (2.5,0.0157)};
\addlegendentry{$\operatorname{Ai}$}
\addplot[black, dashed, line width=0.8pt, smooth] coordinates {(-10,-0.3147) (-9.95,-0.3048) (-9.9,-0.2874) (-9.85,-0.2629) (-9.8,-0.2319) (-9.75,-0.1952) (-9.7,-0.1538) (-9.65,-0.1087) (-9.6,-0.0609) (-9.55,-0.0117) (-9.5,0.0378) (-9.45,0.0864) (-9.4,0.1329) (-9.35,0.1764) (-9.3,0.2157) (-9.25,0.2500) (-9.2,0.2786) (-9.15,0.3007) (-9.1,0.3160) (-9.05,0.3242) (-9,0.3249) (-8.95,0.3184) (-8.9,0.3048) (-8.85,0.2844) (-8.8,0.2578) (-8.75,0.2255) (-8.7,0.1882) (-8.65,0.1469) (-8.6,0.1024) (-8.55,0.0556) (-8.5,0.0078) (-8.45,-0.0403) (-8.4,-0.0875) (-8.35,-0.1329) (-8.3,-0.1755) (-8.25,-0.2145) (-8.2,-0.2490) (-8.15,-0.2785) (-8.1,-0.3023) (-8.05,-0.3200) (-8,-0.3313) (-7.95,-0.3359) (-7.9,-0.3339) (-7.85,-0.3253) (-7.8,-0.3103) (-7.75,-0.2893) (-7.7,-0.2627) (-7.65,-0.2310) (-7.6,-0.1949) (-7.55,-0.1552) (-7.5,-0.1125) (-7.45,-0.0677) (-7.4,-0.0216) (-7.35,0.0249) (-7.3,0.0709) (-7.25,0.1156) (-7.2,0.1582) (-7.15,0.1980) (-7.1,0.2343) (-7.05,0.2663) (-7,0.2938) (-6.95,0.3160) (-6.9,0.3328) (-6.85,0.3439) (-6.8,0.3491) (-6.75,0.3483) (-6.7,0.3417) (-6.65,0.3294) (-6.6,0.3116) (-6.55,0.2887) (-6.5,0.2610) (-6.45,0.2291) (-6.4,0.1935) (-6.35,0.1549) (-6.3,0.1137) (-6.25,0.0708) (-6.2,0.0268) (-6.15,-0.0177) (-6.1,-0.0618) (-6.05,-0.1051) (-6,-0.1467) (-5.95,-0.1861) (-5.9,-0.2228) (-5.85,-0.2562) (-5.8,-0.2859) (-5.75,-0.3114) (-5.7,-0.3325) (-5.65,-0.3488) (-5.6,-0.3602) (-5.55,-0.3665) (-5.5,-0.3678) (-5.45,-0.3640) (-5.4,-0.3553) (-5.35,-0.3418) (-5.3,-0.3237) (-5.25,-0.3013) (-5.2,-0.2750) (-5.15,-0.2451) (-5.1,-0.2121) (-5.05,-0.1763) (-5,-0.1384) (-4.95,-0.0987) (-4.9,-0.0577) (-4.85,-0.0161) (-4.8,0.0257) (-4.75,0.0672) (-4.7,0.1079) (-4.65,0.1474) (-4.6,0.1851) (-4.55,0.2208) (-4.5,0.2539) (-4.45,0.2841) (-4.4,0.3112) (-4.35,0.3349) (-4.3,0.3549) (-4.25,0.3712) (-4.2,0.3835) (-4.15,0.3917) (-4.1,0.3959) (-4.05,0.3961) (-4,0.3922) (-3.95,0.3845) (-3.9,0.3729) (-3.85,0.3577) (-3.8,0.3390) (-3.75,0.3172) (-3.7,0.2924) (-3.65,0.2648) (-3.6,0.2349) (-3.55,0.2028) (-3.5,0.1689) (-3.45,0.1336) (-3.4,0.0971) (-3.35,0.0598) (-3.3,0.0220) (-3.25,-0.0160) (-3.2,-0.0539) (-3.15,-0.0913) (-3.1,-0.1281) (-3.05,-0.1638) (-3,-0.1983) (-2.95,-0.2313) (-2.9,-0.2626) (-2.85,-0.2920) (-2.8,-0.3193) (-2.75,-0.3444) (-2.7,-0.3671) (-2.65,-0.3873) (-2.6,-0.4050) (-2.55,-0.4201) (-2.5,-0.4324) (-2.45,-0.4421) (-2.4,-0.4491) (-2.35,-0.4533) (-2.3,-0.4549) (-2.25,-0.4539) (-2.2,-0.4504) (-2.15,-0.4443) (-2.1,-0.4359) (-2.05,-0.4252) (-2,-0.4123) (-1.95,-0.3974) (-1.9,-0.3805) (-1.85,-0.3618) (-1.8,-0.3414) (-1.75,-0.3195) (-1.7,-0.2962) (-1.65,-0.2716) (-1.6,-0.2460) (-1.55,-0.2193) (-1.5,-0.1918) (-1.45,-0.1636) (-1.4,-0.1347) (-1.35,-0.1054) (-1.3,-0.0758) (-1.25,-0.0459) (-1.2,-0.0158) (-1.15,0.0143) (-1.1,0.0443) (-1.05,0.0743) (-1,0.1040) (-0.95,0.1335) (-0.9,0.1626) (-0.85,0.1914) (-0.8,0.2198) (-0.75,0.2478) (-0.7,0.2753) (-0.65,0.3023) (-0.6,0.3288) (-0.55,0.3548) (-0.5,0.3804) (-0.45,0.4054) (-0.4,0.4300) (-0.35,0.4542) (-0.3,0.4780) (-0.25,0.5014) (-0.2,0.5245) (-0.15,0.5474) (-0.1,0.5700) (-0.05,0.5925) (0,0.6149) (0.05,0.6374) (0.1,0.6599) (0.15,0.6825) (0.2,0.7055) (0.25,0.7287) (0.3,0.7525) (0.35,0.7768) (0.4,0.8018) (0.45,0.8276) (0.5,0.8543) (0.55,0.8821) (0.6,0.9111) (0.65,0.9414) (0.7,0.9733) (0.75,1.0069) (0.8,1.0424) (0.85,1.0800) (0.9,1.1199) (0.95,1.1623) (1,1.2074) (1.05,1.2556) (1.1,1.3071) (1.15,1.3621) (1.2,1.4211) (1.25,1.4844) (1.3,1.5523) (1.35,1.6252) (1.4,1.7037) (1.45,1.7881) (1.5,1.8789) (1.55,1.9769) (1.6,2.0825) (1.65,2.1964) (1.7,2.3194) (1.75,2.4523) (1.8,2.5959) (1.85,2.7512) (1.9,2.9192) (1.95,3.1011) (2,3.2981)};
\addlegendentry{$\operatorname{Bi}$}
\end{axis}
\end{tikzpicture}}\hfill
\sdpanel{B}{\begin{tikzpicture}
\begin{axis}[sdplot, width=0.4\textwidth, height=0.42\textwidth, xmin=0, xmax=2.5, ymin=-1.2, ymax=6,
  xtick={0,0.5,1,1.5,2,2.5}, ytick={-1,0,1,2,3,4,5,6}, xlabel={$t$}, legend pos=north west]
\addplot[black, line width=0.3pt, domain=0:2.5, samples=2, forget plot] {0};
\addplot[black, densely dotted, line width=0.5pt, forget plot] coordinates {(2.0031,-1.2) (2.0031,6)};
\addplot[black, line width=0.8pt, smooth] coordinates {(0,1.0000) (0.025,1.0000) (0.05,1.0000) (0.075,1.0000) (0.1,1.0000) (0.125,1.0000) (0.15,1.0000) (0.175,0.9999) (0.2,0.9999) (0.225,0.9998) (0.25,0.9997) (0.275,0.9995) (0.3,0.9993) (0.325,0.9991) (0.35,0.9987) (0.375,0.9984) (0.4,0.9979) (0.425,0.9973) (0.45,0.9966) (0.475,0.9958) (0.5,0.9948) (0.525,0.9937) (0.55,0.9924) (0.575,0.9909) (0.6,0.9892) (0.625,0.9873) (0.65,0.9852) (0.675,0.9828) (0.7,0.9801) (0.725,0.9771) (0.75,0.9738) (0.775,0.9701) (0.8,0.9661) (0.825,0.9617) (0.85,0.9569) (0.875,0.9517) (0.9,0.9460) (0.925,0.9398) (0.95,0.9331) (0.975,0.9259) (1,0.9181) (1.025,0.9098) (1.05,0.9009) (1.075,0.8913) (1.1,0.8811) (1.125,0.8703) (1.15,0.8587) (1.175,0.8465) (1.2,0.8335) (1.225,0.8198) (1.25,0.8053) (1.275,0.7900) (1.3,0.7739) (1.325,0.7570) (1.35,0.7392) (1.375,0.7206) (1.4,0.7012) (1.425,0.6809) (1.45,0.6597) (1.475,0.6377) (1.5,0.6148) (1.525,0.5911) (1.55,0.5665) (1.575,0.5410) (1.6,0.5147) (1.625,0.4876) (1.65,0.4596) (1.675,0.4309) (1.7,0.4015) (1.725,0.3713) (1.75,0.3404) (1.775,0.3088) (1.8,0.2767) (1.825,0.2440) (1.85,0.2108) (1.875,0.1771) (1.9,0.1431) (1.925,0.1087) (1.95,0.0741) (1.975,0.0393) (2,0.0044) (2.025,-0.0305) (2.05,-0.0653) (2.075,-0.1000) (2.1,-0.1344) (2.125,-0.1684) (2.15,-0.2019) (2.175,-0.2349) (2.2,-0.2671) (2.225,-0.2986) (2.25,-0.3291) (2.275,-0.3586) (2.3,-0.3869) (2.325,-0.4140) (2.35,-0.4396) (2.375,-0.4638) (2.4,-0.4863) (2.425,-0.5070) (2.45,-0.5259) (2.475,-0.5428) (2.5,-0.5576)};
\addlegendentry{$w_1$}
\addplot[black, dashed, line width=0.8pt, smooth] coordinates {(0,0.0000) (0.025,0.0000) (0.05,0.0000) (0.075,0.0001) (0.1,0.0003) (0.125,0.0007) (0.15,0.0011) (0.175,0.0018) (0.2,0.0027) (0.225,0.0038) (0.25,0.0052) (0.275,0.0069) (0.3,0.0090) (0.325,0.0114) (0.35,0.0143) (0.375,0.0176) (0.4,0.0214) (0.425,0.0256) (0.45,0.0304) (0.475,0.0358) (0.5,0.0418) (0.525,0.0484) (0.55,0.0557) (0.575,0.0637) (0.6,0.0724) (0.625,0.0820) (0.65,0.0923) (0.675,0.1035) (0.7,0.1157) (0.725,0.1287) (0.75,0.1428) (0.775,0.1579) (0.8,0.1741) (0.825,0.1914) (0.85,0.2100) (0.875,0.2298) (0.9,0.2509) (0.925,0.2734) (0.95,0.2975) (0.975,0.3230) (1,0.3502) (1.025,0.3792) (1.05,0.4100) (1.075,0.4428) (1.1,0.4776) (1.125,0.5147) (1.15,0.5542) (1.175,0.5963) (1.2,0.6411) (1.225,0.6889) (1.25,0.7399) (1.275,0.7945) (1.3,0.8529) (1.325,0.9155) (1.35,0.9827) (1.375,1.0551) (1.4,1.1331) (1.425,1.2175) (1.45,1.3090) (1.475,1.4087) (1.5,1.5174) (1.525,1.6368) (1.55,1.7683) (1.575,1.9140) (1.6,2.0764) (1.625,2.2588) (1.65,2.4651) (1.675,2.7007) (1.7,2.9728) (1.725,3.2909) (1.75,3.6683) (1.775,4.1245) (1.8,4.6881) (1.825,5.4041) (1.85,6.3465) (1.875,7.6475) (1.9,9.5670) (1.925,12.6977) (1.95,18.7473) (1.975,35.4904)};
\addlegendentry{$y=-w_1'/w_1$}
\end{axis}
\end{tikzpicture}}
\caption{(A) The Airy functions $\operatorname{Ai}$ (solid) and $\operatorname{Bi}$ (dashed) of \cref{def:airy-functions}. Both oscillate for $x<0$; for $x>0$, $\operatorname{Ai}$ decreases to zero and $\operatorname{Bi}$ grows rapidly. (B) The solution $w_1$ of $w''+t^2w=0$ with $w_1(0)=1$, $w_1'(0)=0$ (solid) and the solution $y=-w_1'/w_1$ of the Riccati equation $y'=t^2+y^2$ with $y(0)=0$ (dashed), which blows up at the first positive zero $t^*=2.0031\ldots$ of $w_1$ (dotted line; \cref{prop:riccati-blowup}).\label{fig:airy-functions}}
\end{figure}

\subsection{The Riccati Equation of Chapter~11}

\Cref{ch:first-order-equations} could not solve the Riccati equation $y'=t^2+y^2$, and determined only numerically that its solution with $y(0)=0$ blows up at the time $2.003147359$. By \cref{prop:riccati-linear}, the substitution $y=-w'/w$ relates it to the linear equation
\begin{equation}
w''+t^2w=0 .
\label{eq:riccati-linear}
\end{equation}
In this subsection we keep the variable $t$ of \cref{ch:first-order-equations}. Every point is ordinary for \eqref{eq:riccati-linear}, and with $w=\sum a_nt^n$ the recurrence is $(n+2)(n+1)a_{n+2}=-a_{n-2}$, with $a_2=a_3=0$. So each coefficient determines the one four places later, and the solutions with $(w(0),w'(0))=(1,0)$ and $(0,1)$ are
\begin{equation}
\begin{gathered}
w_1(t)=\sum_{k=0}^\infty\frac{(-1)^k\,t^{4k}}{\prod_{j=1}^k(4j-1)\,4j}=1-\frac{t^4}{12}+\frac{t^8}{672}-\cdots,\\
w_2(t)=\sum_{k=0}^\infty\frac{(-1)^k\,t^{4k+1}}{\prod_{j=1}^k4j\,(4j+1)}=t-\frac{t^5}{20}+\cdots ,
\end{gathered}
\label{eq:riccati-series}
\end{equation}
both convergent on all of $\mathbb R$ by \cref{thm:ordinary-point}.

\begin{proposition}[The blow-up time]
\label{prop:riccati-blowup}
The function $w_1$ has a smallest positive zero $t^*$, and the maximal solution of $y'=t^2+y^2$ with $y(0)=0$ is $y=-w_1'/w_1$ on the interval $(-t^*,t^*)$. In particular it blows up at $t^*$.
\end{proposition}

\begin{proof}
On $(1,\infty)$ the coefficient $t^2$ exceeds $1$, so $w_1$ has a zero in $(2,2+\pi)$ by \cref{cor:zero-spacing}\,(i), and by \cref{lem:simple-zeros} it has a smallest positive zero $t^*$; since $w_1$ is even, $-t^*$ is its largest negative zero, and $w_1>0$ on $(-t^*,t^*)$, as $w_1(0)=1$. By \cref{prop:riccati-linear}, $\tilde y\coloneqq-w_1'/w_1$ solves the Riccati equation on $(-t^*,t^*)$, and $\tilde y(0)=-w_1'(0)/w_1(0)=0$. Let $y$ be the maximal solution, on $(\omega_-,\omega_+)$ (\cref{thm:maximal-solution}). By uniqueness (\cref{cor:local-uniqueness}), $y=\tilde y$ on the common interval. If $\omega_+<t^*$, then $y=\tilde y$ would stay bounded as $t\to\omega_+$, since $\tilde y$ is continuous at $\omega_+$, contradicting the blow-up alternative after \cref{thm:maximal-solution}; so $\omega_+\ge t^*$. As $t$ increases to $t^*$, $w_1(t)$ decreases to $0$ through positive values and $w_1'(t^*)<0$, because a zero of a nontrivial solution is simple (\cref{lem:simple-zeros}); so $\tilde y(t)\to+\infty$. A solution on an interval containing $t^*$ would be continuous at $t^*$, so $\omega_+\le t^*$, and $\omega_+=t^*$. The same argument, or the oddness of $y$ observed in \cref{sec:existence-uniqueness}, gives $\omega_-=-t^*$.
\end{proof}

The proposition turns the computation of the blow-up time into the computation of a zero of an everywhere convergent power series. The terms of $w_1(t)$ for $t$ near $2$ grow to about $16/12$ and then decrease rapidly, since the ratio of consecutive terms is $-t^4/((4k+3)(4k+4))$. Summing the series in ball arithmetic and bisecting gives
\[
t^*=2.003147359426884\ldots,
\]
whose first ten digits are the value obtained in \cref{sec:existence-uniqueness} with the Runge--Kutta method. \Cref{fig:airy-functions}B shows $w_1$ and the Riccati solution. \Cref{sec:bessel} identifies $w_1$ and $w_2$ with Bessel functions of the orders $\mp1/4$, as the note in \cref{sec:direction-fields} announced.

\section{Legendre's Equation}
\label{sec:legendre}

\subsection{The Series Solutions}

\term{Legendre's equation} of degree $\nu$ is
\begin{equation}
(1-x^2)\,y''-2x\,y'+\nu(\nu+1)\,y=0,\qquad\text{that is,}\qquad\bigl((1-x^2)\,y'\bigr)'+\nu(\nu+1)\,y=0 ,
\label{eq:legendre}
\end{equation}
with a real constant $\nu$. Since $\nu$ and $-1-\nu$ give the same product $\nu(\nu+1)$, we may assume $\nu\ge-1/2$. The equation arises when Laplace's equation is solved in spherical coordinates (see the note at the end of this section), with $x=\cos\theta$ for the polar angle $\theta$, so that the physically relevant interval is $[-1,1]$. Its singular points are $\pm1$ (\cref{prop:polynomial-ordinary}), and the point $0$ is ordinary. The second form of the equation, which follows from $((1-x^2)y')'=(1-x^2)y''-2xy'$, will give the orthogonality of its solutions.

\begin{proposition}[Series solutions of Legendre's equation]
\label{prop:legendre-series}
On $(-1,1)$, every solution of \eqref{eq:legendre} is the sum of a power series $\sum a_nx^n$ with radius of convergence at least $1$, whose coefficients satisfy
\begin{equation}
a_{n+2}=-\frac{(\nu-n)(\nu+n+1)}{(n+1)(n+2)}\,a_n\qquad(n\in\mathbb N).
\label{eq:legendre-recurrence}
\end{equation}
The solutions $y_{\mathrm e}$ with $a_0=1$, $a_1=0$ and $y_{\mathrm o}$ with $a_0=0$, $a_1=1$ are even and odd, respectively, and form a fundamental set on $(-1,1)$. If a series among them has infinitely many nonzero coefficients, its radius of convergence is exactly $1$.
\end{proposition}

\begin{proof}
On $(-1,1)$ the normal form has the coefficients $p(x)=-2x/(1-x^2)=-2\sum_kx^{2k+1}$ and $q(x)=\nu(\nu+1)\sum_kx^{2k}$, geometric series of radius $1$, so \cref{thm:ordinary-point} applies with $R=1$. Instead of using \eqref{eq:ordinary-recurrence} it is simpler to substitute $y=\sum a_nx^n$ into the original equation, whose coefficients are polynomials: the coefficient of $x^n$ in $(1-x^2)y''-2xy'+\nu(\nu+1)y$ is
\begin{align*}
&(n+2)(n+1)\,a_{n+2}-n(n-1)\,a_n-2n\,a_n+\nu(\nu+1)\,a_n\\
&\qquad=(n+2)(n+1)\,a_{n+2}+(\nu-n)(\nu+n+1)\,a_n ,
\end{align*}
since $n(n-1)+2n-\nu(\nu+1)=n(n+1)-\nu(\nu+1)=(n-\nu)(n+\nu+1)$. Setting it to zero gives \eqref{eq:legendre-recurrence}. The recurrence links only indices of the same parity, so $y_{\mathrm e}$ contains only even powers and $y_{\mathrm o}$ only odd ones; they form a fundamental set by \cref{thm:ordinary-point}\,(iii). If the even coefficients $a_{2k}$ are all nonzero, the ratios $|a_{2k+2}/a_{2k}|$ tend to $1$ by \eqref{eq:legendre-recurrence}, so the power series $\sum_ka_{2k}s^k$ in $s=x^2$ has radius $1$ (\vcref{1}{thm:cauchy-hadamard}), and $\sum a_{2k}x^{2k}$ diverges for $|x|>1$. The odd series is treated in the same way.
\end{proof}

The factor $\nu-n$ in the recurrence decides everything that follows. If $\nu$ is not an integer, no coefficient after the first vanishes, and both series have radius $1$. If $\nu=n$ is a natural number, the factor vanishes for the index $n$, and the series that contains the power $x^n$ stops there.

\subsection{Legendre Polynomials}

\begin{theorem}[Polynomial solutions]
\label{thm:legendre-polynomials}
Let $n\in\mathbb N$ and $\nu=n$. The polynomial solutions of \eqref{eq:legendre} form a space of dimension $1$, spanned by a polynomial of degree exactly $n$ that is even or odd as $n$ is, and every nonzero polynomial solution is nonzero at $x=1$. The \term{Legendre polynomial} $P_n$ is the polynomial solution with $P_n(1)=1$.
\end{theorem}

\begin{proof}
Take $n$ even; the odd case is the same with the roles of $y_{\mathrm e}$ and $y_{\mathrm o}$ exchanged. In \eqref{eq:legendre-recurrence}, the factor $(n-k)(n+k+1)$ vanishes only for $k=n$. For the even coefficients, $a_0=1$, the factors with $k=0,2,\ldots,n-2$ are nonzero, so $a_n\ne0$, and $a_{n+2}=0$, hence $a_{n+4}=a_{n+6}=\cdots=0$: $y_{\mathrm e}$ is a polynomial of degree $n$. For the odd coefficients no factor vanishes, so $y_{\mathrm o}$ has infinitely many nonzero coefficients. A solution $c_{\mathrm e}y_{\mathrm e}+c_{\mathrm o}y_{\mathrm o}$ has the coefficients of $c_{\mathrm e}y_{\mathrm e}$ in the even places and those of $c_{\mathrm o}y_{\mathrm o}$ in the odd places; if it is a polynomial, its power series about $0$ is the polynomial itself (\vcref{1}{cor:power-series-unique}) and has finitely many nonzero coefficients, so $c_{\mathrm o}=0$. Thus the polynomial solutions are the multiples of $y_{\mathrm e}$.

Let $p$ be a polynomial solution with $p(1)=0$. Differentiating the equation $k$ times by Leibniz's rule (\vcref{1}{thm:leibniz-rule}) gives
\[
(1-x^2)\,p^{(k+2)}-2(k+1)\,x\,p^{(k+1)}+(n-k)(n+k+1)\,p^{(k)}=0 ,
\]
since $D^k\bigl((1-x^2)p''\bigr)=(1-x^2)p^{(k+2)}-2kx\,p^{(k+1)}-k(k-1)\,p^{(k)}$ and $D^k(-2xp')=-2x\,p^{(k+1)}-2k\,p^{(k)}$. At $x=1$ the first term drops out, and $2(k+1)\,p^{(k+1)}(1)=(n-k)(n+k+1)\,p^{(k)}(1)$. By induction all derivatives of $p$ vanish at $1$, and since a polynomial equals its Taylor polynomial at $1$, $p=0$. So every nonzero polynomial solution is nonzero at $1$, and $P_n$ is well defined.
\end{proof}

The recurrence with $\nu=n$ and the normalisation give
\begin{gather*}
P_0=1,\quad P_1=x,\quad P_2=\frac{3x^2-1}2,\quad P_3=\frac{5x^3-3x}2,\\
P_4=\frac{35x^4-30x^2+3}8,\quad P_5=\frac{63x^5-70x^3+15x}8 ,
\end{gather*}
the polynomials that \cref{subsec:legendre} obtained by orthogonalising the monomials. For $P_4$, for instance, the recurrence starts from $a_0=1$ and gives $a_2=-(4\cdot5)/2=-10$ and $a_4=-(2\cdot7)/12\cdot a_2=35/3$, so the even solution is $1-10x^2+(35/3)x^4$, whose value at $1$ is $8/3$; dividing by it gives $P_4$. Since $P_n$ is even or odd, $P_n(-x)=(-1)^nP_n(x)$ and $P_n(-1)=(-1)^n$. \Cref{fig:legendre-functions}A shows $P_0,\ldots,P_5$. A closed formula for all of them was found by Rodrigues in 1816~\cite{rodrigues1816}.

\begin{theorem}[Rodrigues' formula]
\label{thm:rodrigues}
For every $n\in\mathbb N$,
\[
P_n(x)=\frac1{2^nn!}\,\frac{d^n}{dx^n}\,(x^2-1)^n=\frac1{2^n}\sum_{k=0}^{\lfloor n/2\rfloor}(-1)^k\binom nk\binom{2n-2k}nx^{n-2k} .
\]
In particular the leading coefficient of $P_n$ is $k_n\coloneqq(2n)!/(2^n(n!)^2)$.
\end{theorem}

\begin{proof}
Let $u\coloneqq(x^2-1)^n$ and $v\coloneqq u^{(n)}$, a polynomial of degree $n$. Since $u'=2nx(x^2-1)^{n-1}$, we have $(x^2-1)u'=2nxu$. Differentiating this identity $n+1$ times by Leibniz's rule, only the first three terms on the left and the first two on the right survive, because $x^2-1$ and $x$ have vanishing derivatives of orders above $2$ and $1$:
\[
(x^2-1)u^{(n+2)}+2(n+1)xu^{(n+1)}+n(n+1)u^{(n)}=2nxu^{(n+1)}+2n(n+1)u^{(n)} .
\]
Collecting terms gives $(x^2-1)v''+2xv'-n(n+1)v=0$, which is \eqref{eq:legendre} with $\nu=n$, multiplied by $-1$. So $v$ is a polynomial solution, and by \cref{thm:legendre-polynomials} a multiple of $P_n$. To find the factor, write $u=(x-1)^n(x+1)^n$ and apply Leibniz's rule to $v=D^n u$: every term except $\bigl(D^n(x-1)^n\bigr)(x+1)^n=n!\,(x+1)^n$ contains a factor $x-1$, so $v(1)=n!\,2^n$, and $P_n=v/(2^nn!)$. For the explicit sum, expand $u$ by the binomial theorem and differentiate $n$ times:
\begin{gather*}
u=\sum_{k=0}^n\binom nk(-1)^kx^{2n-2k},\\
D^nx^{2n-2k}=\frac{(2n-2k)!}{(n-2k)!}\,x^{n-2k}\quad(2k\le n),
\end{gather*}
while the powers with $2n-2k<n$ disappear. Dividing by $2^nn!$ gives the coefficients of the formula, since
\[
\frac1{2^nn!}\binom nk\frac{(2n-2k)!}{(n-2k)!}=\frac1{2^n}\binom nk\binom{2n-2k}n ,
\]
and the term $k=0$ gives the leading coefficient $(2n)!/(2^n(n!)^2)$.
\end{proof}

\subsection{Orthogonality and the Three-Term Recurrence}

\begin{theorem}[Orthogonality]
\label{thm:legendre-orthogonality}
For $m,n\in\mathbb N$,
\[
\int_{-1}^1P_m(x)\,P_n(x)\,dx=\begin{cases}0,&m\ne n,\\[2pt]\dfrac2{2n+1},&m=n .\end{cases}
\]
Consequently $P_n$ is orthogonal to every polynomial of degree less than $n$ in the inner product $\langle f,g\rangle=\int_{-1}^1fg$, and $P_n$ is the Legendre polynomial of \cref{subsec:legendre}, obtained there by the Gram--Schmidt process.
\end{theorem}

\begin{proof}
By the second form of \eqref{eq:legendre}, $((1-x^2)P_n')'=-n(n+1)P_n$, and the same with $m$. Multiply the equation for $P_n$ by $P_m$, the one for $P_m$ by $P_n$, and subtract:
\begin{align*}
\bigl(m(m+1)-n(n+1)\bigr)P_mP_n&=P_m\bigl((1-x^2)P_n'\bigr)'-P_n\bigl((1-x^2)P_m'\bigr)'\\
&=\Bigl((1-x^2)\bigl(P_mP_n'-P_nP_m'\bigr)\Bigr)' ,
\end{align*}
as one checks by the product rule. Integrating over $[-1,1]$, the right side gives zero, since $1-x^2$ vanishes at both ends. For $m\ne n$ the factor $m(m+1)-n(n+1)=(m-n)(m+n+1)$ is nonzero, so $\int_{-1}^1P_mP_n=0$.

For $m=n$ we use Rodrigues' formula. With $u=(x^2-1)^n$ as before, the derivatives $D^ju$ for $j<n$ contain the factor $(x^2-1)^{n-j}$ and vanish at $\pm1$. Integrating by parts $n$ times, each time moving one derivative from the first factor to the second, the boundary terms vanish and
\[
\int_{-1}^1(D^nu)^2\,dx=(-1)^n\int_{-1}^1u\,D^{2n}u\,dx=(2n)!\int_{-1}^1(1-x^2)^n\,dx ,
\]
since $D^{2n}u=(2n)!$, the polynomial $u$ having the leading term $x^{2n}$. The integral $I_n\coloneqq\int_{-1}^1(1-x^2)^ndx$ satisfies $I_0=2$ and, for $n\ge1$, $I_n=I_{n-1}-\int_{-1}^1x\cdot x(1-x^2)^{n-1}dx=I_{n-1}-I_n/(2n)$, by an integration by parts with the antiderivative $-(1-x^2)^n/(2n)$ of $x(1-x^2)^{n-1}$. Hence $I_n=2nI_{n-1}/(2n+1)$, and by induction $I_n=2^{2n+1}(n!)^2/(2n+1)!$. Therefore
\[
\int_{-1}^1P_n^2=\frac{(2n)!}{(2^nn!)^2}\cdot\frac{2^{2n+1}(n!)^2}{(2n+1)!}=\frac2{2n+1} .
\]
Since $P_0,\ldots,P_{n-1}$ have the degrees $0,\ldots,n-1$, they form a basis of the polynomials of degree less than $n$, and $P_n$ is orthogonal to all of them. The polynomials of degree at most $n$ that are orthogonal to those of lower degree form a subspace of dimension $(n+1)-n=1$ (\cref{thm:orthogonal-decomposition}), and the polynomial of \cref{subsec:legendre} lies in it, by the construction of the Gram--Schmidt process, and has the value $1$ at $x=1$; so it equals $P_n$.
\end{proof}

The proof of orthogonality used only that $P_m$ and $P_n$ solve equations of the form $((1-x^2)y')'+\lambda y=0$ with different values of $\lambda$, and that the factor $1-x^2$ vanishes at the ends. It is an instance of the self-adjointness that the note on Sturm--Liouville problems in \cref{sec:variable-coefficients} described: the operator $y\mapsto-((1-x^2)y')'$ is symmetric with respect to $\int_{-1}^1fg$ on the functions that are smooth on $[-1,1]$, and the Legendre polynomials are its eigenfunctions, with the eigenvalues $n(n+1)$. The same mechanism gives the orthogonality of the eigenvectors of a symmetric matrix (\cref{ch:symmetric-matrices}).

\begin{theorem}[Three-term recurrence]
\label{thm:legendre-recurrence}
For $n\ge1$,
\[
(n+1)\,P_{n+1}(x)=(2n+1)\,x\,P_n(x)-n\,P_{n-1}(x) .
\]
\end{theorem}

\begin{proof}
The polynomial $xP_n$ has degree $n+1$, so it is a combination of the orthogonal basis $P_0,\ldots,P_{n+1}$, with the coefficients $\langle P_k,xP_n\rangle/\|P_k\|^2$ (\cref{thm:orthonormal-expansion}, applied to the normalised basis). For $k\le n-2$, $\langle P_k,xP_n\rangle=\int_{-1}^1(xP_k)P_n=0$, since $xP_k$ has degree less than $n$ (\cref{thm:legendre-orthogonality}). For $k=n$, $\int_{-1}^1xP_n^2=0$, since the integrand is odd. Hence $xP_n=\alpha P_{n+1}+\gamma P_{n-1}$ with constants $\alpha$, $\gamma$. Comparing the coefficients of $x^{n+1}$ gives $k_n=\alpha k_{n+1}$, and by \cref{thm:rodrigues} $k_{n+1}/k_n=(2n+2)(2n+1)/(2(n+1)^2)=(2n+1)/(n+1)$, so $\alpha=(n+1)/(2n+1)$. At $x=1$, where every $P_k$ equals $1$, the identity gives $1=\alpha+\gamma$, so $\gamma=n/(2n+1)$. Multiplying by $2n+1$ gives the recurrence.
\end{proof}

The recurrence, which \cref{subsec:legendre} could only check for small $n$, computes $P_2,\ldots,P_n$ at a given point with about $5n$ arithmetic operations, without ever forming the coefficients of the polynomials, and \cref{sec:series-computation} shows that in floating-point arithmetic it is far more accurate than the evaluation of the coefficient formula. It is the case of the Legendre polynomials of the general three-term recurrence of orthogonal polynomials described in the note on three-term recurrences in \cref{subsec:legendre}; \cref{ch:difference-equations} studies such recurrences and their numerical stability in general.

\subsection{Solutions That Are Not Polynomials}

For $\nu=n$, the solutions other than the multiples of $P_n$ are given by the series that does not terminate, which has radius $1$. \Cref{sec:wronskian} found such a solution for $n=1$ by reduction of order, the function $Q_1(x)=(x/2)\ln((1+x)/(1-x))-1$. For $n=0$, the equation is $((1-x^2)y')'=0$, with the solutions $y'=c/(1-x^2)$, and the odd one with $y'(0)=1$ is
\[
Q_0(x)\coloneqq\frac12\ln\frac{1+x}{1-x}=x+\frac{x^3}3+\frac{x^5}5+\cdots ,
\]
in agreement with the recurrence \eqref{eq:legendre-recurrence}, which for $\nu=0$ reads $a_{n+2}=na_n/(n+2)$. For $\nu=1$ the even series is $1-x^2-x^4/3-x^6/5-\cdots$, which is $-Q_1$. The functions $Q_n$ are the \term{Legendre functions of the second kind}. Both $Q_0$ and $Q_1$ tend to $\pm\infty$ at the ends of the interval (\cref{fig:legendre-functions}B), and this is no accident.

\begin{figure}
\sdpanel{A}{\begin{tikzpicture}
\begin{axis}[sdplot, width=0.52\textwidth, height=0.42\textwidth, xmin=-1, xmax=1, ymin=-1.05, ymax=2.0,
  xtick={-1,-0.5,0,0.5,1}, ytick={-1,-0.5,0,0.5,1}, xlabel={$x$}, sdlegendtop, legend columns=3]
\addplot+[domain=-1:1, samples=2, no marks] {1};
\addlegendentry{$P_0$}
\addplot+[domain=-1:1, samples=2, no marks] {x};
\addlegendentry{$P_1$}
\addplot+[domain=-1:1, samples=81, no marks, smooth] {(3*x*x-1)/2};
\addlegendentry{$P_2$}
\addplot+[domain=-1:1, samples=81, no marks, smooth] {(5*x*x*x-3*x)/2};
\addlegendentry{$P_3$}
\addplot+[domain=-1:1, samples=121, no marks, smooth] {(35*x*x*x*x-30*x*x+3)/8};
\addlegendentry{$P_4$}
\addplot+[domain=-1:1, samples=121, no marks, smooth] {(63*x*x*x*x*x-70*x*x*x+15*x)/8};
\addlegendentry{$P_5$}
\end{axis}
\end{tikzpicture}}\hfill
\sdpanel{B}{\begin{tikzpicture}
\begin{axis}[sdplot, width=0.44\textwidth, height=0.42\textwidth, xmin=-1, xmax=1, ymin=-3, ymax=3,
  xtick={-1,-0.5,0,0.5,1}, ytick={-3,-2,-1,0,1,2,3}, xlabel={$x$}, legend pos=south east]
\addplot[black, line width=0.3pt, domain=-1:1, samples=2, forget plot] {0};
\addplot[black, line width=0.8pt, domain=-0.9975:0.9975, samples=301, smooth] {0.5*ln((1+x)/(1-x))};
\addlegendentry{$Q_0$}
\addplot[black, dashed, line width=0.8pt, domain=-0.9999:0.9999, samples=301, smooth] {0.5*x*ln((1+x)/(1-x))-1};
\addlegendentry{$Q_1$}
\end{axis}
\end{tikzpicture}}
\caption{Solutions of Legendre's equation. (A) The Legendre polynomials $P_0,\ldots,P_5$, normalised by $P_n(1)=1$; $P_n$ has $n$ zeros in $(-1,1)$, and $|P_n|\le1$ there. (B) The solutions $Q_0$ (solid) and $Q_1$ (dashed) for $\nu=0$ and $\nu=1$, which are not multiples of $P_0$ and $P_1$; they are unbounded at both ends of the interval, as \cref{prop:legendre-bounded} requires.\label{fig:legendre-functions}}
\end{figure}

\begin{proposition}[Bounded solutions]
\label{prop:legendre-bounded}
Let $\nu=n\in\mathbb N$. Every solution of \eqref{eq:legendre} on $(-1,1)$ that is not a multiple of $P_n$ is unbounded near $1$ and near $-1$. Thus the solutions bounded on $(-1,1)$ are exactly the multiples of $P_n$.
\end{proposition}

\begin{proof}
The polynomial $P_n$ has finitely many zeros and $P_n(1)=1$, so there is $b\in(0,1)$ with $P_n>0$ on $[b,1]$; let $M\coloneqq\max_{[b,1]}P_n^2$. On $(b,1)$ we apply reduction of order (\cref{prop:reduction-order}) to the normal form, whose coefficient $a_1(x)=-2x/(1-x^2)$ has the antiderivative $\ln(1-x^2)$; so $e^{-P}=(1-b^2)/(1-x^2)$ with $P(x)=\int_b^xa_1$, and
\[
y_2(x)\coloneqq P_n(x)\int_b^x\frac{1-b^2}{(1-s^2)\,P_n(s)^2}\,ds
\]
is a solution on $(b,1)$ that forms a fundamental set with $P_n$. Since $1-s^2\le2(1-s)$, the integrand is at least $(1-b^2)/(2M(1-s))$, whose integral from $b$ to $x$ is $((1-b^2)/(2M))\ln((1-b)/(1-x))$, which tends to $\infty$ as $x\to1^-$; and $P_n(x)\to1$. So $y_2(x)\to\infty$. A solution $y$ on $(-1,1)$ that is not a multiple of $P_n$ is, on $(b,1)$, a combination $c_1P_n+c_2y_2$; here $c_2\ne0$, since otherwise $y-c_1P_n$ would be a solution vanishing on $(b,1)$, hence everywhere (\cref{cor:higher-existence}). Since $P_n$ is bounded, $|y(x)|\to\infty$ as $x\to1^-$. Finally, the equation is unchanged by $x\mapsto-x$, so $x\mapsto y(-x)$ is also a solution, and it is not a multiple of $P_n(x)$, since $P_n(-x)=(-1)^nP_n(x)$. By what was just proved it is unbounded near $1$, so $y$ is unbounded near $-1$.
\end{proof}

When $\nu$ is not an integer, neither series terminates, and the conclusion is even stronger.

\begin{proposition}[Non-integral degree]
\label{prop:legendre-nonintegral}
If $\nu$ is not an integer, every nonzero solution of \eqref{eq:legendre} on $(-1,1)$ is unbounded.
\end{proposition}

\begin{proof}[sketch]
The points $\pm1$ are regular singular points in the sense of \cref{sec:frobenius}, each with the double indicial root $0$. By \cref{thm:frobenius}, near $x=1$ there is a solution $u$ given by a power series in $1-x$ with $u(1)=1$, and every solution that is not a multiple of $u$ contains a term $c\ln(1-x)$ with $c\ne0$ and is unbounded near $1$. The solution $u$ is the hypergeometric function $F(-\nu,\nu+1;1;(1-x)/2)$ of \cref{sec:classical-equations}. A hypergeometric series $F(a,b;a+b;s)$ behaves as $s\to1^-$ like $-\Gamma(a+b)\ln(1-s)/(\Gamma(a)\Gamma(b))$~\cite[Eq.~15.4.21]{olver2010nist}; here $a=-\nu$ and $b=\nu+1$, the factor $1/(\Gamma(-\nu)\Gamma(\nu+1))$ is not zero when $\nu$ is not an integer, and $u$ is unbounded near $x=-1$~\cite[Chap.~XV]{whittaker1927}, \cite[Chap.~7]{lebedev1972}. So no nonzero solution is bounded near both ends.
\end{proof}

The two propositions explain why the degree $\nu$ is an integer in the applications. When Laplace's equation is solved on a sphere, the polar angle $\theta$ runs over $[0,\pi]$, and $x=\cos\theta$ over the closed interval $[-1,1]$, including the poles $\theta=0$ and $\theta=\pi$, where a physical potential must be finite. A solution of Legendre's equation that is bounded on $(-1,1)$ exists only for $\nu\in\mathbb N$, and then it is a multiple of $P_n$. The condition of boundedness thus selects the eigenvalues $n(n+1)$ of the operator $-((1-x^2)y')'$, just as the boundary conditions of \cref{prop:boundary-value} selected the eigenvalues $n^2$ of $-y''$.

\begin{note}{Laplace's equation and separation of variables}
In spherical coordinates $(r,\theta,\varphi)$, with the polar angle $\theta$ measured from the positive $z$-axis as in the Notation, the Laplacian of a function $u$ that does not depend on $\varphi$ is
\[
\Delta u=\frac1{r^2}\frac{\partial}{\partial r}\Bigl(r^2\frac{\partial u}{\partial r}\Bigr)+\frac1{r^2\sin\theta}\frac{\partial}{\partial\theta}\Bigl(\sin\theta\frac{\partial u}{\partial\theta}\Bigr)
\]
\cite{courant1953}. A product $u=R(r)\Theta(\theta)$ is harmonic exactly when $(r^2R')'/R=-(\sin\theta\,\Theta')'/(\sin\theta\,\Theta)$, and since the left side depends only on $r$ and the right side only on $\theta$, both equal a constant $\lambda$. With $x=\cos\theta$ and $\Theta(\theta)=y(\cos\theta)$, the chain rule gives $(\sin\theta\,\Theta')'/\sin\theta=((1-x^2)y')'$, so $y$ solves Legendre's equation with $\nu(\nu+1)=\lambda$, and the radial equation $r^2R''+2rR'-\lambda R=0$ is an Euler--Cauchy equation (\cref{prop:euler-cauchy}), with the solutions $r^n$ and $r^{-n-1}$ for $\lambda=n(n+1)$. The harmonic functions $r^nP_n(\cos\theta)$ and $r^{-n-1}P_n(\cos\theta)$ are the building blocks of the potentials of gravitation and electrostatics with axial symmetry; Legendre introduced his polynomials in his work on the attraction of spheroids~\cite{legendre1785}. In cylindrical coordinates the same separation leads to Bessel's equation (\cref{sec:bessel}). Partial differential equations are not treated in this series; Courant and Hilbert~\cite{courant1953} develop the method of separation of variables.
\end{note}

\section{Regular Singular Points and the Method of Frobenius}
\label{sec:frobenius}

\subsection{Regular Singular Points}

At a singular point the theorem of \cref{sec:ordinary-points} says nothing, and the solutions need not be analytic there. The Euler--Cauchy equation $x^2y''+axy'+by=0$ of \cref{prop:euler-cauchy} shows what can happen: its solutions on $(0,\infty)$ are the powers $x^r$ for the roots $r$ of the indicial polynomial $r(r-1)+ar+b$, together with $x^r\ln x$ for a double root, and products of powers with $\cos(\beta\ln x)$ and $\sin(\beta\ln x)$ for complex roots. These functions are not analytic at $0$ unless $r\in\mathbb N$, but they are not wild either: they are powers, possibly multiplied by a logarithm. The equations whose solutions behave in this way near a singular point are those that look near it like an Euler--Cauchy equation.

\begin{definition}[Regular singular point]
\label{def:regular-singular}
A singular point $x_0$ of $y''+p(x)y'+q(x)y=0$ is a \term{regular singular point} if the functions $(x-x_0)\,p(x)$ and $(x-x_0)^2\,q(x)$, suitably defined at $x_0$, are analytic at $x_0$. Otherwise it is an \term{irregular singular point}.
\end{definition}

Take $x_0=0$. Multiplying the equation by $x^2$ and writing $P(x)\coloneqq xp(x)$ and $Q(x)\coloneqq x^2q(x)$, a regular singular point at $0$ means that the equation has the form
\begin{equation}
x^2y''+x\,P(x)\,y'+Q(x)\,y=0,\quad P(x)=\sum_{k=0}^\infty P_kx^k,\quad Q(x)=\sum_{k=0}^\infty Q_kx^k ,
\label{eq:frobenius-form}
\end{equation}
with power series $P$ and $Q$ of positive radius. If $P$ and $Q$ are constants, this is the Euler--Cauchy equation. In general, near $0$ the coefficients are close to $P_0$ and $Q_0$, and one may hope that the solutions are close to those of the Euler--Cauchy equation with $a=P_0$ and $b=Q_0$. This suggests the \term{indicial polynomial}
\begin{equation}
F(r)\coloneqq r(r-1)+P_0\,r+Q_0 ,
\label{eq:indicial}
\end{equation}
and its roots $r_1$ and $r_2$, the \term{indicial roots} or exponents at $0$. Since $F(r)=r^2+(P_0-1)r+Q_0$, they satisfy $r_1+r_2=1-P_0$.

Bessel's equation $x^2y''+xy'+(x^2-\nu^2)y=0$ is of the form \eqref{eq:frobenius-form} with $P=1$ and $Q(x)=x^2-\nu^2$, so $0$ is a regular singular point with $F(r)=r^2-\nu^2$ and the indicial roots $\pm\nu$. Legendre's equation has regular singular points at $\pm1$: near $x=1$, $(x-1)p(x)=2x/(1+x)$ and $(x-1)^2q(x)=-(x-1)\nu(\nu+1)/(1+x)$ are analytic at $1$, with the values $1$ and $0$ there, so the indicial polynomial at $1$ is $r(r-1)+r=r^2$, with the double root $0$. The equation $x^3y''+y=0$ has an irregular singular point at $0$, since $x^2q(x)=1/x$ is not analytic there.

\subsection{The Frobenius Series}

Following Frobenius~\cite{frobenius1873}, we look for solutions of \eqref{eq:frobenius-form} on $(0,R)$ of the form $x^r\sum_na_nx^n$, a power times a power series. Powers with real exponents are differentiated by $(x^r)'=rx^{r-1}$ for $x>0$ (\vcref{1}{thm:power-derivatives}), and \vcref{1}{thm:exponent-laws} gives $x^rx^s=x^{r+s}$. Let $L$ denote the operator on the left of \eqref{eq:frobenius-form}, $Ly\coloneqq x^2y''+xPy'+Qy$.

\begin{lemma}[Substituting a Frobenius series]
\label{lem:frobenius-substitution}
Let $P$ and $Q$ be as in \eqref{eq:frobenius-form} with radii of convergence at least $R>0$, let $A(x)=\sum_na_nx^n$ be a power series with radius at least $R$, and let $r$ be real. Then on $(0,R)$:
\begin{enumerate}
\item $L(x^rA)=x^r\sum_{n=0}^\infty c_nx^n$ with
\[
c_n\coloneqq F(r+n)\,a_n+\sum_{k=0}^{n-1}\bigl((r+k)\,P_{n-k}+Q_{n-k}\bigr)a_k ,
\]
and the series $\sum c_nx^n$ has radius at least $R$;
\item for every $u\in C^2(0,R)$, $L(u\ln x)=(\ln x)\,Lu+2xu'+(P-1)\,u$.
\end{enumerate}
\end{lemma}

\begin{proof}
(i) Let $y=x^rA$. By the product rule, $xy'=x^r(rA+xA')$ and $x^2y''=x^r\bigl(r(r-1)A+2rxA'+x^2A''\bigr)$, so
\[
Ly=x^r\bigl(x^2A''+2rxA'+r(r-1)A+P\,(rA+xA')+QA\bigr).
\]
The bracket is a power series with radius at least $R$, by \vcref{1}{thm:termwise-diff,prop:formal-analytic}. The coefficient of $x^n$ in $x^2A''+2rxA'+r(r-1)A$ is $\bigl(n(n-1)+2rn+r(r-1)\bigr)a_n=(r+n)(r+n-1)\,a_n$, and that in $P(rA+xA')+QA$ is $\sum_{k=0}^n\bigl(P_{n-k}(r+k)+Q_{n-k}\bigr)a_k$. The term $k=n$ of this sum is $(P_0(r+n)+Q_0)a_n$, and together with the first part it gives $F(r+n)\,a_n$. (ii) With $y=u\ln x$, $y'=u'\ln x+u/x$ and $y''=u''\ln x+2u'/x-u/x^2$, so $x^2y''=x^2u''\ln x+2xu'-u$, $xPy'=xPu'\ln x+Pu$, and adding $Qu\ln x$ gives the claim.
\end{proof}

By part (i), $x^rA$ solves \eqref{eq:frobenius-form} on $(0,R)$ if and only if all $c_n$ vanish (\vcref{1}{cor:power-series-unique}). With $a_0\ne0$, the condition $c_0=F(r)a_0=0$ says that $r$ must be an indicial root, and the conditions $c_n=0$ for $n\ge1$ form the recurrence
\begin{equation}
F(r+n)\,a_n=-\sum_{k=0}^{n-1}\bigl((r+k)\,P_{n-k}+Q_{n-k}\bigr)a_k\qquad(n\ge1).
\label{eq:frobenius-recurrence}
\end{equation}
It determines $a_n$ from the earlier coefficients as long as $F(r+n)\ne0$. Since $F(s)=(s-r_1)(s-r_2)$, this can fail only if $r+n$ is the other indicial root, that is, if the two roots differ by the positive integer $n$. For the larger root $r_1$ it never fails, since $r_1+n$ exceeds both roots; for the smaller root $r_2$ it fails exactly when $r_1-r_2$ is a positive integer. The case of a double root is also special, because then the recurrence produces only one solution. These are the three cases of the theorem.

\subsection{The Theorem of Frobenius}

\begin{theorem}[Frobenius]
\label{thm:frobenius}
Let $P$ and $Q$ be power series with radii of convergence at least $R>0$, and let the indicial roots $r_1\ge r_2$ of \eqref{eq:indicial} be real. On $(0,R)$, the equation \eqref{eq:frobenius-form} has the following fundamental sets of solutions, in which $A(x)=\sum_na_nx^n$ and $B(x)=\sum_nb_nx^n$ are power series with radius of convergence at least $R$ and $a_0=1$.
\begin{enumerate}
\item In all cases there is the solution $y_1=x^{r_1}A(x)$, whose coefficients are given by \eqref{eq:frobenius-recurrence} with $r=r_1$.
\item If $r_1-r_2\notin\mathbb N$: $y_1$ and $y_2=x^{r_2}B(x)$, with $b_0=1$ and the coefficients given by \eqref{eq:frobenius-recurrence} with $r=r_2$.
\item If $r_1=r_2$: $y_1$ and $y_2=y_1\ln x+x^{r_1}B(x)$, with $b_0=0$.
\item If $r_1-r_2=N$ is a positive integer: $y_1$ and $y_2=Cy_1\ln x+x^{r_2}B(x)$, with $b_0=1$ and a constant $C$, which may or may not be zero.
\end{enumerate}
\end{theorem}

\begin{proof}
Fix $\rho$ with $0<\rho<R$, let $M$ bound the sequences $(|P_k|\rho^k)$ and $(|Q_k|\rho^k)$, and fix $\lambda>1$. Let $d\coloneqq r_1-r_2\ge0$.

(i) and (ii). Let $r$ be $r_1$, or $r_2$ in case (ii). Then $F(r+n)=n(n+d)$ for $r=r_1$ and $F(r+n)=n(n-d)$ for $r=r_2$, and in both cases $F(r+n)\ne0$ for $n\ge1$, so \eqref{eq:frobenius-recurrence} defines the coefficients $a_n$ (or $b_n$) from $a_0=1$. Let $\gamma_n\coloneqq|F(r+n)|$; then $\gamma_n/n=|n\pm d|\to\infty$. Since $|(r+k)P_{n-k}+Q_{n-k}|\le(|r|+1+k)M\rho^{k-n}$, the numbers $b_n\coloneqq|a_n|\rho^n$ satisfy
\[
\gamma_n\,b_n\le\sum_{k=0}^{n-1}M(|r|+1+k)\,b_k\qquad(n\ge1),
\]
which is \eqref{eq:majorant-inequality} with $\alpha=M(|r|+1)$, $\beta=M$ and $D=0$. By \cref{lem:coefficient-growth}, $|a_n|\rho^n\le K\lambda^n$ with a constant $K$, and since $\rho<R$ and $\lambda>1$ are arbitrary, $A$ has radius at least $R$. By \cref{lem:frobenius-substitution}\,(i), $L(x^rA)=x^r\sum c_nx^n$ with $c_0=F(r)=0$ and $c_n=0$ for $n\ge1$ by the recurrence; so $x^rA$ is a solution on $(0,R)$.

For (iii) and (iv) we need the function $G$ defined by $2xy_1'+(P-1)y_1=x^{r_1}G(x)$. By the computation in the proof of \cref{lem:frobenius-substitution}, $xy_1'=x^{r_1}(r_1A+xA')$, so $G=2r_1A+2xA'+(P-1)A$, a power series with radius at least $R$ and coefficients
\[
g_n=\bigl(2(r_1+n)-1\bigr)a_n+\sum_{k=0}^nP_{n-k}\,a_k ,\qquad g_0=2r_1-1+P_0=r_1-r_2=d,
\]
using $r_1+r_2=1-P_0$. From $|a_k|\rho^k\le K\lambda^k$, with the constant $K$ of (i) for $r=r_1$, we get $|g_n|\rho^n\le(2|r_1|+2n+1)K\lambda^n+\sum_{k\le n}MK\lambda^k\le(\alpha'+\beta'n)\lambda^n$ with constants $\alpha'$, $\beta'$, since $\sum_{k\le n}\lambda^k\le\lambda^{n+1}/(\lambda-1)$. By \cref{lem:frobenius-substitution}\,(ii), for every constant $C'$ and every series $B$ with radius at least $R$,
\begin{equation}
L\bigl(C'y_1\ln x+x^{s}B\bigr)=C'x^{r_1}G+L(x^sB),
\label{eq:log-solution}
\end{equation}
since $Ly_1=0$.

(iii) Let $r_1=r_2$, so that $d=0$, $g_0=0$ and $F(r_1+n)=n^2$. Take $C'=1$ and $s=r_1$ in \eqref{eq:log-solution}. By \cref{lem:frobenius-substitution}\,(i), the right side is $x^{r_1}\sum_n\bigl(g_n+F(r_1+n)b_n+\sum_{k<n}((r_1+k)P_{n-k}+Q_{n-k})b_k\bigr)x^n$. For $n=0$ the coefficient is $g_0+F(r_1)b_0=0$, and for $n\ge1$ it vanishes if we define
\[
b_n\coloneqq-\frac1{n^2}\Bigl(g_n+\sum_{k=0}^{n-1}\bigl((r_1+k)P_{n-k}+Q_{n-k}\bigr)b_k\Bigr),\qquad b_0\coloneqq0 .
\]
The numbers $|b_n|\rho^n$ satisfy \eqref{eq:majorant-inequality} with $\gamma_n=n^2$, $\alpha=\max(M(|r_1|+1),\alpha')$, $\beta=\max(M,\beta')$ and $D=1$, so $B$ has radius at least $R$ by \cref{lem:coefficient-growth}, and $y_2=y_1\ln x+x^{r_1}B$ is a solution on $(0,R)$.

(iv) Let $d=N\ge1$, so that $r_2=r_1-N$, $g_0=N$ and $F(r_2+n)=n(n-N)$. Take $s=r_2$ and $C'=C$ in \eqref{eq:log-solution}, and write $x^{r_1}G=x^{r_2}\sum_ng_{n-N}x^n$ with $g_j\coloneqq0$ for $j<0$. The coefficient of $x^{r_2+n}$ is
\[
n(n-N)\,b_n+\sum_{k=0}^{n-1}\bigl((r_2+k)P_{n-k}+Q_{n-k}\bigr)b_k+C\,g_{n-N}.
\]
For $n=0$ it vanishes with $b_0=1$. For $1\le n<N$, $g_{n-N}=0$ and $n(n-N)\ne0$, so it vanishes for a unique $b_n$. For $n=N$ the first term is $0$ whatever $b_N$ is, and the coefficient is $\sum_{k<N}((r_2+k)P_{N-k}+Q_{N-k})b_k+CN$, which vanishes for exactly one value of $C$, namely
\[
C=-\frac1N\sum_{k=0}^{N-1}\bigl((r_2+k)P_{N-k}+Q_{N-k}\bigr)b_k ;
\]
we take this $C$, and $b_N\coloneqq0$. For $n>N$, $n(n-N)\ne0$ again determines $b_n$. With $\gamma_n=n(n-N)$ for $n>N$, the numbers $|b_n|\rho^n$ satisfy \eqref{eq:majorant-inequality} for $n>N$ with $\alpha=\max(M(|r_2|+1),\alpha')$, $\beta=\max(M,\beta')$ and $D=|C|\rho^N$, because $|g_{n-N}|\rho^n\le\rho^N(\alpha'+\beta'n)\lambda^n$; so $B$ has radius at least $R$, and $y_2$ is a solution.

Finally, in each case $y_1$ and $y_2$ are linearly independent. Suppose $c_1y_1+c_2y_2=0$ on $(0,R)$. In case (ii), dividing by $x^{r_2}$ gives $c_1x^dA(x)+c_2B(x)=0$, and letting $x\to0^+$, where $x^d\to0$ because $d>0$, gives $c_2B(0)=c_2=0$. In case (iv), dividing by $x^{r_2}$ gives $(c_1+c_2C\ln x)\,x^NA(x)+c_2B(x)=0$, and since $x^N\ln x\to0$ as $x\to0^+$ (\vcref{1}{prop:ln-growth}), again $c_2=0$. In case (iii), dividing by $x^{r_1}$ gives $(c_1+c_2\ln x)A(x)+c_2B(x)=0$; if $c_2\ne0$, then $\ln x=-(c_1A(x)+c_2B(x))/(c_2A(x))$ for small $x>0$, where $A(x)\ne0$ because $A(0)=1$, and the right side has a finite limit as $x\to0^+$ while $\ln x\to-\infty$; so $c_2=0$. In all cases $c_1y_1=0$ then forces $c_1=0$, since $y_1=x^{r_1}A$ with $A(0)=1$ is not the zero function. Two independent solutions of an equation of the second order form a fundamental set (\cref{thm:solution-space}), applied to the normal form, whose coefficients $P(x)/x$ and $Q(x)/x^2$ are continuous on $(0,R)$.
\end{proof}

The proof is constructive: in each case it gives a recurrence for the coefficients of the second solution. In case (iii) the solution contains a logarithm unavoidably, since $y_1\ln x$ appears with the coefficient $1$. In case (iv), the logarithm appears exactly when the constant $C$ is nonzero; the examples below show that both possibilities occur. The second solution can also be obtained by reduction of order (\cref{prop:reduction-order}). The normal form has the coefficient $a_1=P(x)/x$, and $e^{-\int a_1}=x^{-P_0}\exp\bigl(-\int(P-P_0)/x\bigr)$, where $(P-P_0)/x$ is a power series; so
\[
y_2=y_1\int\frac{x^{-P_0}\exp\bigl(-\int(P(x)-P_0)/x\,dx\bigr)}{x^{2r_1}A(x)^2}\,dx ,
\]
and since $2r_1+P_0=1+d$, the integrand is $x^{-1-d}$ times a power series. Integrating term by term gives powers $x^{k-d}$, except that a term with $x^{-1}$, which occurs only when $d$ is a natural number, integrates to a logarithm. This explains the forms of the theorem. Frobenius himself obtained the logarithmic solutions in another way, by differentiating the series $x^r\sum a_n(r)x^n$ with respect to the parameter $r$~\cite{frobenius1873}, \cite[Chap.~XVI]{ince1926}.

\subsection{Complex Roots and Negative \texorpdfstring{$x$}{x}}

If $P$ and $Q$ are real but $F$ has complex roots $\gamma\pm i\beta$ with $\beta>0$, the difference of the roots is not an integer, and the proof of case (ii) applies verbatim with complex exponents, the modulus replacing the absolute value: there are two complex solutions $x^{\gamma\pm i\beta}\sum_na_n^\pm x^n$, where $x^{\gamma+i\beta}\coloneqq x^\gamma e^{i\beta\ln x}=x^\gamma\bigl(\cos(\beta\ln x)+i\sin(\beta\ln x)\bigr)$ and the derivative $(x^{s})'=sx^{s-1}$ holds for complex $s$ by \cref{lem:complex-calculus}. The coefficients of the second solution are the complex conjugates of those of the first, and the real and imaginary parts of the first solution form a real fundamental set, as for the Euler--Cauchy equation $x^2y''+xy'+y=0$ in \cref{sec:constant-coefficients}, with the solutions $\cos(\ln x)$ and $\sin(\ln x)$. Such solutions oscillate infinitely often as $x\to0^+$.

On $(-R,0)$, the substitution $x=-s$ transforms \eqref{eq:frobenius-form} into an equation of the same form in $s$, with the coefficients $P(-s)$ and $Q(-s)$ and the same indicial polynomial, so the theorem holds with $|x|^{r}$ and $\ln|x|$ in place of $x^r$ and $\ln x$.

\subsection{Examples}

The equation
\[
2x^2y''-xy'+(1+x)\,y=0
\]
has, after division by $2$, the form \eqref{eq:frobenius-form} with $P=-1/2$ and $Q(x)=(1+x)/2$, polynomials with $R=\infty$. The indicial polynomial is $F(r)=r(r-1)-r/2+1/2=(r-1)(r-1/2)$, with the roots $r_1=1$ and $r_2=1/2$, whose difference is not an integer. Only $Q_1=1/2$ is nonzero among the higher coefficients, so \eqref{eq:frobenius-recurrence} reads $F(r+n)a_n=-a_{n-1}/2$, that is,
\[
a_n=-\frac{a_{n-1}}{n(2n+1)}\quad(r=1),\qquad b_n=-\frac{b_{n-1}}{n(2n-1)}\quad(r=\tfrac12).
\]
With $a_0=b_0=1$, induction gives $a_n=(-1)^n2^n/(2n+1)!$ and $b_n=(-1)^n2^n/(2n)!$, because $2n(2n+1)\,(2n-1)!=(2n+1)!$ and $2n(2n-1)\,(2n-2)!=(2n)!$. The two solutions are therefore
\begin{gather*}
y_1=x\sum_{n=0}^\infty\frac{(-1)^n(2x)^n}{(2n+1)!}=\sqrt{\frac x2}\,\sin\sqrt{2x},\\
y_2=\sqrt x\sum_{n=0}^\infty\frac{(-1)^n(2x)^n}{(2n)!}=\sqrt x\,\cos\sqrt{2x},
\end{gather*}
by the series of the sine and the cosine with the argument $\sqrt{2x}$. The first is $x$ times a power series, since its exponent $r_1=1$ is a natural number, and so extends to a function analytic at $0$. The second is not analytic at $0$: it tends to $0$ as $x\to0^+$, but its derivative tends to $\infty$.

The equation $xy''+y=0$, or $x^2y''+x\cdot y=0$ after multiplication by $x$, has $P=0$ and $Q(x)=x$, so $F(r)=r(r-1)$, with the roots $1$ and $0$, which differ by $N=1$. The recurrence for $r_1=1$ is $(n+1)n\,a_n=-a_{n-1}$, so
\[
y_1=\sum_{n=0}^\infty\frac{(-1)^n\,x^{n+1}}{n!\,(n+1)!}=x-\frac{x^2}2+\frac{x^3}{12}-\frac{x^4}{144}+\cdots .
\]
For the second solution, case (iv) of the proof gives $C=-Q_1b_0/N=-1$, so the logarithm is present. With $g_n=(2n+1)a_n$ (here $P=0$ and $r_1=1$) and $b_1=0$, the coefficients satisfy $n(n-1)b_n=-b_{n-1}+g_{n-1}$ for $n\ge2$, which gives $b_2=-3/4$ and $b_3=7/36$:
\[
y_2=-y_1\ln x+1-\frac34x^2+\frac7{36}x^3-\cdots .
\]
\Cref{sec:bessel} shows that $y_1=\sqrt x\,J_1(2\sqrt x)$, a Bessel function. For Bessel's equation of order $1/2$, on the other hand, the indicial roots $\pm1/2$ also differ by $N=1$, but there $C=0$ and the second solution is a pure Frobenius series (\cref{prop:bessel-half}). Legendre's equation at $x=1$ illustrates case (iii): the double root $0$ gives one solution analytic at $1$ and another containing $\ln(1-x)$, as $Q_0=(\ln(1+x)-\ln(1-x))/2$ does for $\nu=0$.

\subsection{Irregular Singular Points}

At an irregular singular point the method can fail completely. The simplest example is of the first order. Substituting $y=\sum a_nx^n$ into
\[
x^2y'=y-x
\]
gives $a_0=0$, $a_1=1$ and $a_{n+1}=na_n$ for $n\ge1$, so that $a_n=(n-1)!$ and the formal solution is
\[
\sum_{n=0}^\infty n!\,x^{n+1}=x+x^2+2x^3+6x^4+\cdots ,
\]
whose radius of convergence is $0$. The equation is a perfectly good linear equation of the first order on $(-\infty,0)$, where its normal form $y'-y/x^2=-1/x$ has continuous coefficients, and its solutions exist by \cref{thm:linear-first-order}. The divergent series is related to one of them. With Euler's function $F(s)=\int_0^\infty e^{-t}/(1+st)\,dt$ of \vcref{1}{sec:asymptotic-expansions}, differentiation under the integral sign (\vcref{1}{thm:parameter-improper}) gives $s^2F'(s)=-\int_0^\infty s^2te^{-t}(1+st)^{-2}\,dt$ for $s>0$. Since $\partial_t(1+st)^{-1}=-s(1+st)^{-2}$, an integration by parts in $t$, whose boundary terms vanish, turns this into $s^2F'(s)=-s\int_0^\infty(1-t)e^{-t}/(1+st)\,dt$, so that
\[
s^2F'(s)+(1+s)F(s)=\int_0^\infty\frac{(1+st)\,e^{-t}}{1+st}\,dt=1\qquad(s>0),
\]
and a short computation shows that $y(x)\coloneqq xF(-x)$ solves $x^2y'=y-x$ for $x<0$. By \vcref{1}{prop:euler-series}, $F(s)\sim\sum_k(-1)^kk!\,s^k$ as $s\to0^+$, so $y(x)\sim\sum_kk!\,x^{k+1}$ as $x\to0^-$: the divergent series is the asymptotic expansion of a true solution. This is typical of irregular singular points, and Bessel's and Airy's equations both have one at infinity, as the next subsection shows.

\begin{note}{When divergent series are useful}
The partial sums of $\sum_kk!\,x^{k+1}$ approximate the solution $xF(-x)$ the better the smaller $|x|$ is, as long as one stops near the smallest term, as \vcref{1}{fig:euler-series} showed for $F$; summing further makes things worse. Such asymptotic expansions describe the solutions of equations near an irregular singular point, and they are often the best way to compute special functions where the convergent series are useless, for example the Bessel functions for large $x$ (\cref{thm:bessel-asymptotic}). Their theory is the subject of Olver's book~\cite{olver1997asymptotics} and of the first chapters of Volume~XI.
\end{note}

\subsection{The Point at Infinity}

The behaviour of solutions for large $|x|$ can be studied by the substitution $x=1/s$, which maps a neighbourhood of infinity to a neighbourhood of $s=0$. If $y(x)=Y(1/x)$, then $y'=-s^2Y'$ and $y''=s^4Y''+2s^3Y'$, where $Y$ is evaluated at $s=1/x$, and $y''+p(x)y'+q(x)y=0$ becomes
\[
Y''+\Bigl(\frac2s-\frac{p(1/s)}{s^2}\Bigr)Y'+\frac{q(1/s)}{s^4}\,Y=0 .
\]
The point at infinity is called ordinary, regular singular or irregular singular according as $s=0$ is for this equation. For Legendre's equation, $p(1/s)=-2s/(s^2-1)$ and $q(1/s)=\nu(\nu+1)s^2/(s^2-1)$, so that the new coefficients are $2s/(s^2-1)$ and $\nu(\nu+1)/(s^2(s^2-1))$, and infinity is a regular singular point, except for $\nu=0$, when it is an ordinary point. For Airy's equation the new coefficient of $Y$ is $-1/s^5$, and for Bessel's equation it is $(1-\nu^2s^2)/s^4$; in both cases $s^2$ times it is not analytic at $0$, and infinity is an irregular singular point. This is why the solutions of these equations have no simple form at infinity, and why their behaviour there is described by asymptotic expansions rather than by convergent series.

\section{Bessel's Equation}
\label{sec:bessel}

\term{Bessel's equation} of order $\nu$ is
\begin{equation}
x^2y''+x\,y'+(x^2-\nu^2)\,y=0 ,
\label{eq:bessel}
\end{equation}
with a real constant $\nu\ge0$. It arises whenever a problem with circular or cylindrical symmetry is solved by separation of variables. For the vibrations of a circular drum of radius $a$, whose displacement $u(r,\varphi,t)$ obeys the wave equation $u_{tt}=c^2\Delta u$, the Laplacian in polar coordinates, $u_{rr}+u_r/r+u_{\varphi\varphi}/r^2$ (\vcref{1}{sec:chain-rule-n}), and the trial solution $u=R(r)\cos(m\varphi)\cos(\omega t)$ give $R''+R'/r+(\omega^2/c^2-m^2/r^2)R=0$, which becomes \eqref{eq:bessel} with $\nu=m$ in the variable $x=\omega r/c$. We return to the drum at the end of the section.

As \cref{sec:frobenius} showed, $0$ is a regular singular point of \eqref{eq:bessel}, with $P=1$, $Q(x)=x^2-\nu^2$ and the indicial roots $\pm\nu$, and infinity is an irregular singular point. Since $Q_2=1$ is the only higher coefficient of $P$ and $Q$ that is not zero, the recurrence \eqref{eq:frobenius-recurrence} has just two terms:
\[
(r+n-\nu)(r+n+\nu)\,a_n=-a_{n-2}\qquad(n\ge1),
\]
with $a_{-1}\coloneqq0$. For $r=\nu$ it reads $n(n+2\nu)a_n=-a_{n-2}$; it gives $a_1=0$, so all odd coefficients vanish, and $a_{2m}=-a_{2m-2}/(4m(m+\nu))$. Iterating, $a_{2m}=(-1)^ma_0/(4^mm!\,(\nu+1)(\nu+2)\cdots(\nu+m))$. The product in the denominator is $\Gamma(\nu+m+1)/\Gamma(\nu+1)$, by the functional equation $\Gamma(s+1)=s\Gamma(s)$ (\vcref{1}{thm:gamma}), and the choice $a_0=1/(2^\nu\Gamma(\nu+1))$ makes the formula symmetric. For the other root $-\nu$ the same formula with $-\nu$ in place of $\nu$ requires the gamma function at negative arguments.

\subsection{The Gamma Function for Negative Arguments}

\Vcref{1}{def:gamma} defined $\Gamma(s)$ by an integral that converges only for $s>0$. The functional equation $\Gamma(s)=\Gamma(s+1)/s$ extends it.

\begin{definition}[Gamma function for negative arguments]
\label{def:gamma-extension}
For real $s<0$ that is not an integer, let $k\in\mathbb N$ be such that $s+k>0$, and
\[
\Gamma(s)\coloneqq\frac{\Gamma(s+k)}{s(s+1)\cdots(s+k-1)} .
\]
For all real $s$ let $1/\Gamma(s)$ denote the reciprocal of $\Gamma(s)$ if $s\notin\{0,-1,-2,\ldots\}$, and $0$ if $s\in\{0,-1,-2,\ldots\}$.
\end{definition}

The definition does not depend on $k$: replacing $k$ by $k+1$ multiplies the numerator by $s+k$, by \vcref{1}{thm:gamma}, and the denominator by the same factor. The extended function satisfies $\Gamma(s+1)=s\Gamma(s)$ for all real $s\notin\{0,-1,-2,\ldots\}$, and it has no zeros, since $\Gamma>0$ on $(0,\infty)$. As $s$ approaches a nonpositive integer, $|\Gamma(s)|\to\infty$, so the convention for $1/\Gamma$ makes $1/\Gamma$ continuous, and it gives the identity
\begin{equation}
\frac1{\Gamma(s)}=\frac s{\Gamma(s+1)}\qquad\text{for all real }s,
\label{eq:rgamma}
\end{equation}
which at $s=0,-1,-2,\ldots$ says $0=0$. For example $\Gamma(-1/2)=\Gamma(1/2)/(-1/2)=-2\sqrt\pi$, by \vcref{1}{cor:gaussian-consequences}. Volume~III extends $\Gamma$ to complex arguments, and Volume~XI computes it.

\subsection{Bessel Functions of the First Kind}

\begin{theorem}[Bessel functions of the first kind]
\label{thm:bessel-j}
For every real $\nu$ and $x>0$ the series
\begin{equation}
J_\nu(x)\coloneqq\sum_{m=0}^\infty\frac{(-1)^m}{m!\,\Gamma(m+\nu+1)}\Bigl(\frac x2\Bigr)^{2m+\nu}
\label{eq:bessel-j}
\end{equation}
converges, and $J_\nu$ solves Bessel's equation \eqref{eq:bessel} of order $|\nu|$ on $(0,\infty)$. More precisely, $x^{-\nu}J_\nu(x)$ is the sum of a power series in $x^2$ with radius of convergence $\infty$. For $\nu\ge0$, $J_\nu$ is the solution $y_1$ of \cref{thm:frobenius}\,(i), multiplied by $1/(2^\nu\Gamma(\nu+1))$.
\end{theorem}

\begin{proof}
Write $J_\nu(x)=x^\nu\sum_md_mx^{2m}$ with $d_m\coloneqq(-1)^m/(2^{2m+\nu}m!\,\Gamma(m+\nu+1))$. Only finitely many $d_m$ vanish, those with $m+\nu+1\in\{0,-1,\ldots\}$, and for the others $|d_{m+1}/d_m|=1/(4(m+1)|m+\nu+1|)\to0$, by $\Gamma(s+1)=s\Gamma(s)$. So the power series $\sum d_ms^m$ has radius $\infty$ (\vcref{1}{thm:cauchy-hadamard}), and so has $\sum d_mx^{2m}$. To see that $J_\nu$ is a solution, apply \cref{lem:frobenius-substitution}\,(i) with $r=\nu$, $P=1$, $Q=x^2-\nu^2$ and the coefficients $a_{2m}=d_m$, $a_{2m+1}=0$. Here $F(s)=s^2-\nu^2$, so $c_0=F(\nu)a_0=0$, $c_1=F(\nu+1)a_1=0$, and for $n\ge2$ the coefficient is $c_n=F(\nu+n)a_n+a_{n-2}=n(n+2\nu)a_n+a_{n-2}$, which vanishes for odd $n$. For $n=2m$, by \eqref{eq:rgamma} with $s=m+\nu$,
\begin{align*}
2m(2m+2\nu)\,a_{2m}&=\frac{(-1)^m\,4m\,(m+\nu)}{2^{2m+\nu}\,m!\,\Gamma(m+\nu+1)}\\
&=\frac{(-1)^m}{2^{2m-2+\nu}\,(m-1)!\,\Gamma(m+\nu)}=-a_{2m-2},
\end{align*}
so $c_{2m}=0$ as well. Bessel's equation depends on $\nu$ only through $\nu^2$, so $J_\nu$ solves the equation of order $|\nu|$. For $\nu\ge0$ the coefficients are those of the recurrence \eqref{eq:frobenius-recurrence} with $r=r_1=\nu$, scaled by $a_0=1/(2^\nu\Gamma(\nu+1))$.
\end{proof}

For $\nu=0$ and $\nu=1$,
\begin{gather*}
J_0(x)=\sum_{m=0}^\infty\frac{(-1)^m}{(m!)^2}\Bigl(\frac x2\Bigr)^{2m}=1-\frac{x^2}4+\frac{x^4}{64}-\cdots,\\
J_1(x)=\sum_{m=0}^\infty\frac{(-1)^m}{m!\,(m+1)!}\Bigl(\frac x2\Bigr)^{2m+1}=\frac x2-\frac{x^3}{16}+\cdots .
\end{gather*}
For a natural number $n$, $J_n$ is an entire power series, odd or even as $n$ is, and solves the equation on all of $\mathbb R$; its behaviour near $0$ is $J_n(x)\sim(x/2)^n/n!$. The function $\sqrt x\,J_1(2\sqrt x)=\sum_m(-1)^mx^{m+1}/(m!\,(m+1)!)$ is the solution $y_1$ of $xy''+y=0$ found in \cref{sec:frobenius}; the change of variable $x\mapsto2\sqrt x$ is one of many that bring equations to Bessel's form~\cite[Chap.~IV]{watson1944}. \Cref{fig:bessel-functions} shows $J_0$, $J_1$ and $J_2$.

\begin{theorem}[A fundamental set for non-integral order]
\label{thm:bessel-basis}
If $\nu>0$ is not an integer, $J_\nu$ and $J_{-\nu}$ form a fundamental set of Bessel's equation of order $\nu$ on $(0,\infty)$, and
\[
W(J_\nu,J_{-\nu})(x)=-\frac2{x\,\Gamma(\nu)\,\Gamma(1-\nu)}\qquad(x>0).
\]
\end{theorem}

\begin{proof}
On $(0,\infty)$ the normal form of \eqref{eq:bessel} has the coefficient $a_1(x)=1/x$, so by Abel's formula (\cref{thm:abel}) $W(x)=W(1)\exp(-\int_1^xds/s)=W(1)/x$, and $xW(x)$ is a constant $K$. To find it, write $J_\nu=x^\nu A$ and $J_{-\nu}=x^{-\nu}B$ with the power series $A$ and $B$ of \cref{thm:bessel-j}, for which $A(0)=1/(2^\nu\Gamma(\nu+1))$ and $B(0)=2^\nu/\Gamma(1-\nu)$; both are nonzero, since $\Gamma$ is finite at $\nu+1$ and, $\nu$ not being an integer, at $1-\nu$. Then
\begin{align*}
W&=x^\nu A\bigl(-\nu x^{-\nu-1}B+x^{-\nu}B'\bigr)-\bigl(\nu x^{\nu-1}A+x^\nu A'\bigr)x^{-\nu}B\\
&=-\frac{2\nu AB}x+AB'-A'B ,
\end{align*}
so $K=xW(x)=-2\nu A(x)B(x)+x(AB'-A'B)(x)$ for every $x>0$. Letting $x\to0^+$, the power series being continuous at $0$, gives $K=-2\nu A(0)B(0)=-2\nu/(\Gamma(\nu+1)\Gamma(1-\nu))=-2/(\Gamma(\nu)\Gamma(1-\nu))$, using $\Gamma(\nu+1)=\nu\Gamma(\nu)$. Since $\nu>0$ and $A(0)B(0)\ne0$, $K\ne0$, and the two solutions form a fundamental set by \cref{thm:wronskian-solutions}.
\end{proof}

With the reflection formula $\Gamma(\nu)\Gamma(1-\nu)=\pi/\sin(\pi\nu)$, quoted in the note in \vcref{1}{sec:general-powers} and proved in Volume~III, the Wronskian is $-2\sin(\nu\pi)/(\pi x)$; the proof above does not need it. The theorem covers the orders $\nu=1/2,3/2,\ldots$, for which the indicial roots $\pm\nu$ differ by the integer $2\nu$: in these cases the constant $C$ of \cref{thm:frobenius}\,(iv) is zero, and no logarithm appears. The reason is visible in the recurrence: for $r=-\nu$ the factor that vanishes is $n(n-2\nu)$ at the odd index $n=2\nu$, where the coefficient is zero anyway. For integral order the situation is different.

\begin{proposition}[Integral order]
\label{prop:bessel-negative-integer}
For $n\in\mathbb N$, $J_{-n}=(-1)^nJ_n$. For $n\ge1$, the constant $C$ of \cref{thm:frobenius}\,(iv) for Bessel's equation of order $n$ is not zero, so every solution that is not a multiple of $J_n$ contains a logarithmic term.
\end{proposition}

\begin{proof}
In the series \eqref{eq:bessel-j} for $J_{-n}$, the terms with $m<n$ vanish, since $1/\Gamma(m-n+1)=0$ for $m-n+1\le0$. Substituting $m=k+n$ in the others,
\begin{align*}
J_{-n}(x)&=\sum_{k=0}^\infty\frac{(-1)^{k+n}}{(k+n)!\,\Gamma(k+1)}\Bigl(\frac x2\Bigr)^{2k+n}\\
&=(-1)^n\sum_{k=0}^\infty\frac{(-1)^k}{k!\,(k+n)!}\Bigl(\frac x2\Bigr)^{2k+n}=(-1)^nJ_n(x).
\end{align*}
For the second statement, follow the proof of \cref{thm:frobenius}\,(iv) with $r_2=-n$ and $N=2n$. The coefficients $b_{2k}$ for $k<n$ satisfy $2k(2k-2n)\,b_{2k}=-b_{2k-2}$ with $b_0=1$, so they are all nonzero, and the equation for the index $N=2n$ is $b_{2n-2}+CN=0$, since $Q_2=1$ is the only nonzero coefficient that enters. Hence $C=-b_{2n-2}/(2n)\ne0$.
\end{proof}

\subsection{Bessel Functions of the Second Kind}

For integral order a second solution must contain a logarithm. For $n=0$, \cref{thm:frobenius}\,(iii) applies, and the computation can be carried out completely.

\begin{proposition}[The logarithmic solution of order zero]
\label{prop:bessel-y0}
The function
\[
y_2(x)\coloneqq J_0(x)\ln x+\sum_{m=1}^\infty\frac{(-1)^{m+1}H_m}{(m!)^2}\Bigl(\frac x2\Bigr)^{2m},
\]
with the harmonic numbers $H_m=1+1/2+\cdots+1/m$, is the second solution of Bessel's equation of order $0$ given by \cref{thm:frobenius}\,(iii), and $J_0$, $y_2$ form a fundamental set on $(0,\infty)$.
\end{proposition}

\begin{proof}
Here $r_1=r_2=0$, $P=1$, $Q=x^2$, $F(s)=s^2$, and $J_0=A$ has the coefficients $a_{2m}=(-1)^m/(4^m(m!)^2)$ and $a_{2m+1}=0$. The coefficients of $G=2xA'+(P-1)A=2xA'$ are $g_n=2na_n$. The recurrence of case (iii) is $n^2b_n=-g_n-b_{n-2}$, since $Q_2=1$ is the only nonzero coefficient with positive index, with $b_0=0$ and $b_{-1}\coloneqq0$. It gives $b_1=0$, hence $b_n=0$ for all odd $n$, and $4m^2b_{2m}=-4m\,a_{2m}-b_{2m-2}$. We show $b_{2m}=(-1)^{m+1}H_m/(4^m(m!)^2)$ by induction; it holds for $m=0$ with $H_0=0$. If it holds for $m-1$, then $b_{2m-2}=(-1)^mH_{m-1}\cdot4m^2/(4^m(m!)^2)$, and
\[
4m^2\,b_{2m}=-\frac{(-1)^m\bigl(4m+4m^2H_{m-1}\bigr)}{4^m(m!)^2}=-\frac{(-1)^m\,4m^2\,H_m}{4^m(m!)^2},
\]
since $4m+4m^2H_{m-1}=4m^2(H_{m-1}+1/m)$. Dividing by $4m^2$ gives the claim, and $x^0B=\sum b_{2m}x^{2m}$ is the series in the proposition. \Cref{thm:frobenius}\,(iii), with $R=\infty$, shows that $y_2$ is a solution on $(0,\infty)$ and that $J_0$, $y_2$ form a fundamental set.
\end{proof}

The second solution is normalised in a conventional way that makes its behaviour for large $x$ resemble that of $J_\nu$, with a sine in place of a cosine (\cref{thm:bessel-asymptotic}).

\begin{definition}[Bessel functions of the second kind]
\label{def:bessel-y}
For $\nu>0$ not an integer, the \term{Bessel function of the second kind} is
\[
Y_\nu(x)\coloneqq\frac{J_\nu(x)\cos(\nu\pi)-J_{-\nu}(x)}{\sin(\nu\pi)}\qquad(x>0),
\]
and for integral order $n$ it is $Y_n(x)\coloneqq\lim_{\nu\to n}Y_\nu(x)$. In particular, with Euler's constant $\gamma=0.5772\ldots$ (\vcref{1}{thm:euler-constant}),
\begin{equation}
Y_0(x)=\frac2\pi\Bigl(\ln\frac x2+\gamma\Bigr)J_0(x)+\frac2\pi\sum_{m=1}^\infty\frac{(-1)^{m+1}H_m}{(m!)^2}\Bigl(\frac x2\Bigr)^{2m} .
\label{eq:bessel-y0}
\end{equation}
\end{definition}

For non-integral $\nu$, $Y_\nu$ is a solution of Bessel's equation, and $J_\nu$, $Y_\nu$ form a fundamental set, since the coefficient of $J_{-\nu}$ is not zero. The definition for integral order rests on the following theorem, which we do not prove in full; by \cref{prop:bessel-y0}, the function \eqref{eq:bessel-y0} is $(2/\pi)\bigl(y_2+(\gamma-\ln2)J_0\bigr)$, a solution independent of $J_0$, so that only the identification with the limit is taken on trust in the case $n=0$.

\begin{theorem}[Integral order]
\label{thm:bessel-y-limit}
For every $n\in\mathbb N$ and $x>0$, the limit $Y_n(x)=\lim_{\nu\to n}Y_\nu(x)$ exists, $Y_n$ solves Bessel's equation of order $n$, and $W(J_\nu,Y_\nu)(x)=2/(\pi x)$ for every $\nu\ge0$. For $n=0$ the limit is \eqref{eq:bessel-y0}, and for $n\ge1$
\begin{align*}
Y_n(x)&=\frac2\pi\ln\Bigl(\frac x2\Bigr)J_n(x)-\frac1\pi\sum_{k=0}^{n-1}\frac{(n-k-1)!}{k!}\Bigl(\frac x2\Bigr)^{2k-n}\\
&\qquad-\frac1\pi\sum_{k=0}^\infty\bigl(\psi(k+1)+\psi(n+k+1)\bigr)\frac{(-1)^k}{k!\,(n+k)!}\Bigl(\frac x2\Bigr)^{2k+n},
\end{align*}
where $\psi(m+1)\coloneqq H_m-\gamma$.
\end{theorem}

\begin{proof}[sketch]
For non-integral $\nu$, $W(J_\nu,Y_\nu)=-W(J_\nu,J_{-\nu})/\sin(\nu\pi)=2/(\pi x)$ by \cref{thm:bessel-basis} and the reflection formula. As $\nu\to n$, numerator and denominator of $Y_\nu$ tend to $0$, by \cref{prop:bessel-negative-integer}, and l'H\^opital's rule gives $Y_n=\bigl(\partial_\nu J_\nu-(-1)^n\partial_\nu J_{-\nu}\bigr)/\pi$ at $\nu=n$. The derivatives are computed by differentiating the series \eqref{eq:bessel-j} term by term with respect to $\nu$, which requires the derivative of $1/\Gamma$, the digamma function $\psi=\Gamma'/\Gamma$ with $\psi(m+1)=H_m-\gamma$, and its limits at the poles of $\Gamma$; this produces the formula. Justifying the differentiation term by term and the limits requires uniform convergence in $\nu$; Watson~\cite[Chap.~III]{watson1944} carries out the details, and the formula is \cite[Eq.~10.8.1]{olver2010nist}.
\end{proof}

Near $0$, $Y_0(x)\sim(2/\pi)\ln x$ and, for $n\ge1$, $Y_n(x)\sim-(n-1)!\,(2/x)^n/\pi$: every solution of integral order that is not a multiple of $J_n$ is unbounded at $0$. \Cref{fig:bessel-functions} shows $Y_0$. The values computed from \eqref{eq:bessel-y0} in ball arithmetic agree with those of python-flint's own Bessel function of the second kind to all twelve digits compared, for instance $Y_0(2)=0.510375672650$.

\begin{figure}
\sdfloat{%
\begin{tikzpicture}
\begin{axis}[sdplot, width=0.7\textwidth, height=0.4\textwidth, xmin=0, xmax=15, ymin=-1.05, ymax=1.25,
  xtick={0,2.5,5,7.5,10,12.5,15}, ytick={-1,-0.5,0,0.5,1}, xlabel={$x$}, sdlegendtop]
\addplot[black, line width=0.3pt, domain=0:15, samples=2, forget plot] {0};
\addplot[black, line width=0.8pt, smooth] coordinates {(0,1.0000) (0.1,0.9975) (0.2,0.9900) (0.3,0.9776) (0.4,0.9604) (0.5,0.9385) (0.6,0.9120) (0.7,0.8812) (0.8,0.8463) (0.9,0.8075) (1,0.7652) (1.1,0.7196) (1.2,0.6711) (1.3,0.6201) (1.4,0.5669) (1.5,0.5118) (1.6,0.4554) (1.7,0.3980) (1.8,0.3400) (1.9,0.2818) (2,0.2239) (2.1,0.1666) (2.2,0.1104) (2.3,0.0555) (2.4,0.0025) (2.5,-0.0484) (2.6,-0.0968) (2.7,-0.1424) (2.8,-0.1850) (2.9,-0.2243) (3,-0.2601) (3.1,-0.2921) (3.2,-0.3202) (3.3,-0.3443) (3.4,-0.3643) (3.5,-0.3801) (3.6,-0.3918) (3.7,-0.3992) (3.8,-0.4026) (3.9,-0.4018) (4,-0.3971) (4.1,-0.3887) (4.2,-0.3766) (4.3,-0.3610) (4.4,-0.3423) (4.5,-0.3205) (4.6,-0.2961) (4.7,-0.2693) (4.8,-0.2404) (4.9,-0.2097) (5,-0.1776) (5.1,-0.1443) (5.2,-0.1103) (5.3,-0.0758) (5.4,-0.0412) (5.5,-0.0068) (5.6,0.0270) (5.7,0.0599) (5.8,0.0917) (5.9,0.1220) (6,0.1506) (6.1,0.1773) (6.2,0.2017) (6.3,0.2238) (6.4,0.2433) (6.5,0.2601) (6.6,0.2740) (6.7,0.2851) (6.8,0.2931) (6.9,0.2981) (7,0.3001) (7.1,0.2991) (7.2,0.2951) (7.3,0.2882) (7.4,0.2786) (7.5,0.2663) (7.6,0.2516) (7.7,0.2346) (7.8,0.2154) (7.9,0.1944) (8,0.1717) (8.1,0.1475) (8.2,0.1222) (8.3,0.0960) (8.4,0.0692) (8.5,0.0419) (8.6,0.0146) (8.7,-0.0125) (8.8,-0.0392) (8.9,-0.0653) (9,-0.0903) (9.1,-0.1142) (9.2,-0.1367) (9.3,-0.1577) (9.4,-0.1768) (9.5,-0.1939) (9.6,-0.2090) (9.7,-0.2218) (9.8,-0.2323) (9.9,-0.2403) (10,-0.2459) (10.1,-0.2490) (10.2,-0.2496) (10.3,-0.2477) (10.4,-0.2434) (10.5,-0.2366) (10.6,-0.2276) (10.7,-0.2164) (10.8,-0.2032) (10.9,-0.1881) (11,-0.1712) (11.1,-0.1528) (11.2,-0.1330) (11.3,-0.1121) (11.4,-0.0902) (11.5,-0.0677) (11.6,-0.0446) (11.7,-0.0213) (11.8,0.0020) (11.9,0.0250) (12,0.0477) (12.1,0.0697) (12.2,0.0908) (12.3,0.1108) (12.4,0.1296) (12.5,0.1469) (12.6,0.1626) (12.7,0.1766) (12.8,0.1887) (12.9,0.1988) (13,0.2069) (13.1,0.2129) (13.2,0.2167) (13.3,0.2183) (13.4,0.2177) (13.5,0.2150) (13.6,0.2101) (13.7,0.2032) (13.8,0.1943) (13.9,0.1836) (14,0.1711) (14.1,0.1570) (14.2,0.1414) (14.3,0.1245) (14.4,0.1065) (14.5,0.0875) (14.6,0.0679) (14.7,0.0476) (14.8,0.0271) (14.9,0.0064) (15,-0.0142)};
\addlegendentry{$J_0$}
\addplot[black, dashed, line width=0.8pt, smooth] coordinates {(0,0.0000) (0.1,0.0499) (0.2,0.0995) (0.3,0.1483) (0.4,0.1960) (0.5,0.2423) (0.6,0.2867) (0.7,0.3290) (0.8,0.3688) (0.9,0.4059) (1,0.4401) (1.1,0.4709) (1.2,0.4983) (1.3,0.5220) (1.4,0.5419) (1.5,0.5579) (1.6,0.5699) (1.7,0.5778) (1.8,0.5815) (1.9,0.5812) (2,0.5767) (2.1,0.5683) (2.2,0.5560) (2.3,0.5399) (2.4,0.5202) (2.5,0.4971) (2.6,0.4708) (2.7,0.4416) (2.8,0.4097) (2.9,0.3754) (3,0.3391) (3.1,0.3009) (3.2,0.2613) (3.3,0.2207) (3.4,0.1792) (3.5,0.1374) (3.6,0.0955) (3.7,0.0538) (3.8,0.0128) (3.9,-0.0272) (4,-0.0660) (4.1,-0.1033) (4.2,-0.1386) (4.3,-0.1719) (4.4,-0.2028) (4.5,-0.2311) (4.6,-0.2566) (4.7,-0.2791) (4.8,-0.2985) (4.9,-0.3147) (5,-0.3276) (5.1,-0.3371) (5.2,-0.3432) (5.3,-0.3460) (5.4,-0.3453) (5.5,-0.3414) (5.6,-0.3343) (5.7,-0.3241) (5.8,-0.3110) (5.9,-0.2951) (6,-0.2767) (6.1,-0.2559) (6.2,-0.2329) (6.3,-0.2081) (6.4,-0.1816) (6.5,-0.1538) (6.6,-0.1250) (6.7,-0.0953) (6.8,-0.0652) (6.9,-0.0349) (7,-0.0047) (7.1,0.0252) (7.2,0.0543) (7.3,0.0826) (7.4,0.1096) (7.5,0.1352) (7.6,0.1592) (7.7,0.1813) (7.8,0.2014) (7.9,0.2192) (8,0.2346) (8.1,0.2476) (8.2,0.2580) (8.3,0.2657) (8.4,0.2708) (8.5,0.2731) (8.6,0.2728) (8.7,0.2697) (8.8,0.2641) (8.9,0.2559) (9,0.2453) (9.1,0.2324) (9.2,0.2174) (9.3,0.2004) (9.4,0.1816) (9.5,0.1613) (9.6,0.1395) (9.7,0.1166) (9.8,0.0928) (9.9,0.0684) (10,0.0435) (10.1,0.0184) (10.2,-0.0066) (10.3,-0.0313) (10.4,-0.0555) (10.5,-0.0789) (10.6,-0.1012) (10.7,-0.1224) (10.8,-0.1422) (10.9,-0.1603) (11,-0.1768) (11.1,-0.1913) (11.2,-0.2039) (11.3,-0.2143) (11.4,-0.2225) (11.5,-0.2284) (11.6,-0.2320) (11.7,-0.2333) (11.8,-0.2323) (11.9,-0.2290) (12,-0.2234) (12.1,-0.2157) (12.2,-0.2060) (12.3,-0.1943) (12.4,-0.1807) (12.5,-0.1655) (12.6,-0.1487) (12.7,-0.1307) (12.8,-0.1114) (12.9,-0.0912) (13,-0.0703) (13.1,-0.0489) (13.2,-0.0271) (13.3,-0.0052) (13.4,0.0166) (13.5,0.0380) (13.6,0.0590) (13.7,0.0791) (13.8,0.0984) (13.9,0.1165) (14,0.1334) (14.1,0.1488) (14.2,0.1626) (14.3,0.1747) (14.4,0.1850) (14.5,0.1934) (14.6,0.1999) (14.7,0.2043) (14.8,0.2066) (14.9,0.2069) (15,0.2051)};
\addlegendentry{$J_1$}
\addplot[black, densely dotted, line width=0.9pt, smooth] coordinates {(0,0.0000) (0.1,0.0012) (0.2,0.0050) (0.3,0.0112) (0.4,0.0197) (0.5,0.0306) (0.6,0.0437) (0.7,0.0588) (0.8,0.0758) (0.9,0.0946) (1,0.1149) (1.1,0.1366) (1.2,0.1593) (1.3,0.1830) (1.4,0.2074) (1.5,0.2321) (1.6,0.2570) (1.7,0.2817) (1.8,0.3061) (1.9,0.3299) (2,0.3528) (2.1,0.3746) (2.2,0.3951) (2.3,0.4139) (2.4,0.4310) (2.5,0.4461) (2.6,0.4590) (2.7,0.4696) (2.8,0.4777) (2.9,0.4832) (3,0.4861) (3.1,0.4862) (3.2,0.4835) (3.3,0.4780) (3.4,0.4697) (3.5,0.4586) (3.6,0.4448) (3.7,0.4283) (3.8,0.4093) (3.9,0.3879) (4,0.3641) (4.1,0.3383) (4.2,0.3105) (4.3,0.2811) (4.4,0.2501) (4.5,0.2178) (4.6,0.1846) (4.7,0.1506) (4.8,0.1161) (4.9,0.0813) (5,0.0466) (5.1,0.0121) (5.2,-0.0217) (5.3,-0.0547) (5.4,-0.0867) (5.5,-0.1173) (5.6,-0.1464) (5.7,-0.1737) (5.8,-0.1990) (5.9,-0.2221) (6,-0.2429) (6.1,-0.2612) (6.2,-0.2769) (6.3,-0.2899) (6.4,-0.3001) (6.5,-0.3074) (6.6,-0.3119) (6.7,-0.3135) (6.8,-0.3123) (6.9,-0.3082) (7,-0.3014) (7.1,-0.2920) (7.2,-0.2800) (7.3,-0.2656) (7.4,-0.2490) (7.5,-0.2303) (7.6,-0.2097) (7.7,-0.1875) (7.8,-0.1638) (7.9,-0.1389) (8,-0.1130) (8.1,-0.0864) (8.2,-0.0593) (8.3,-0.0320) (8.4,-0.0047) (8.5,0.0223) (8.6,0.0488) (8.7,0.0745) (8.8,0.0993) (8.9,0.1228) (9,0.1448) (9.1,0.1653) (9.2,0.1840) (9.3,0.2008) (9.4,0.2154) (9.5,0.2279) (9.6,0.2380) (9.7,0.2458) (9.8,0.2512) (9.9,0.2542) (10,0.2546) (10.1,0.2527) (10.2,0.2483) (10.3,0.2416) (10.4,0.2327) (10.5,0.2216) (10.6,0.2085) (10.7,0.1936) (10.8,0.1769) (10.9,0.1586) (11,0.1390) (11.1,0.1183) (11.2,0.0966) (11.3,0.0741) (11.4,0.0512) (11.5,0.0279) (11.6,0.0046) (11.7,-0.0185) (11.8,-0.0413) (11.9,-0.0635) (12,-0.0849) (12.1,-0.1053) (12.2,-0.1245) (12.3,-0.1424) (12.4,-0.1587) (12.5,-0.1734) (12.6,-0.1862) (12.7,-0.1972) (12.8,-0.2061) (12.9,-0.2130) (13,-0.2177) (13.1,-0.2203) (13.2,-0.2208) (13.3,-0.2191) (13.4,-0.2152) (13.5,-0.2094) (13.6,-0.2015) (13.7,-0.1917) (13.8,-0.1801) (13.9,-0.1668) (14,-0.1520) (14.1,-0.1358) (14.2,-0.1185) (14.3,-0.1001) (14.4,-0.0808) (14.5,-0.0609) (14.6,-0.0405) (14.7,-0.0199) (14.8,0.0008) (14.9,0.0214) (15,0.0416)};
\addlegendentry{$J_2$}
\addplot[black!45, line width=0.8pt, smooth] coordinates {(0.1,-1.5342) (0.2,-1.0811) (0.3,-0.8073) (0.4,-0.6060) (0.5,-0.4445) (0.6,-0.3085) (0.7,-0.1907) (0.8,-0.0868) (0.9,0.0056) (1,0.0883) (1.1,0.1622) (1.2,0.2281) (1.3,0.2865) (1.4,0.3379) (1.5,0.3824) (1.6,0.4204) (1.7,0.4520) (1.8,0.4774) (1.9,0.4968) (2,0.5104) (2.1,0.5183) (2.2,0.5208) (2.3,0.5181) (2.4,0.5104) (2.5,0.4981) (2.6,0.4813) (2.7,0.4605) (2.8,0.4359) (2.9,0.4079) (3,0.3769) (3.1,0.3431) (3.2,0.3071) (3.3,0.2691) (3.4,0.2296) (3.5,0.1890) (3.6,0.1477) (3.7,0.1061) (3.8,0.0645) (3.9,0.0234) (4,-0.0169) (4.1,-0.0561) (4.2,-0.0938) (4.3,-0.1296) (4.4,-0.1633) (4.5,-0.1947) (4.6,-0.2235) (4.7,-0.2494) (4.8,-0.2723) (4.9,-0.2921) (5,-0.3085) (5.1,-0.3216) (5.2,-0.3313) (5.3,-0.3374) (5.4,-0.3402) (5.5,-0.3395) (5.6,-0.3354) (5.7,-0.3282) (5.8,-0.3177) (5.9,-0.3044) (6,-0.2882) (6.1,-0.2694) (6.2,-0.2483) (6.3,-0.2251) (6.4,-0.1999) (6.5,-0.1732) (6.6,-0.1452) (6.7,-0.1162) (6.8,-0.0864) (6.9,-0.0563) (7,-0.0259) (7.1,0.0042) (7.2,0.0339) (7.3,0.0628) (7.4,0.0907) (7.5,0.1173) (7.6,0.1424) (7.7,0.1658) (7.8,0.1872) (7.9,0.2065) (8,0.2235) (8.1,0.2381) (8.2,0.2501) (8.3,0.2595) (8.4,0.2662) (8.5,0.2702) (8.6,0.2715) (8.7,0.2700) (8.8,0.2659) (8.9,0.2592) (9,0.2499) (9.1,0.2383) (9.2,0.2245) (9.3,0.2086) (9.4,0.1907) (9.5,0.1712) (9.6,0.1502) (9.7,0.1279) (9.8,0.1045) (9.9,0.0804) (10,0.0557) (10.1,0.0307) (10.2,0.0056) (10.3,-0.0193) (10.4,-0.0437) (10.5,-0.0675) (10.6,-0.0904) (10.7,-0.1122) (10.8,-0.1326) (10.9,-0.1516) (11,-0.1688) (11.1,-0.1843) (11.2,-0.1977) (11.3,-0.2091) (11.4,-0.2183) (11.5,-0.2252) (11.6,-0.2299) (11.7,-0.2322) (11.8,-0.2322) (11.9,-0.2298) (12,-0.2252) (12.1,-0.2184) (12.2,-0.2095) (12.3,-0.1986) (12.4,-0.1858) (12.5,-0.1712) (12.6,-0.1551) (12.7,-0.1375) (12.8,-0.1187) (12.9,-0.0989) (13,-0.0782) (13.1,-0.0569) (13.2,-0.0352) (13.3,-0.0134) (13.4,0.0085) (13.5,0.0301) (13.6,0.0512) (13.7,0.0717) (13.8,0.0913) (13.9,0.1099) (14,0.1272) (14.1,0.1431) (14.2,0.1575) (14.3,0.1703) (14.4,0.1812) (14.5,0.1903) (14.6,0.1974) (14.7,0.2025) (14.8,0.2056) (14.9,0.2065) (15,0.2055)};
\addlegendentry{$Y_0$}
\end{axis}
\end{tikzpicture}}{The Bessel functions $J_0$ (solid), $J_1$ (dashed), $J_2$ (dotted) and $Y_0$ (grey). Near $0$, $J_n(x)$ behaves like $(x/2)^n/n!$ and $Y_0(x)$ like $(2/\pi)\ln x$; for large $x$ all of them oscillate with an amplitude that decreases like $x^{-1/2}$ and with zeros about $\pi$ apart (\cref{thm:bessel-zeros,thm:bessel-asymptotic}).\label{fig:bessel-functions}}
\end{figure}

\subsection{Half-Integral Order}

For the orders $\pm1/2$ the series can be summed in closed form. This completes the example after \cref{prop:normal-form}, where the normal form $w''+w=0$ showed that the solutions of order $1/2$ are $(c_1\cos x+c_2\sin x)/\sqrt x$.

\begin{proposition}[Bessel functions of order one half]
\label{prop:bessel-half}
For $x>0$,
\[
J_{1/2}(x)=\sqrt{\frac2{\pi x}}\,\sin x,\qquad J_{-1/2}(x)=\sqrt{\frac2{\pi x}}\,\cos x .
\]
\end{proposition}

\begin{proof}
By \vcref{1}{cor:gaussian-consequences}, $\Gamma(k+1/2)=(2k)!\sqrt\pi/(4^kk!)$ for $k\in\mathbb N$. With $k=m+1$,
\[
m!\,\Gamma\Bigl(m+\frac32\Bigr)=\frac{m!\,(2m+2)!\sqrt\pi}{4^{m+1}(m+1)!}=\frac{(2m+2)!\sqrt\pi}{4^{m+1}(m+1)}=\frac{(2m+1)!\sqrt\pi}{2\cdot4^m},
\]
so the term of index $m$ in \eqref{eq:bessel-j} with $\nu=1/2$ is
\[
\frac{(-1)^m\,2\cdot4^m}{(2m+1)!\,\sqrt\pi}\Bigl(\frac x2\Bigr)^{2m+1/2}=\sqrt{\frac2{\pi x}}\,\frac{(-1)^m\,x^{2m+1}}{(2m+1)!},
\]
and the sum is $\sqrt{2/(\pi x)}\sin x$ by \vcref{1}{thm:sin-cos-series}. For $\nu=-1/2$, $m!\,\Gamma(m+1/2)=(2m)!\sqrt\pi/4^m$ by the same formula, and the terms are $\sqrt{2/(\pi x)}\,(-1)^mx^{2m}/(2m)!$.
\end{proof}

By the recurrence relations of the next subsection, every Bessel function of order $n+1/2$, $n\in\mathbb Z$, is an elementary function: for instance $J_{3/2}(x)=J_{1/2}(x)/x-J_{-1/2}(x)=\sqrt{2/(\pi x)}\,(\sin x/x-\cos x)$. Multiplied by $\sqrt{\pi/(2x)}$, they are the \term{spherical Bessel functions} of physics, which arise when the wave equation is separated in spherical coordinates. By \cref{thm:bessel-basis}, $W(J_{1/2},J_{-1/2})=-2/(x\Gamma(1/2)^2)=-2/(\pi x)$, as the closed forms confirm.

\subsection{Recurrence Relations}

\begin{proposition}[Recurrence relations]
\label{prop:bessel-recurrences}
For every real $\nu$ and $x>0$:
\begin{enumerate}
\item $(x^\nu J_\nu)'=x^\nu J_{\nu-1}$ and $(x^{-\nu}J_\nu)'=-x^{-\nu}J_{\nu+1}$;
\item $J_{\nu-1}(x)+J_{\nu+1}(x)=(2\nu/x)\,J_\nu(x)$ and $J_{\nu-1}(x)-J_{\nu+1}(x)=2J_\nu'(x)$.
\end{enumerate}
In particular $J_0'=-J_1$.
\end{proposition}

\begin{proof}
(i) By \cref{thm:bessel-j}, the function $x^{-\nu}J_\nu(x)$ is the sum $\sum_md_mx^{2m}$ with the coefficients $d_m=(-1)^m/(2^{2m+\nu}m!\,\Gamma(m+\nu+1))$, a power series of infinite radius, which may be differentiated term by term:
\begin{align*}
(x^{-\nu}J_\nu)'&=\sum_{m\ge1}\frac{(-1)^m\,2m\,x^{2m-1}}{2^{2m+\nu}\,m!\,\Gamma(m+\nu+1)}\\
&=\sum_{k\ge0}\frac{(-1)^{k+1}x^{2k+1}}{2^{2k+\nu+1}\,k!\,\Gamma(k+\nu+2)}=-x^{-\nu}J_{\nu+1}(x),
\end{align*}
with $m=k+1$. Similarly $x^\nu J_\nu=x^{2\nu}\sum_md_mx^{2m}$, so $(x^\nu J_\nu)'=\sum_md_m(2m+2\nu)x^{2m+2\nu-1}$, and by \eqref{eq:rgamma} with $s=m+\nu$, $d_m(2m+2\nu)=(-1)^m/(2^{2m+\nu-1}m!\,\Gamma(m+\nu))$, which is the coefficient of $x^{2m+2(\nu-1)+1}$ in $x^{\nu}J_{\nu-1}=x^{2\nu-1}\cdot x^{-(\nu-1)}J_{\nu-1}$. (ii) By the product rule, the two identities of (i) say $J_\nu'+(\nu/x)J_\nu=J_{\nu-1}$ and $J_\nu'-(\nu/x)J_\nu=-J_{\nu+1}$. Subtracting and adding them gives (ii). For $\nu=0$, the second identity of (i) is $J_0'=-J_1$.
\end{proof}

The first identity of (ii) is a three-term recurrence in the order: from $J_0(x)$ and $J_1(x)$ it produces $J_2(x)$, $J_3(x)$, and so on, with a few operations each. Whether this is a good way to compute $J_n(x)$ in floating-point arithmetic is a question about the stability of recurrences, which \cref{ch:difference-equations} takes up; Volume~XI uses the answer to compute Bessel functions of integral order. The functions $Y_\nu$ satisfy the same relations: for non-integral $\nu$ this follows from the definition, since $\cos(\nu\pi)$ and $\sin(\nu\pi)$ change sign when $\nu$ changes by $1$, and for integral order by passing to the limit~\cite[Chap.~III]{watson1944}.

\begin{note}{Generating functions}
The Bessel functions of integral order are the coefficients of a Laurent series in $t$:
\[
\exp\Bigl(\frac x2\Bigl(t-\frac1t\Bigr)\Bigr)=\sum_{n=-\infty}^\infty J_n(x)\,t^n\qquad(t\ne0).
\]
To see this, multiply the two absolutely convergent series $e^{xt/2}=\sum_j(x/2)^jt^j/j!$ and $e^{-x/(2t)}=\sum_k(-1)^k(x/2)^kt^{-k}/k!$ and collect the products with $j-k=n$, which is permitted for absolutely convergent double series~\cite[Chap.~II]{knopp1990}; the coefficient of $t^n$ for $n\ge0$ is $\sum_k(-1)^k(x/2)^{n+2k}/(k!\,(n+k)!)=J_n(x)$, and for negative $n$ the same computation gives the coefficient $(-1)^nJ_{|n|}(x)$, which is $J_n(x)$ by \cref{prop:bessel-negative-integer}. With $t=e^{i\theta}$ the left side is $e^{ix\sin\theta}$, and the coefficients of this Fourier series are the integrals
\[
J_n(x)=\frac1{2\pi}\int_{-\pi}^\pi\cos(n\theta-x\sin\theta)\,d\theta ,
\]
Bessel's own definition~\cite{bessel1824}, \cite[Chap.~II]{watson1944}. The Legendre polynomials have the generating function $(1-2xt+t^2)^{-1/2}=\sum_nP_n(x)t^n$ for $|x|\le1$ and $|t|<1$, the expansion of the reciprocal distance $1/\|\mathbf a-\mathbf b\|$ between two points in powers of the ratio of their distances from the origin~\cite[Chap.~XV]{whittaker1927}. Generating functions of this kind were the subject of \vcref{1}{sec:gf}.
\end{note}

\subsection{Zeros}

The zeros of Bessel functions determine the frequencies of drums and the modes of waveguides, and they are tabulated in every handbook. Their basic properties follow from Sturm's theorems.

\begin{theorem}[Zeros of solutions of Bessel's equation]
\label{thm:bessel-zeros}
Let $\nu\ge0$, and let $y$ be a nontrivial real solution of \eqref{eq:bessel} on $(0,\infty)$.
\begin{enumerate}
\item $y$ has infinitely many positive zeros, only finitely many in each interval $[a,b]\subseteq(0,\infty)$, and they form an increasing sequence tending to $\infty$.
\item The distance between consecutive zeros tends to $\pi$.
\item For $0\le\nu<1/2$ consecutive zeros are less than $\pi$ apart, for $\nu>1/2$ more than $\pi$ apart, and for $\nu=1/2$ exactly $\pi$ apart.
\end{enumerate}
\end{theorem}

\begin{proof}
On $(0,\infty)$ the normal form has $a_1(x)=1/x$, which is continuously differentiable, and by \cref{prop:normal-form} (with $x_0=1$) the function $w=\sqrt x\,y$ solves
\[
w''+q(x)\,w=0,\qquad q(x)=1+\frac{1/4-\nu^2}{x^2},
\]
and has the same zeros as $y$. By \cref{lem:simple-zeros}, the zeros of $w$ in each closed bounded subinterval are finitely many. They do not accumulate at $0$ either. The indicial roots $\pm\nu$ are real, so by \cref{thm:frobenius} $y=c_1y_1+c_2y_2$ with real constants and the real solutions $y_1=x^\nu A(x)$ and $y_2$ of the theorem. If $c_2=0$, then $y$ has no zeros near $0$, since $A(0)=1$. If $c_2\ne0$ and $\nu>0$, then $x^\nu y(x)\to c_2B(0)=c_2$ as $x\to0^+$, in case (ii) as in case (iv), because $x^{2\nu}$ and $x^{2\nu}\ln x$ tend to $0$; if $c_2\ne0$ and $\nu=0$, then $y(x)/\ln x\to c_2A(0)=c_2$. In every case $y$ has no zeros on some interval $(0,\varepsilon)$. Let $0<\delta<1$. Since $q(x)\to1$, there is $X>0$ with $(1-\delta)^2<q(x)<(1+\delta)^2$ for $x\ge X$. By \cref{cor:zero-spacing}\,(i) with $\omega=1-\delta$, applied on $(X,\infty)$, $w$ has a zero in every interval $(a,a+\pi/(1-\delta))$ with $a>X$; so there are infinitely many zeros, and consecutive zeros $a<b$ beyond $X$ satisfy $b-a<\pi/(1-\delta)$. By \cref{cor:zero-spacing}\,(ii) with $\omega=1+\delta$, they also satisfy $b-a>\pi/(1+\delta)$. Since $\delta$ was arbitrary, this proves (i) and (ii). For (iii), if $\nu<1/2$, then $q>1$ on $(0,\infty)$, and \cref{cor:zero-spacing}\,(i) with $\omega=1$ shows that after each zero $a$ the next one comes before $a+\pi$; if $\nu>1/2$, then $q<1$, and part (ii) of the corollary gives distances larger than $\pi$. If $\nu=1/2$, $w''+w=0$, and $w=c\sin(x-\alpha)$ has its zeros exactly $\pi$ apart.
\end{proof}

The theorem applies to $J_\nu$, $Y_\nu$ and every other real solution. \Cref{tab:bessel-zeros} lists the first ten positive zeros $j_{0,k}$ of $J_0$, computed with python-flint. The distances between consecutive zeros are less than $\pi$, as (iii) says for $\nu=0$, and they increase towards $\pi$. The zeros are close to $(k-1/4)\pi$, which is where the cosine in the asymptotic form of \cref{thm:bessel-asymptotic} vanishes, and the differences decrease roughly like $1/(8(k-1/4)\pi)$, the next term of McMahon's expansion of the zeros~\cite[Sect.~10.21(vi)]{olver2010nist}.

\begin{table}
\sdfloat{%
\begin{sdtable}{colspec={rllll}, row{1}={font=\small\bfseries}}
\toprule
$k$ & $j_{0,k}$ & $(k-1/4)\pi$ & difference & $j_{0,k+1}-j_{0,k}$ \\
\midrule
$1$ & $2.404825558$ & $2.356194490$ & $4.86\times10^{-2}$ & $3.115253$ \\
$2$ & $5.520078110$ & $5.497787144$ & $2.23\times10^{-2}$ & $3.133650$ \\
$3$ & $8.653727913$ & $8.639379797$ & $1.43\times10^{-2}$ & $3.137807$ \\
$4$ & $11.79153444$ & $11.78097245$ & $1.06\times10^{-2}$ & $3.139383$ \\
$5$ & $14.93091771$ & $14.92256510$ & $8.35\times10^{-3}$ & $3.140146$ \\
$6$ & $18.07106397$ & $18.06415776$ & $6.91\times10^{-3}$ & $3.140573$ \\
$7$ & $21.21163663$ & $21.20575041$ & $5.89\times10^{-3}$ & $3.140835$ \\
$8$ & $24.35247153$ & $24.34734307$ & $5.13\times10^{-3}$ & $3.141008$ \\
$9$ & $27.49347913$ & $27.48893572$ & $4.54\times10^{-3}$ & $3.141127$ \\
$10$ & $30.63460647$ & $30.63052837$ & $4.08\times10^{-3}$ & \\
\bottomrule
\end{sdtable}}{The first ten positive zeros $j_{0,k}$ of $J_0$, computed with the ball arithmetic of python-flint, their distance from $(k-1/4)\pi$, and the distances between consecutive zeros, which are less than $\pi=3.141593$ and increase towards it (\cref{thm:bessel-zeros}).\label{tab:bessel-zeros}}
\end{table}

The zeros give the frequencies of a circular drum. For the radially symmetric vibrations, $m=0$ in the separation at the beginning of the section, the function $R$ must be bounded at the centre $r=0$, which excludes $Y_0$ and leaves $R(r)=J_0(\omega r/c)$; the membrane is fixed at the rim, so $J_0(\omega a/c)=0$, and the possible angular frequencies are $\omega_k=c\,j_{0,k}/a$. The overtones of the drum are therefore not integral multiples of the fundamental, as those of a string are: $\omega_2/\omega_1=j_{0,2}/j_{0,1}=2.2954\ldots$ and $\omega_3/\omega_1=3.5985\ldots$. Rayleigh's treatise on sound discusses the vibrations of membranes in these terms~\cite{rayleigh1877}.

\subsection{Large Arguments}

For large $x$ the normal form $w''+(1+(1/4-\nu^2)/x^2)w=0$ is close to $w''+w=0$, which suggests that $\sqrt x\,J_\nu(x)$ is close to a sinusoid. This is true, with constants that the power series do not reveal.

\begin{theorem}[Behaviour for large arguments]
\label{thm:bessel-asymptotic}
For every real $\nu$, as $x\to\infty$,
\begin{gather*}
J_\nu(x)=\sqrt{\frac2{\pi x}}\Bigl(\cos\Bigl(x-\frac{\nu\pi}2-\frac\pi4\Bigr)+O\Bigl(\frac1x\Bigr)\Bigr),\\
Y_\nu(x)=\sqrt{\frac2{\pi x}}\Bigl(\sin\Bigl(x-\frac{\nu\pi}2-\frac\pi4\Bigr)+O\Bigl(\frac1x\Bigr)\Bigr).
\end{gather*}
\end{theorem}

\begin{proof}[sketch]
The function $w=\sqrt x\,J_\nu$ solves $w''+w=\varepsilon(x)w$ with $\varepsilon(x)=(\nu^2-1/4)/x^2$. Treating the right side as a forcing term, variation of parameters (\cref{thm:variation-two}) gives, for $x\ge1$,
\[
w(x)=A\cos x+B\sin x-\int_x^\infty\sin(x-s)\,\varepsilon(s)\,w(s)\,ds
\]
with constants $A$ and $B$, the integral converging because $\varepsilon$ is integrable on $[1,\infty)$ and $w$ turns out to be bounded, by Gr\"onwall's inequality (\cref{lem:gronwall}). The integral is $O(\int_x^\infty s^{-2}ds)=O(1/x)$, so $w(x)=a\cos(x-\varphi)+O(1/x)$ for some amplitude $a$ and phase $\varphi$. The values $a=\sqrt{2/\pi}$ and $\varphi=\nu\pi/2+\pi/4$ cannot be found from the equation alone, since every solution has this form; they follow from an integral representation of $J_\nu$, such as Bessel's integral of the note on generating functions for integral $\nu$, by the method of stationary phase. For non-integral $\nu$ the statement for $Y_\nu$ follows from those for $J_{\pm\nu}$ by the definition and the addition formula of the cosine, and for integral order by passing to the limit. Watson~\cite[Chap.~VII]{watson1944} and Olver~\cite[Chap.~7]{olver1997asymptotics} give complete proofs and the full asymptotic expansions, and Volume~XI develops them.
\end{proof}

The amplitude $\sqrt{2/(\pi x)}$ is visible in \cref{fig:bessel-functions}, and the zeros of the cosine, at $x=(k-1/4)\pi$ for $\nu=0$, are the approximations of \cref{tab:bessel-zeros}. The theorem also shows why the series \eqref{eq:bessel-j} is a poor way to compute $J_0(x)$ for large $x$: the sum is of size $x^{-1/2}$, while the terms, as \cref{sec:series-computation} measures, are enormous.

\subsection{The Riccati Equation Completed}

The solutions $w_1$ and $w_2$ of $w''+t^2w=0$ in \eqref{eq:riccati-series} are Bessel functions of the orders $\mp1/4$, as the note in \cref{sec:direction-fields} announced.

\begin{proposition}[Bessel functions of order one quarter]
\label{prop:riccati-bessel}
For $t>0$,
\begin{gather*}
w_1(t)=\frac{\Gamma(3/4)}{\sqrt2}\,\sqrt t\,J_{-1/4}\Bigl(\frac{t^2}2\Bigr),\\
w_2(t)=\sqrt2\,\Gamma(5/4)\,\sqrt t\,J_{1/4}\Bigl(\frac{t^2}2\Bigr).
\end{gather*}
Consequently the blow-up time of the solution of $y'=t^2+y^2$ with $y(0)=0$ is $t^*=\sqrt{2j}$, where $j$ is the smallest positive zero of $J_{-1/4}$.
\end{proposition}

\begin{proof}
By \eqref{eq:bessel-j} with $\nu=-1/4$ and $x=t^2/2$, the term of index $m$ of $\sqrt t\,J_{-1/4}(t^2/2)$ is
\[
\sqrt t\,\frac{(-1)^m}{m!\,\Gamma(m+3/4)}\Bigl(\frac{t^2}4\Bigr)^{2m-1/4}=\frac{(-1)^m\,\sqrt2\,t^{4m}}{16^m\,m!\,\Gamma(m+3/4)} .
\]
By the functional equation of the gamma function, $\Gamma(m+3/4)$ equals $\Gamma(3/4)\prod_{j=1}^m(j-1/4)=\Gamma(3/4)\,4^{-m}\prod_{j=1}^m(4j-1)$, and $4^mm!=\prod_{j=1}^m4j$, so the term multiplied by $\Gamma(3/4)/\sqrt2$ is $(-1)^mt^{4m}/\prod_{j=1}^m(4j-1)\,4j$, the term of index $m$ of $w_1$ in \eqref{eq:riccati-series}. The computation for $w_2$, with $\nu=1/4$, $\Gamma(m+5/4)=\Gamma(5/4)\,4^{-m}\prod_{j=1}^m(4j+1)$ and the factor $(t^2/4)^{1/4}\sqrt t=t/\sqrt2$, is the same. By \cref{prop:riccati-blowup} the blow-up time is the smallest positive zero of $w_1$, which is the smallest $t>0$ with $J_{-1/4}(t^2/2)=0$.
\end{proof}

With python-flint, the smallest positive zero of $J_{-1/4}$ is $j=2.006299671789450\ldots$, and $\sqrt{2j}=2.003147359426884\ldots$, the value found in \cref{sec:airy} from the power series of $w_1$. More generally, the substitution $w=\sqrt t\,Z\bigl(2t^{(m+2)/2}/(m+2)\bigr)$ transforms $w''+t^mw=0$, for $m\ne-2$, into Bessel's equation of order $1/(m+2)$ for $Z$, as one checks by the chain rule; for $m=2$ this is the proposition. The order $1/(m+2)$ is half an odd integer, positive or negative, exactly when $m=-4k/(2k\pm1)$ with $k\in\mathbb N$, for which $m+2$ is $2/(2k+1)$ or $-2/(2k-1)$. These values are Liouville's exceptional exponents for $y'=y^2+t^m$ quoted in \cref{sec:direction-fields}, apart from $m=-2$, for which $w''+t^{-2}w=0$ is an Euler--Cauchy equation. So for $m\ne-2$ the Riccati equation is solvable by quadratures exactly when the Bessel functions that solve the corresponding linear equation are elementary. Watson proves both halves of this correspondence: that Bessel's equation has solutions expressible in finite terms only when twice the order is an odd integer, and that the Riccati equation is solvable in finite terms only in the exceptional cases~\cite[Chap.~IV]{watson1944}.

\section{Other Classical Equations: A Brief Look}
\label{sec:classical-equations}

\subsection{Polynomial Solutions}

The Legendre polynomials arose because a factor in the recurrence \eqref{eq:legendre-recurrence} vanishes for one index and cuts the series off. Three other classical equations have the same feature.

\term{Hermite's equation} $y''-2xy'+2\lambda y=0$ has no singular point, and substituting $y=\sum a_kx^k$ gives
\[
a_{k+2}=\frac{2(k-\lambda)}{(k+1)(k+2)}\,a_k .
\]
For $\lambda=n\in\mathbb N$ the series with the parity of $n$ stops at the degree $n$. Scaled to the leading coefficient $2^n$, these are the \term{Hermite polynomials} $H_0=1$, $H_1=2x$, $H_2=4x^2-2$, $H_3=8x^3-12x$; for $n=2$, the recurrence gives $a_2=-2a_0$, so the even solution is $1-2x^2$, and multiplying by $-2$ gives $H_2$. They describe the stationary states of the quantum harmonic oscillator, and their weight $e^{-x^2}$ is, up to scaling, the normal density.

\term{Chebyshev's equation} $(1-x^2)y''-xy'+n^2y=0$ has the singular points $\pm1$ and the recurrence $a_{k+2}=(k^2-n^2)a_k/((k+1)(k+2))$. Its polynomial solutions are the \term{Chebyshev polynomials} $T_n$, characterised by $T_n(\cos\theta)=\cos n\theta$: differentiating $y(\cos\theta)=\cos n\theta$ twice with respect to $\theta$ gives $\sin^2\theta\,y''(\cos\theta)-\cos\theta\,y'(\cos\theta)=-n^2\cos n\theta$, which is the equation at $x=\cos\theta$. The addition formula $\cos(n+1)\theta+\cos(n-1)\theta=2\cos\theta\cos n\theta$ gives $T_{n+1}=2xT_n-T_{n-1}$, and $T_0=1$, $T_1=x$, $T_2=2x^2-1$, $T_3=4x^3-3x$; that $T_n$ is a polynomial of degree $n$ follows from this recurrence by induction.

\term{Laguerre's equation} $xy''+(1-x)y'+ny=0$ has a regular singular point at $0$: multiplied by $x$, it is \eqref{eq:frobenius-form} with $P(x)=1-x$ and $Q(x)=nx$, so $F(r)=r^2$ and $0$ is a double indicial root. The recurrence \eqref{eq:frobenius-recurrence} with $r=0$ is $k^2a_k=(k-1-n)\,a_{k-1}$, which stops at $k=n+1$, and with $a_0=1$ it gives the \term{Laguerre polynomial}
\[
L_n(x)=\sum_{k=0}^n\frac{(-1)^k}{k!}\binom nk\,x^k ,
\]
for instance $L_1=1-x$ and $L_2=1-2x+x^2/2$.

\begin{table}
\sdfloat{%
\begin{sdtable}{colspec={llll}, row{1}={font=\small\bfseries}}
\toprule
polynomials & equation in self-adjoint form & interval & weight $w$ \\
\midrule
Legendre $P_n$ & $((1-x^2)\,y')'+n(n+1)\,y=0$ & $[-1,1]$ & $1$ \\
Chebyshev $T_n$ & $(\sqrt{1-x^2}\,y')'+n^2w\,y=0$ & $[-1,1]$ & $(1-x^2)^{-1/2}$ \\
Hermite $H_n$ & $(e^{-x^2}y')'+2n\,w\,y=0$ & $\mathbb R$ & $e^{-x^2}$ \\
Laguerre $L_n$ & $(xe^{-x}y')'+n\,w\,y=0$ & $[0,\infty)$ & $e^{-x}$ \\
\bottomrule
\end{sdtable}}{The classical orthogonal polynomials as polynomial solutions of differential equations (\cref{prop:classical-polynomials}), with the equations written as $(\sigma y')'+\lambda wy=0$. Polynomials of different degrees in each family are orthogonal in the inner product $\int wpq$ over the interval. The normalisations are $P_n(1)=T_n(1)=L_n(0)=1$ and the leading coefficient $2^n$ for $H_n$.\label{tab:orthogonal-polynomials}}
\end{table}

\begin{proposition}[Classical orthogonal polynomials]
\label{prop:classical-polynomials}
For each of the four equations of \cref{tab:orthogonal-polynomials} and each $n\in\mathbb N$, the polynomial solutions form a space of dimension $1$, spanned by a polynomial of degree $n$. Polynomials of different degrees in the same family are orthogonal with respect to $\langle p,q\rangle_w=\int pqw$ over the interval of the table.
\end{proposition}

\begin{proof}
For Legendre's equation this is \cref{thm:legendre-polynomials,thm:legendre-orthogonality}. For Hermite's and Chebyshev's equations the argument of \cref{thm:legendre-polynomials} applies: the factors $k-n$ and $k^2-n^2$ in the recurrences vanish only for $k=n$, so the series with the parity of $n$ is a polynomial of degree $n$, and the other series has infinitely many nonzero coefficients and is not a polynomial. For Laguerre's equation, the polynomial $L_n$ is the Frobenius solution of \cref{thm:frobenius}\,(i), and every solution independent of it contains the term $L_n\ln x$, by part (iii), and is not a polynomial. For the orthogonality, each equation can be written as $(\sigma y')'=-\lambda wy$ with the functions $\sigma$ and $w$ and the eigenvalues $\lambda$ of the table; for instance $(xe^{-x}y')'=e^{-x}(xy''+(1-x)y')$. As in the proof of \cref{thm:legendre-orthogonality}, two polynomial solutions $p_m$ and $p_n$ satisfy $(\lambda_m-\lambda_n)\,wp_mp_n=\bigl(\sigma(p_mp_n'-p_np_m')\bigr)'$. Integrate over a closed interval inside the interval of the table and let its ends tend to those of the table. The boundary terms tend to $0$, because $\sigma$ vanishes at $\pm1$ for Chebyshev and at $0$ for Laguerre, and $e^{-x^2}$ and $xe^{-x}$ times a polynomial tend to $0$ at infinity (\vcref{1}{prop:exp-growth}); the integrals $\int wp_mp_n$ converge, by \vcref{1}{prop:arctan-integral} for $(1-x^2)^{-1/2}$, which has the antiderivative $\arcsin x$, and by \vcref{1}{cor:gaussian-consequences,thm:gamma} for the moments of $e^{-x^2}$ and $e^{-x}$. Since $\lambda_m\ne\lambda_n$ for $m\ne n$, the integral vanishes.
\end{proof}

All four families satisfy three-term recurrences, by the argument of the note on three-term recurrences in \cref{subsec:legendre}, and their zeros are the nodes of the Gauss quadrature rules for the corresponding weights. The Chebyshev polynomials moreover play a central part in polynomial approximation on an interval. Szeg\H{o}~\cite{szego1975} is the classical reference; Volume~XI uses Chebyshev series to approximate special functions, and Volume~XII builds Gauss quadrature on the Legendre polynomials and treats approximation by Chebyshev polynomials.

\subsection{The Hypergeometric Equation}

The equations of this chapter have one or two singular points in the finite plane. The \term{hypergeometric equation}
\begin{equation}
x(1-x)\,y''+\bigl(c-(a+b+1)x\bigr)\,y'-ab\,y=0 ,
\label{eq:hypergeometric}
\end{equation}
with real constants $a$, $b$ and $c$, has in general three regular singular points, and it contains many of the others as special cases. We write $(a)_n\coloneqq a(a+1)\cdots(a+n-1)$, with $(a)_0=1$, for the rising factorial.

\begin{theorem}[The hypergeometric equation]
\label{thm:hypergeometric}
\leavevmode
\begin{enumerate}
\item Every point other than $0$ and $1$ is an ordinary point of \eqref{eq:hypergeometric}. The points $0$ and $1$, as well as the point at infinity, are regular singular points, except for special values of $a$, $b$ and $c$, for which some of them are ordinary; the indicial roots at $0$ are $0$ and $1-c$.
\item If $c\notin\{0,-1,-2,\ldots\}$, the \term{hypergeometric series}
\[
F(a,b;c;x)\coloneqq\sum_{n=0}^\infty\frac{(a)_n(b)_n}{(c)_n\,n!}\,x^n
\]
has a radius of convergence at least $1$, and its sum solves \eqref{eq:hypergeometric} on $(-1,1)$. It is a polynomial if $a$ or $b$ is a nonpositive integer, and its radius is exactly $1$ otherwise.
\end{enumerate}
\end{theorem}

\begin{proof}
(i) The leading coefficient $x(1-x)$ vanishes only at $0$ and $1$, so every other point is ordinary (\cref{prop:polynomial-ordinary}). The normal form has $p(x)=(c-(a+b+1)x)/(x(1-x))$ and $q(x)=-ab/(x(1-x))$, so $xp(x)=(c-(a+b+1)x)/(1-x)$ and $x^2q(x)=-abx/(1-x)$ are analytic at $0$, with the values $c$ and $0$ there; so $0$ is a regular singular point, unless $p$ and $q$ themselves are analytic there, and the indicial polynomial is $r(r-1)+cr=r(r-1+c)$. At $x=1$ the computation is the same with $1-x$ in place of $x$. The point at infinity is checked with the substitution $x=1/s$ of \cref{sec:frobenius}: since $x(1-x)=(s-1)/s^2$, the new coefficients are $2/s-(cs-(a+b+1))/(s(s-1))$ and $-ab/(s^2(s-1))$, and $s$ times the first and $s^2$ times the second are analytic at $s=0$. (ii) Substituting $y=\sum a_nx^n$ and collecting the coefficient of $x^n$ gives
\begin{align*}
&(n+1)n\,a_{n+1}-n(n-1)\,a_n+c(n+1)\,a_{n+1}-(a+b+1)n\,a_n-ab\,a_n\\
&\qquad=(n+1)(n+c)\,a_{n+1}-(n+a)(n+b)\,a_n ,
\end{align*}
so a power series solves the equation exactly when $a_{n+1}=(n+a)(n+b)\,a_n/((n+1)(n+c))$, and with $a_0=1$ this gives the coefficients of $F$; the condition on $c$ keeps the denominators nonzero. If a factor $n+a$ or $n+b$ vanishes, the series is a polynomial; otherwise the ratio of consecutive coefficients tends to $1$, and the radius is $1$ by \vcref{1}{thm:cauchy-hadamard}. In both cases the computation is justified on $(-1,1)$ by \vcref{1}{thm:termwise-diff,prop:formal-analytic}.
\end{proof}

Many functions are hypergeometric series: the geometric series is $F(1,b;b;x)$, $(1-x)^{-a}=F(a,b;b;x)$ by the binomial series, and $\ln(1+x)=xF(1,1;2;-x)$, as the recurrence $a_{n+1}/a_n=(n+1)^2/((n+1)(n+2))$ confirms. The Legendre polynomials are hypergeometric as well: the substitution $x=1-2s$ turns Legendre's equation with $\nu=n$ into
\[
s(1-s)\,Y''+(1-2s)\,Y'+n(n+1)\,Y=0 ,
\]
which is \eqref{eq:hypergeometric} with $a=-n$, $b=n+1$ and $c=1$. The polynomial $F(-n,n+1;1;s)$ has degree $n$ and the value $1$ at $s=0$, that is, at $x=1$, so by \cref{thm:legendre-polynomials}
\[
P_n(x)=F\Bigl(-n,\,n+1;\,1;\,\frac{1-x}2\Bigr) .
\]
In the same way $T_n(x)=F(-n,n;1/2;(1-x)/2)$. For non-integral $\nu$, $F(-\nu,\nu+1;1;(1-x)/2)$ is the solution of Legendre's equation that is analytic at $x=1$, used in the sketch of \cref{prop:legendre-nonintegral}. The functions $F$ and the confluent hypergeometric functions, which arise when two singular points merge and describe the Bessel, Hermite and Laguerre functions, are computed in Volume~XI; the stardust library provides them as \code{special.hyp2f1} and \code{special.hyp1f1}.

\begin{note}{Fuchsian equations and Riemann's P-symbol}
An equation all of whose singular points, including infinity, are regular is called \emph{Fuchsian}, after Fuchs, who showed in 1866 that the regular singular points are exactly those near which all solutions grow at most like a power of the distance~\cite{fuchs1866}. Riemann observed in 1857 that a Fuchsian equation of the second order with exactly three singular points is determined by their positions and the six indicial roots, and that a M\"obius transformation of $x$ moving the singular points to $0$, $1$ and $\infty$, together with a factor $x^\alpha(1-x)^\beta$, turns it into the hypergeometric equation; he wrote the solutions with a symbol listing the singular points and their exponents. Gauss had studied the hypergeometric series in 1812. Gray~\cite{gray2000} tells this history, and Whittaker and Watson~\cite[Chaps.~X and~XIV]{whittaker1927} give the theory; it belongs to the complex domain, where the solutions continued around a singular point change by a linear map, the \emph{monodromy}, and Volume~III provides the tools.
\end{note}

\subsection{Sturm--Liouville Problems Again}

The Legendre, Chebyshev, Hermite and Laguerre polynomials, and the functions $J_0(j_{0,k}r/a)$ of the vibrating drum, are eigenfunctions of operators of the form $y\mapsto-(\sigma y')'/w$ under conditions at the ends of an interval: boundedness for Legendre's equation at $\pm1$, growth slower than $e^{x^2/2}$ for Hermite's, the boundary condition $J_0(j_{0,k})=0$ for the drum. These are the Sturm--Liouville problems of the note in \cref{sec:variable-coefficients}, and in each case the eigenfunctions are orthogonal in the inner product $\int fgw$ by the argument of \cref{thm:legendre-orthogonality}; for the drum the weight is $r$, and $\int_0^arJ_0(j_{0,k}r/a)J_0(j_{0,l}r/a)\,dr=0$ for $k\ne l$. The expansion of an arbitrary function in such eigenfunctions, which generalises the Fourier series and the expansion in Legendre polynomials, is part of the spectral theory of Volume~III.

\section{Computing with Series: A First Look}
\label{sec:series-computation}

\subsection{Summing by the Recurrence}

A power series solution comes with a recurrence for its coefficients, and for the series of this chapter the recurrence relates consecutive terms by a simple factor. For $J_\nu$ the terms $t_m$ of \eqref{eq:bessel-j} satisfy
\[
t_{m+1}=\rho_m\,t_m,\qquad\rho_m=-\frac{(x/2)^2}{(m+1)(m+\nu+1)},
\]
and for the Airy series of \cref{prop:airy-series} the factors are $x^3/((3k-1)\,3k)$ and $x^3/(3k\,(3k+1))$. Each new term therefore costs a few arithmetic operations, and no factorials or powers need to be formed. \Cref{alg:series-recurrence} states the method. Its stopping rule follows the advice of \vcref{1}{alg:sum-series}, to stop on a bound for the remainder rather than on the size of the last term alone: if $|\rho_k|$ does not increase with $k$ and $|\rho_{m-1}|\le1/2$, then $|t_{m+j}|\le2^{-j}|t_m|$ for $j\ge1$, and the remainder after $t_m$ is at most $\sum_{j\ge1}2^{-j}|t_m|=|t_m|$. When the loop stops, the truncation error is therefore at most $u|s|$. For $J_\nu$ with $\nu\ge0$ the factors $|\rho_m|$ decrease, as required, and they fall below $1/2$ once $(m+1)(m+\nu+1)\ge x^2/2$.

\begin{algorithm}
\begin{algorithmic}[1]
\Require the first term $t_0$ and the ratios $\rho_m=t_{m+1}/t_m$ of a series $\sum_mt_m$, with $|\rho_m|$ nonincreasing in $m$ and tending to $0$; the unit roundoff $u$
\Ensure the partial sum $s=\sum_{k\le m}t_k$ for the first $m$ at which the remainder is at most $u|s|$
\Name{SumByRecurrence}{t_0,\rho,u}
\State $t\gets t_0$; $s\gets t_0$; $m\gets0$
\Repeat
  \State $t\gets\rho_m\,t$ \Comment{$t=t_{m+1}$}
  \State $s\gets s+t$
  \State $m\gets m+1$
\Until{$|\rho_{m-1}|\le1/2$ and $|t|\le u\,|s|$}
\State \Return $s$
\end{algorithmic}
\caption{Summing a series whose terms are generated by a recurrence. The test on $\rho_{m-1}$ ensures that the remaining terms decrease at least geometrically, so that their sum is at most $|t|$; rounding errors are not taken into account.}
\label{alg:series-recurrence}
\end{algorithm}

In exact arithmetic the algorithm computes $J_0(x)$ to any accuracy for every $x$, since the series converges everywhere. In floating-point arithmetic it does not, and the reason is the size of the terms.

\subsection{Cancellation}

We summed the series of $J_0(x)$ by \cref{alg:series-recurrence} in binary64 arithmetic and compared the results with values computed by the ball arithmetic of python-flint (version 0.9.0), which are correct to far more digits than shown. For comparison we also evaluated \code{special.j0}, the Bessel function of order $0$ of the stardust library, on the machine and with the version of the library described in \cref{sec:systems-practice}. \Cref{tab:j0-series} shows the results, and \cref{fig:j0-series-error} the error at $160$ points. Here and in the following tables the error of a computed number is measured against the binary64 number nearest to the exact value, so that an error $0$ means a correctly rounded result.

\begin{table}
\sdfloat{%
\begin{sdtable}{colspec={rrlllll}, row{1}={font=\small\bfseries}}
\toprule
$x$ & terms & largest term & $J_0(x)$ & ratio & error, series & error, \code{j0} \\
\midrule
$1$ & $10$ & $1.00$ & $0.765198$ & $1.3$ & $0$ & $1.5\times10^{-16}$ \\
$5$ & $19$ & $9.77$ & $-0.177597$ & $55$ & $9.4\times10^{-16}$ & $0$ \\
$10$ & $27$ & $6.78\times10^{2}$ & $-0.245936$ & $2.8\times10^{3}$ & $9.3\times10^{-14}$ & $1.1\times10^{-16}$ \\
$20$ & $42$ & $7.59\times10^{6}$ & $0.167025$ & $4.5\times10^{7}$ & $4.7\times10^{-10}$ & $1.7\times10^{-16}$ \\
$30$ & $57$ & $1.12\times10^{11}$ & $-0.0863680$ & $1.3\times10^{12}$ & $1.4\times10^{-4}$ & $0$ \\
$40$ & $71$ & $1.86\times10^{15}$ & $0.00736689$ & $2.5\times10^{17}$ & $13$ & $2.4\times10^{-16}$ \\
\bottomrule
\end{sdtable}}{The series of $J_0(x)$ summed in binary64 by \cref{alg:series-recurrence}: the number of terms added, the largest term in absolute value, the value $J_0(x)$, the ratio of the largest term to $|J_0(x)|$, and the relative error of the sum and of \code{special.j0}. At $x=40$ the sum is $-0.0901$, of the wrong sign.\label{tab:j0-series}}
\end{table}

\begin{figure}
\sdfloat{%
\begin{tikzpicture}
\begin{axis}[sdplot, ymode=log, width=0.62\textwidth, height=0.42\textwidth, xmin=0, xmax=40.5, ymin=1e-18, ymax=10,
  xtick={0,10,20,30,40}, ytick={1e-16,1e-12,1e-8,1e-4,1}, xlabel={$x$}, legend pos=north west]
\addplot[black, only marks, mark=*, mark size=0.9pt] coordinates {(0.5,1.11e-16) (0.75,1.11e-16) (1.25,1.11e-16) (2,5.55e-17) (2.25,1.39e-17) (2.75,2.78e-17) (3,1.11e-16) (3.25,5.55e-17) (4,5.55e-17) (4.25,4.44e-16) (4.5,5.55e-17) (4.75,5.00e-16) (5,1.67e-16) (5.25,4.16e-17) (5.5,1.40e-16) (5.75,3.47e-16) (6,5.55e-17) (6.25,1.14e-15) (6.5,3.89e-16) (6.75,8.88e-16) (7,1.67e-15) (7.25,2.83e-15) (7.5,5.55e-17) (7.75,4.69e-15) (8,1.78e-15) (8.25,3.58e-15) (8.5,1.04e-14) (8.75,2.49e-15) (9,1.88e-14) (9.25,6.47e-15) (9.5,5.16e-15) (9.75,2.31e-14) (10,2.30e-14) (10.25,2.07e-14) (10.5,9.99e-16) (10.75,1.12e-13) (11,1.13e-13) (11.25,1.54e-14) (11.5,1.61e-13) (11.75,1.75e-13) (12,2.17e-13) (12.25,4.62e-13) (12.5,8.00e-13) (12.75,3.27e-13) (13,5.90e-13) (13.25,4.85e-13) (13.5,1.53e-12) (13.75,2.73e-12) (14,1.03e-12) (14.25,2.23e-12) (14.5,1.34e-13) (14.75,4.56e-12) (15,2.61e-12) (15.25,1.25e-11) (15.5,1.39e-11) (15.75,4.77e-12) (16,3.38e-12) (16.25,1.83e-12) (16.5,1.94e-11) (16.75,1.81e-11) (17,4.57e-11) (17.25,4.63e-11) (17.5,1.40e-12) (17.75,5.09e-11) (18,1.16e-10) (18.25,6.18e-11) (18.5,9.72e-11) (18.75,2.47e-10) (19,2.35e-10) (19.25,4.69e-11) (19.5,1.56e-11) (19.75,8.25e-11) (20,7.77e-11) (20.25,6.96e-10) (20.5,2.78e-10) (20.75,2.23e-11) (21,1.69e-09) (21.25,3.22e-09) (21.5,5.41e-10) (21.75,6.23e-09) (22,1.34e-09) (22.25,8.80e-09) (22.5,6.06e-10) (22.75,8.25e-09) (23,2.28e-08) (23.25,6.81e-09) (23.5,1.50e-09) (23.75,8.65e-09) (24,2.02e-08) (24.25,2.04e-08) (24.5,6.56e-08) (24.75,3.83e-08) (25,4.95e-09) (25.25,2.04e-08) (25.5,1.86e-07) (25.75,8.29e-08) (26,1.11e-07) (26.25,2.06e-07) (26.5,5.98e-07) (26.75,2.14e-07) (27,8.58e-07) (27.25,8.64e-08) (27.5,6.32e-08) (27.75,7.21e-07) (28,1.32e-06) (28.25,2.06e-06) (28.5,1.80e-06) (28.75,2.13e-06) (29,5.99e-07) (29.25,5.92e-06) (29.5,5.41e-06) (29.75,7.33e-06) (30,1.23e-05) (30.25,3.83e-06) (30.5,9.40e-07) (30.75,9.66e-06) (31,2.32e-05) (31.25,3.26e-05) (31.5,8.33e-06) (31.75,6.23e-05) (32,1.19e-04) (32.25,3.86e-05) (32.5,6.87e-05) (32.75,3.40e-05) (33,1.45e-04) (33.25,3.68e-04) (33.5,2.97e-04) (33.75,2.83e-04) (34,4.85e-04) (34.25,2.58e-04) (34.5,1.60e-03) (34.75,2.72e-04) (35,8.04e-04) (35.25,1.15e-04) (35.5,8.48e-04) (35.75,3.00e-03) (36,1.24e-03) (36.25,8.32e-04) (36.5,4.14e-03) (36.75,4.44e-03) (37,4.13e-03) (37.25,5.60e-03) (37.5,8.23e-03) (37.75,1.36e-02) (38,7.34e-03) (38.25,1.67e-02) (38.5,2.41e-02) (38.75,4.47e-02) (39,1.22e-02) (39.25,4.00e-03) (39.5,7.76e-02) (39.75,1.77e-01) (40,9.75e-02)};
\addlegendentry{error of the sum}
\addplot[black, dashed, line width=0.7pt] coordinates {(0.25,1.11e-16) (0.5,1.11e-16) (0.75,1.11e-16) (1,1.11e-16) (1.25,1.11e-16) (1.5,1.11e-16) (1.75,1.11e-16) (2,1.11e-16) (2.25,1.41e-16) (2.5,1.73e-16) (2.75,2.10e-16) (3,2.50e-16) (3.25,2.93e-16) (3.5,3.40e-16) (3.75,3.90e-16) (4,4.44e-16) (4.25,5.66e-16) (4.5,7.11e-16) (4.75,8.83e-16) (5,1.08e-15) (5.25,1.32e-15) (5.5,1.59e-15) (5.75,1.90e-15) (6,2.25e-15) (6.25,2.87e-15) (6.5,3.63e-15) (6.75,4.56e-15) (7,5.67e-15) (7.25,7.00e-15) (7.5,8.58e-15) (7.75,1.04e-14) (8,1.26e-14) (8.25,1.62e-14) (8.5,2.05e-14) (8.75,2.59e-14) (9,3.24e-14) (9.25,4.04e-14) (9.5,5.00e-14) (9.75,6.15e-14) (10,7.53e-14) (10.25,9.64e-14) (10.5,1.23e-13) (10.75,1.55e-13) (11,1.95e-13) (11.25,2.44e-13) (11.5,3.05e-13) (11.75,3.78e-13) (12,4.66e-13) (12.25,5.97e-13) (12.5,7.61e-13) (12.75,9.65e-13) (13,1.22e-12) (13.25,1.53e-12) (13.5,1.92e-12) (13.75,2.39e-12) (14,2.96e-12) (14.25,3.80e-12) (14.5,4.84e-12) (14.75,6.15e-12) (15,7.79e-12) (15.25,9.82e-12) (15.5,1.23e-11) (15.75,1.54e-11) (16,1.92e-11) (16.25,2.46e-11) (16.5,3.15e-11) (16.75,4.00e-11) (17,5.07e-11) (17.25,6.40e-11) (17.5,8.06e-11) (17.75,1.01e-10) (18,1.27e-10) (18.25,1.62e-10) (18.5,2.07e-10) (18.75,2.64e-10) (19,3.35e-10) (19.25,4.24e-10) (19.5,5.35e-10) (19.75,6.72e-10) (20,8.43e-10) (20.25,1.08e-09) (20.5,1.38e-09) (20.75,1.76e-09) (21,2.24e-09) (21.25,2.83e-09) (21.5,3.58e-09) (21.75,4.51e-09) (22,5.67e-09) (22.25,7.27e-09) (22.5,9.30e-09) (22.75,1.19e-08) (23,1.51e-08) (23.25,1.91e-08) (23.5,2.42e-08) (23.75,3.06e-08) (24,3.85e-08) (24.25,4.93e-08) (24.5,6.31e-08) (24.75,8.05e-08) (25,1.02e-07) (25.25,1.30e-07) (25.5,1.65e-07) (25.75,2.08e-07) (26,2.63e-07) (26.25,3.37e-07) (26.5,4.31e-07) (26.75,5.50e-07) (27,7.01e-07) (27.25,8.90e-07) (27.5,1.13e-06) (27.75,1.43e-06) (28,1.80e-06) (28.25,2.31e-06) (28.5,2.96e-06) (28.75,3.78e-06) (29,4.82e-06) (29.25,6.13e-06) (29.5,7.78e-06) (29.75,9.85e-06) (30,1.24e-05) (30.25,1.60e-05) (30.5,2.04e-05) (30.75,2.61e-05) (31,3.33e-05) (31.25,4.24e-05) (31.5,5.38e-05) (31.75,6.82e-05) (32,8.63e-05) (32.25,1.11e-04) (32.5,1.42e-04) (32.75,1.81e-04) (33,2.31e-04) (33.25,2.94e-04) (33.5,3.74e-04) (33.75,4.74e-04) (34,6.01e-04) (34.25,7.70e-04) (34.5,9.87e-04) (34.75,1.26e-03) (35,1.61e-03) (35.25,2.05e-03) (35.5,2.61e-03) (35.75,3.31e-03) (36,4.19e-03) (36.25,5.38e-03) (36.5,6.89e-03) (36.75,8.81e-03) (37,1.12e-02) (37.25,1.43e-02) (37.5,1.82e-02) (37.75,2.32e-02) (38,2.94e-02) (38.25,3.77e-02) (38.5,4.83e-02) (38.75,6.17e-02) (39,7.88e-02) (39.25,1.00e-01) (39.5,1.28e-01) (39.75,1.63e-01) (40,2.06e-01)};
\addlegendentry{$u$ times the largest term}
\end{axis}
\end{tikzpicture}}{The absolute error of the series \eqref{eq:bessel-j} for $J_0$, summed in binary64 by \cref{alg:series-recurrence} at $x=0.25,0.5,\ldots,40$ (dots; points at which the sum is correctly rounded are not shown), and $u$ times the largest term of the series (dashed).\label{fig:j0-series-error}}
\end{figure}

For small $x$ the series is excellent: at $x=1$ ten terms give the correctly rounded value. As $x$ grows, the terms alternate in sign and grow to a largest term before they decrease, and the sum is far smaller than the largest term: at $x=20$ the largest term is $7.6\times10^6$ and the sum $0.167$, at $x=40$ the largest term is $1.9\times10^{15}$ and the sum $0.0074$. Each term carries a rounding error of up to $u$ times its size, and these errors are not reduced by the cancellation that makes the sum small. The measured errors follow the ratio of the largest term to the value: at all $160$ points of \cref{fig:j0-series-error} the absolute error was less than $2u$ times the largest term, and at $x=40$ no digit is correct. The sum of the absolute values of the terms is $\sum_m(x/2)^{2m}/(m!)^2$, the series of the modified Bessel function $I_0(x)$, which at $x=40$ is $1.5\times10^{16}$; for large $x$ it grows like $e^x/\sqrt{2\pi x}$~\cite[Sect.~10.40]{olver2010nist}, while $|J_0(x)|$ is at most of the order $x^{-1/2}$ by \cref{thm:bessel-asymptotic}. This is the phenomenon that \vcref{1}{sec:expansions} met for $e^{-20}$ and \cref{sec:computing-matrix-functions} for the exponential of a matrix (\cref{tab:exp-mvl}): a convergent series whose terms are much larger than its sum cannot be summed accurately in floating-point arithmetic, whatever the order of summation.

The library function, in contrast, has a relative error of at most $2.4\times10^{-16}$ at the six points of the table. It does not sum the power series for large arguments; how special functions are computed accurately over their whole domain, by asymptotic expansions, recurrences, continued fractions and polynomial approximations chosen region by region, is the subject of Volume~XI.

The Airy function shows the same effect more drastically. For $x>0$, $\operatorname{Ai}(x)=c_1f(x)-c_2g(x)$ is a small difference of the two series $f$ and $g$, which grow like $\operatorname{Bi}(x)$, while $\operatorname{Ai}$ decays (\cref{sec:airy}). \Cref{tab:airy-series} shows the result of summing $c_1f-c_2g$ in the manner of \cref{alg:series-recurrence}, the two series together: the $k$th terms of $f$ and $g$ are generated by their recurrences, the difference of the $k$th terms of $c_1f$ and $c_2g$ is added to the sum, and the loop stops when both ratios are at most $1/2$ and that difference is at most $u$ times the sum. The results are compared with python-flint and with the function \code{special.airy} of the library. At $x=10$, where $\operatorname{Ai}(x)=1.1\times10^{-10}$ and the largest term is $2.3\times10^7$, the computed sum is $-5.2\times10^{-8}$, of the wrong sign and almost $500$ times too large in modulus. For negative $x$ the loss is also present, but smaller at the same $|x|$, because $|\operatorname{Ai}|$ does not decrease exponentially there.

\begin{table}
\sdfloat{%
\begin{sdtable}{colspec={rrlllll}, row{1}={font=\small\bfseries}}
\toprule
$x$ & terms & largest term & $\operatorname{Ai}(x)$ & ratio & error, series & error, \code{airy} \\
\midrule
$1$ & $9$ & $0.355$ & $1.35292\times10^{-1}$ & $2.6$ & $6.2\times10^{-16}$ & $4.1\times10^{-16}$ \\
$2$ & $13$ & $0.518$ & $3.49241\times10^{-2}$ & $15$ & $4.4\times10^{-15}$ & $1.2\times10^{-15}$ \\
$4$ & $21$ & $8.41$ & $9.51564\times10^{-4}$ & $8.8\times10^{3}$ & $2.7\times10^{-12}$ & $3.4\times10^{-16}$ \\
$6$ & $30$ & $4.78\times10^{2}$ & $9.94769\times10^{-6}$ & $4.8\times10^{7}$ & $2.3\times10^{-8}$ & $5.1\times10^{-16}$ \\
$8$ & $40$ & $7.12\times10^{4}$ & $4.69221\times10^{-8}$ & $1.5\times10^{12}$ & $1.6\times10^{-3}$ & $1.4\times10^{-16}$ \\
$10$ & $49$ & $2.29\times10^{7}$ & $1.10475\times10^{-10}$ & $2.1\times10^{17}$ & $4.7\times10^{2}$ & $2.3\times10^{-16}$ \\
$-10$ & $45$ & $2.29\times10^{7}$ & $4.02412\times10^{-2}$ & $5.7\times10^{8}$ & $1.1\times10^{-7}$ & $1.7\times10^{-16}$ \\
$-20$ & $89$ & $6.27\times10^{23}$ & $-1.76406\times10^{-1}$ & $3.6\times10^{24}$ & $4.2\times10^{8}$ & $0$ \\
\bottomrule
\end{sdtable}}{The Airy function $\operatorname{Ai}(x)=c_1f(x)-c_2g(x)$ computed in binary64 from the series of \cref{prop:airy-series} by \cref{alg:series-recurrence}: the number of terms, the largest term $c_1|a_{3k}x^{3k}|$ or $c_2|a_{3k+1}x^{3k+1}|$, the value, the ratio of the largest term to $|\operatorname{Ai}(x)|$, and the relative errors of the sum and of \code{special.airy}. The number of terms counts the pairs of terms of $f$ and $g$.\label{tab:airy-series}}
\end{table}

\subsection{Recurrences Can Do Better}

The cancellation is not a property of the functions but of the representation. The Legendre polynomials give a clean example, because two exact representations are available: the coefficients of Rodrigues' formula (\cref{thm:rodrigues}) and the three-term recurrence (\cref{thm:legendre-recurrence}). We evaluated $P_n(x)$ at the binary64 number $x$ nearest to $0.9$ in both ways, by Horner's rule with the coefficients rounded to binary64, and by the recurrence started from $P_0=1$ and $P_1=x$, and compared with the exact rational values of $P_n(x)$. \Cref{tab:legendre-evaluation} shows the results. The coefficients of $P_n$ alternate in sign and grow rapidly with $n$; the largest coefficient of $P_{50}$ is $1.9\times10^{17}$, larger than $2^{53}$, so that most of them are not even representable exactly: $25$ of the $26$ nonzero coefficients of $P_{50}$ are not binary64 numbers. The terms $c_kx^k$ at $x=0.9$ reach $4.9\times10^{15}$ for $n=50$, while $P_{50}(0.9)=-0.170$, and the evaluation from the coefficients has no correct digit. The recurrence never forms large numbers: its errors stayed below $8\times10^{-16}$ for all five degrees.

\begin{table}
\sdfloat{%
\begin{sdtable}{colspec={rlllll}, row{1}={font=\small\bfseries}}
\toprule
$n$ & $P_n(0.9)$ & largest $|c_kx^k|$ & ratio & error, coefficients & error, recurrence \\
\midrule
$10$ & $-0.2631456$ & $1.87\times10^{2}$ & $7.1\times10^{2}$ & $4.2\times10^{-14}$ & $2.1\times10^{-16}$ \\
$20$ & $-0.1493082$ & $3.46\times10^{5}$ & $2.3\times10^{6}$ & $1.7\times10^{-10}$ & $7.4\times10^{-16}$ \\
$30$ & $0.2018124$ & $7.60\times10^{8}$ & $3.8\times10^{9}$ & $8.9\times10^{-8}$ & $0$ \\
$40$ & $0.03698742$ & $1.82\times10^{12}$ & $4.9\times10^{13}$ & $6.2\times10^{-3}$ & $5.6\times10^{-16}$ \\
$50$ & $-0.1700377$ & $4.90\times10^{15}$ & $2.9\times10^{16}$ & $2.2$ & $4.9\times10^{-16}$ \\
\bottomrule
\end{sdtable}}{The Legendre polynomial $P_n(x)=\sum_kc_kx^k$ at $x=0.9$ in binary64: the largest term $|c_kx^k|$, its ratio to $|P_n(0.9)|$, and the relative errors of Horner's rule with the coefficients of \cref{thm:rodrigues} rounded to binary64 and of the three-term recurrence of \cref{thm:legendre-recurrence}. Exact values at the binary64 number nearest to $0.9$ were computed in rational arithmetic.\label{tab:legendre-evaluation}}
\end{table}

The experiments teach the lesson of this section: a power series is a definition of a function, and a proof that it exists, but not always an algorithm for computing it. Where the terms are not much larger than the sum, as for $J_0(x)$ at $x=1$ and $x=5$ or for $\operatorname{Ai}(1)$ in the tables, the series is an excellent method, and it is used in this way. Elsewhere other representations take over: the asymptotic forms of \cref{thm:bessel-asymptotic} for large arguments, recurrences in the order or the degree, run in the direction in which they are stable, continued fractions, and approximations fitted to each region. \Cref{ch:difference-equations} studies the stability of recurrences, which decides the direction, and Volume~XI builds the special functions of the stardust library from these pieces.

\bigskip
Near an ordinary point, every solution of a linear equation whose coefficients are power series is itself a power series, which converges at least as far as the series of the coefficients; its coefficients follow from a recurrence, obtained by substituting the series and comparing coefficients, and their growth is controlled by a simple majorant argument. Systems and equations of higher order behave in the same way. For Airy's equation the series converge everywhere and give the functions $\operatorname{Ai}$ and $\operatorname{Bi}$, and for $w''+t^2w=0$ they give the solution of the Riccati equation of \cref{ch:first-order-equations} and its blow-up time. For Legendre's equation the series terminate exactly when the degree is a natural number, and the resulting polynomials are the Legendre polynomials of \cref{ch:inner-products}, with Rodrigues' formula, the orthogonality relations and the three-term recurrence; all other solutions are unbounded at the ends of $(-1,1)$. Near a regular singular point the solutions are powers times power series, possibly with a logarithm, according to the indicial roots, as Frobenius's theorem states; at an irregular singular point the formal series may diverge. Bessel's equation has a regular singular point at $0$, and the method gives the Bessel functions $J_\nu$, the functions $Y_\nu$ for integral order with their logarithms, the elementary functions of half-integral order, the recurrence relations and the zeros, spaced asymptotically $\pi$ apart. The classical orthogonal polynomials and the hypergeometric function arise in the same way. Finally, the series are reliable algorithms only where their terms are not much larger than their sums.

The next chapter turns from differential equations to their discrete counterparts, the difference equations and recurrences. Linear recurrences with constant coefficients are solved by their characteristic roots, as linear differential equations were in \cref{ch:linear-higher-order}, and the three-term recurrences of this chapter, for the Legendre polynomials and for the Bessel functions of integral order, are examples of the question that matters most in computation: in which direction a recurrence may be run without amplifying its rounding errors.

\chapternotes

The method of majorants, which proves the convergence of power series solutions by comparison, goes back to Cauchy; Hille~\cite[Chap.~2]{hille1976} presents it and its use for equations in the complex domain. Fuchs characterised in 1866 the singular points at which all solutions have at most power-like singularities, the regular singular points of \cref{sec:frobenius}~\cite{fuchs1866}, and Frobenius gave the method of \cref{thm:frobenius} in 1873~\cite{frobenius1873}; Gray~\cite{gray2000} tells the history of the theory of linear differential equations in the complex domain from Riemann to Poincar\'e. Bessel introduced the functions that bear his name in 1824, in a study of planetary perturbations, through the integral of the note on generating functions~\cite{bessel1824}; functions of the same kind had appeared earlier in the work of Daniel Bernoulli and Euler, and Watson~\cite[Chap.~I]{watson1944} describes these early appearances. Legendre introduced his polynomials in a memoir on the attraction of spheroids~\cite{legendre1785}, and Rodrigues found his formula in 1816~\cite{rodrigues1816}. Airy's function appeared in his study of the intensity of light near a caustic~\cite{airy1838}.

The series solution of linear equations is treated at the level of this chapter by Boyce, DiPrima and Meade~\cite[Chap.~5]{boyce2017} and, with complete proofs, by Coddington~\cite[Chaps.~3 and~4]{coddington1961}. Ince~\cite{ince1926} remains a rich source of methods and examples, and Coddington and Levinson~\cite[Chaps.~4 and~5]{coddington1955}, Hartman~\cite{hartman2002} and Teschl~\cite[Chap.~4]{teschl2012} treat singular points in the complex domain, where the theory is most natural. For the special functions themselves, Watson's treatise~\cite{watson1944} is the classic on Bessel functions, Whittaker and Watson~\cite{whittaker1927} cover the hypergeometric and Legendre functions, and Lebedev~\cite{lebedev1972}, Temme~\cite{temme1996} and Andrews, Askey and Roy~\cite{andrews1999} are modern introductions; Szeg\H{o}~\cite{szego1975} is the standard work on orthogonal polynomials. The NIST Handbook~\cite{olver2010nist} collects the formulas, and Olver~\cite{olver1997asymptotics} develops the asymptotic methods that take over where the series fail. Vall\'ee and Soares~\cite{vallee2004} describe the Airy functions and their applications in physics, and Gautschi~\cite{gautschi1967} wrote the classical survey of the computation of special functions by recurrences.

\chapterreferences
